\PassOptionsToPackage{no-math}{fontspec}
% texlate: native XeTeX font and PDF-driver capabilities
\AddToHook{package/microtype/after}{%
\DeclareMicrotypeSet{texlate-native}{encoding={TU,EU1,EU2}}%
\UseMicrotypeSet[protrusion]{texlate-native}%
}
% breakurl already has a native-PDF branch; limit its PDF-mode override to loading.
\AddToHook{package/breakurl/before}{%
\RequirePackage{xkeyval,ifpdf}%
\let\TeXlateSavedIfpdf\ifpdf\let\ifpdf\iftrue
}
\AddToHook{package/breakurl/after}{\let\ifpdf\TeXlateSavedIfpdf}
% This class's optional arXiv check assumes every non-pdfTeX engine writes DVI.
\PassOptionsToClass{nopdfoutputerror,allowfontchageintitle}{quantumarticle}
% Embedded PostScript can silently disappear with restricted XeTeX drivers.
% Record actual drawing operations; merely loading an unused package is harmless.
\AddToHook{package/pstricks/after}{%
\ifcsname pst@object\endcsname
\expandafter\let\expandafter\TeXlatePstObject\csname pst@object\endcsname
\expandafter\def\csname pst@object\endcsname#1{%
\typeout{TeXlate-PostScript-object: #1}\TeXlatePstObject{#1}}%
\fi
}
 
% !TEX TS-program = latex
% kate: default-dictionary en_GB;
\documentclass[
    affil-it, % affiliations in italic
    %affil-sl, % affiliations in sl
    auth-lg, % authors in large
    %auth-sc, % authors in sc
    british, % dates in British format
    compression, % compressed title page
    contents, % show table of contents
    custom-packages={}, % Load listed custom packages
    %draft-mode, % only title and date are needed
    references={bigone}, % Name of bibliography
]{lib/preprint}
\usepackage[fontset=fandol,UTF8]{ctex}  % [texlate injected]
% texlate: twin named-dest for shared-counter theorems
\makeatletter
\def\TeXlate@thmtwin{%
  \@ifundefined{MakeLinkTarget}{}{%
  \@ifundefined{@currentcounter}{}{%
  \@ifundefined{theH\@currentcounter}{}{%
    \ifx\@currenvir\@currentcounter
      \@ifundefined{alias@ctr@\@currenvir}{}{%
        \edef\TeXlate@twin{\csname alias@ctr@\@currenvir\endcsname}%
        \TeXlate@emit}%
    \else
      \@ifundefined{the\@currenvir}{}{%
        \let\TeXlate@twin\@currenvir
        \TeXlate@emit}%
    \fi}}}%
}%
\def\TeXlate@emit{%
  \begingroup
  \let\TeXlate@save\@currentHref
  \edef\TeXlate@name{\TeXlate@twin.\csname theH\@currentcounter\endcsname}%
  \MakeLinkTarget*{\TeXlate@name}%
  \global\let\@currentHref\TeXlate@save
  \endgroup}%
\AddToHook{cmd/@begintheorem/before}{\TeXlate@thmtwin}%
\AddToHook{cmd/@opargbegintheorem/before}{\TeXlate@thmtwin}%
\makeatother


%
% Local settings and commands
%

\newcommand{\boundary}{\flat}
\newcommand{\grp}[1]{\underline{#1}}
\renewcommand{\Re}{\mathrm{Re}}
\renewcommand{\Im}{\mathrm{Im}}
\newcommand{\coker}{\mathrm{coker}}
\renewcommand{\pr}{\mathrm{pr}}

%
% Comments
%

% Use \eqref for equations and \cref for anything else

%
% Document
%

\begin{document}

    % Hypersetup
    \hypersetup{
        pdftitle = {Explicit Non-Abelian Gerbes with Connections},
        pdfauthor = {Dominik Rist,Christian Saemann,Martin Wolf},
        pdfkeywords = {},
    }
    
    % Date of preprint
    \date{\today}
    
    % All emails in the order of appearance
    \email{dr40@hw.ac.uk,c.saemann@hw.ac.uk,m.wolf@surrey.ac.uk}
    
    % All preprint numbers in the order of appearance
    \preprint{EMPG--22--02,DMUS--MP--22/01}
    
    % Title
    \title{Explicit Non-Abelian Gerbes with Connections} 
    
    % All authors
    \author[a]{Dominik~Rist\,\orcidlink{0000-0002-1817-3458}\,}
    \author[a]{Christian~Saemann\,\orcidlink{0000-0002-5273-3359}\,}
    \author[b]{Martin~Wolf\,\orcidlink{0009-0002-8192-3124}\,}
    
    % All affiliations
    \affil[a]{Maxwell Institute for Mathematical Sciences, Department of Mathematics,\\ Heriot--Watt University, Edinburgh EH14 4AS, United Kingdom}
    \affil[b]{School of Mathematics and Physics,\\ University of Surrey, Guildford GU2 7XH, United Kingdom}
    
    % Abstract
    \abstract{We define the notion of adjustment for strict Lie 2-groups and provide the complete cocycle description for non-Abelian gerbes with connections whose structure 2-group is an adjusted 2-group. Most importantly, we depart from the common fake-flat connections and employ adjusted connections. This is an important generalisation that is needed for physical applications especially in the context of supergravity. We give a number of explicit examples; in particular, we lift the spin structure on $S^4$, corresponding to an instanton--anti-instanton pair, to a string structure, a 2-group bundle with connection. We also outline how categorified forms of Bogomolny monopoles known as self-dual strings can be obtained via a Penrose--Ward transform of string bundles over twistor space.}
    
    % Acknowledgements
    \acknowledgements{We thank Hyungrok Kim and Paul Skerritt for fruitful discussions. We are particularly grateful to David Roberts for his contribution to the `Workshop on Higher Gauge Theory and Higher Quantization' in 2014 and for a number of helpful discussions related to signs in the literature. The work of CS was partially supported by the Leverhulme Research Project Grant RPG-2018-329 `The Mathematics of M5-Branes'.}
    
    % Data and Licence Statement
    \datalicencemanagement{No additional research data beyond the data presented and cited in this work are needed to validate the research findings in this work. For the purpose of open access, the authors have applied a Creative Commons Attribution (CC BY) licence to any Author Accepted Manuscript version arising.}
    
    % Body
    \begin{body}
        
        \section{这是译文}   
        
 这是译文 $B$-这是译文. 
        
 这是译文~\cite{Murray:9407015,Murray:2007ps} 这是译文~\cite{Hitchin:1999fh,Chatterjee:1998}. 这是译文 $B$-这是译文~\cite{Bergshoeff:1981um,Chapline:1982ww} 这是译文~\cite{Nikolaus:2018qop}. 这是译文~\cite{Kim:2022opr}. 这是译文~\cite{Saemann:2017zpd,Saemann:2019dsl,Rist:2020uaa}; 这是译文~\cite{Fiorenza:2020hiq} 这是译文 $\sfString(3)$ 这是译文.\footnote{这是译文~\cite{Jurco:2019woz} 这是译文.} 这是译文~\cite{Saemann:2012uq,Saemann:2013pca,Jurco:2014mva,Jurco:2016qwv,Samann:2017dah}.
        
 这是译文.~\cite{Nikolaus:2011ag} 这是译文~\cite{Breen:math0106083} 这是译文~\cite{Aschieri:2003mw}, 这是译文. 
        
这是译文~\cite{Breen:math0106083,Aschieri:2003mw}, 这是译文\footnote{这是译文~\cite{Breen:math0106083,Aschieri:2003mw}, 这是译文 $\delta_{ij}$.} 这是译文 $L_\infty$-这是译文~\cite{Sati:2008eg}. 这是译文~\cite{Baez:0511710,Schreiber:0705.0452,Schreiber:2008aa} 这是译文~\cite{Baez:0511710}, 这是译文 $F$ 这是译文~\cite{Schreiber:0705.0452,Schreiber:2008aa}. 这是译文~\cite{Gastel:2018joi,Saemann:2019dsl}, 这是译文.
        
 这是译文~\cite{Baez:0511710}, 这是译文~\cite{Breen:math0106083,Aschieri:2003mw}. 这是译文 \uline{这是译文} 这是译文. 
        
 这是译文~\cite{Bergshoeff:1981um,Chapline:1982ww}. 这是译文~\cite{Sati:2008eg,Sati:2009ic,Fiorenza:2010mh}. 这是译文 (这是译文) 这是译文~这是译文.
        
 这是译文.~这是译文~\cite{Baez:0511710}. 这是译文~\cite{Saemann:2019dsl}, 这是译文~\cite{Kim:2019owc}.\footnote{这是译文~\cite{Kim:2019owc} 这是译文 \v 这是译文 \v 这是译文.} 这是译文~\cite{Borsten:2021ljb}, 这是译文.
        
 这是译文~\cite{Kim:2022opr}. 这是译文.
        
 这是译文 $n$-这是译文 $X$, 这是译文 $\sfSO(n)$-这是译文 $\sfSpin(n)$-这是译文 $X$. 这是译文 $LX$ 这是译文 $X$ 这是译文\footnote{这是译文 $X$ 这是译文 $n\geq5$. 这是译文 $X$ 这是译文.} 这是译文 $X$ 这是译文 $\sfString(n)$-这是译文~\cite{mclaughlin1992orientation}. 这是译文, $\sfString(n)$ 这是译文 $\sfSpin(n)$, 这是译文 $\sfSpin(n)$. 
        
 这是译文 $\sfString(n)$ 这是译文.~\cite{Nikolaus:2011zg}, 这是译文~\cite{Baez:2005sn}. 这是译文~\cite{Schommer-Pries:0911.2483}.\footnote{这是译文~\cite{Demessie:2016ieh} 这是译文 (这是译文) 这是译文.} 这是译文.
        
这是译文:
        \begin{enumerate}[(i)]\itemsep-2pt
            \item 这是译文 \cref{ssec:adjustedCrossedModules}, 这是译文.
            \item 这是译文\footnote{这是译文'.} 这是译文 (这是译文) 这是译文 \cref{sec:adjustedDescentData}.
            \item 这是译文 $\sfG$ 这是译文 $\caL\sfG$ 这是译文 $\sfG$, 这是译文 \cref{ssec:adjustedCrossedModules}. 这是译文 \cref{sec:principalAdjustedHigherPrincipal}.
            \item 这是译文 $S^4$, 这是译文 $\sfSpin(4)$-这是译文 $\sfSpin(5)\rightarrow S^4\cong\sfSpin(5)/\sfSpin(4)$, 这是译文 $\caL\sfSpin(4)$-这是译文 \cref{sec:Spin4Bundle}.
            \item 这是译文 \cref{sec:instantonAntiInstantonStringLift}. 这是译文 $\sfString(3)/\sfString(2)\cong S^3$, 这是译文 $\sfString(3)$ 这是译文 $S^3$.
            \item 这是译文 \cref{sec:instantonAntiInstantonStringLift}. 这是译文 $\caL\sfSpin(4)$-这是译文 $S^4$ 这是译文 $\sfString(4)$-这是译文.\footnote{这是译文 $\sfString(4)$ 这是译文 $\sfSpin(4)\cong\sfSU(2)\times\sfSU(2)$, 这是译文~\cite{Bagger:2007jr,Gustavsson:2007vu}.} 
            \item 这是译文~\cite{Saemann:2017rjm}, 这是译文 \cref{sec:nonAbelianSelfDualString}.
            \item 这是译文 \cref{sec:commentsTwistorApproach}, 这是译文~\cite{Saemann:2017rjm}.
        \end{enumerate}
这是译文 \cref{app:proofs}.
        
 这是译文.

        \begin{remark}
            这是译文 $X,Y,\ldots$ 这是译文, $\Omega^k(X)$ 这是译文 $k$-这是译文 $X$. 这是译文 $\Omega^0(X)=\scC^\infty(X)$. 
            
 这是译文 $\sfG,\sfH,\ldots$ 这是译文 (这是译文) 这是译文 $\unit$. 这是译文 $g$ 这是译文 $X$ 这是译文 $g^{-1}Xg$ 这是译文.
        \end{remark}

        \section{这是译文}\label{sec:higherPrincipalBundles}
        
 这是译文.~\cite{Schreiber:2008aa}. 这是译文~\cite{Saemann:2019dsl}, 这是译文~\cite{Kim:2019owc,Borsten:2021ljb}, 这是译文. 

        \subsection{这是译文}
        
        \paragraph{这是译文.}
 这是译文 \uline{这是译文} $\caG\coloneqq(\sfH\overset{\sft}{\longrightarrow}\sfG,\acton)$ 这是译文 $\sfG$ 这是译文 $\sfH$, 这是译文 $\acton$ 这是译文 $\sfG$ 这是译文 $\sfH$, 这是译文 $\sft:\sfH\rightarrow\sfG$ 这是译文         \begin{equation}\label{eq:crossedModuleConditions}
            \sft(g\acton h_1)\ =\ g\sft(h_1)g^{-1}
            \eand
            \sft(h_1)\acton h_2\ =\ h_1h_2h_1^{-1}
        \end{equation}  这是译文 $g\in\sfG$ 这是译文 $h_{1,2}\in\sfH$; 这是译文 \uline{这是译文}. 这是译文.~\cite{Baez:0307200}. 
        
 这是译文, $\caG_{1,2}=(\sfH_{1,2}\overset{\sft_{1,2}}{\longrightarrow}\sfG_{1,2},\acton_{1,2})$, 这是译文 \uline{这是译文} $\phi:\caG_1\rightarrow \caG_2$ 这是译文 $\phi_\sfG:\sfG_1\rightarrow \sfG_2$ 这是译文 $\phi_\sfH:\sfH_1\rightarrow \sfH_2$ 这是译文 
        \begin{equation}
			\begin{aligned}
				(\sft_2\circ \phi_\sfH)(h)=(\phi_\sfG\circ \sft_1)(h)
                \eand 
                \phi_\sfH(g\acton_1 h)=\phi_\sfG(g)\acton_2\phi_\sfH(h)
			\end{aligned}
		\end{equation}  这是译文 $g\in \sfG_1$ 这是译文 $h\in \sfH_1$. 这是译文 \uline{这是译文} 这是译文 $\ker(\sft_1)\subseteq\sfH_1$ 这是译文 $\coker(t_1)=\sfG_1/\im(\sft)$.
        
 这是译文, \uline{这是译文} 这是译文~\cite{Aldrovandi:0808.3627}. 这是译文, $\caG_{1,2}=(\sfH_{1,2}\overset{\sft_{1,2}}{\longrightarrow}\sfG_{1,2},\acton_{1,2})$, 这是译文 \uline{这是译文} 这是译文         \begin{subequations}\label{eq:butterfly}
            \begin{equation}
                \begin{tikzcd}
                    \sfH_1 \arrow[dd,"\sft_1",swap] \arrow[dr,"\lambda_1"]& & \sfH_2 \arrow[dl,"\lambda_2",swap] \arrow[dd,"\sft_2"]
                    \\
                    & \sfE \arrow[dl,"\gamma_1",swap] \arrow[dr,"\gamma_2"]& 
                    \\
                    \sfG_1 & & \sfG_2
                \end{tikzcd}
            \end{equation}
            where $\sfE$ is a Lie group, $\lambda_{1,2}$ and $\gamma_{1,2}$ are morphisms of Lie groups, the NE--SW diagonal is a short exact sequence (i.e.~a Lie group extension), and the NW--SE diagonal is a complex. In addtion, the actions need to compatible in the sense of
            \begin{equation}
                \lambda_{1,2}(\gamma_{1,2}(e)\acton_{1,2}h_{1,2})\ =\ e\lambda_{1,2}(h_{1,2})e^{-1} 
            \end{equation}
        \end{subequations}  这是译文 $e\in\sfE$ 这是译文 $h_{1,2}\in\sfH_{1,2}$. 这是译文 \uline{这是译文} 这是译文 $\caG_1$ 这是译文 $\caG_2$ \uline{这是译文}. 
        
 这是译文 $\caG_{1,2}=(\sfH_{1,2}\overset{\sft}{\longrightarrow}\sfG_{1,2},\acton_{1,2})$ 这是译文~\cite[Remark 8.5]{Noohi:0506313} 这是译文         \begin{equation}
            \begin{tikzcd}
                & \caK \arrow[dl,"\phi"'] \arrow[dr,"\psi"]
                \\
                \caG_1 & & \caG_2
            \end{tikzcd}
		\end{equation}  这是译文 $\caK$ 这是译文 $\sfH_1\times \sfH_2\rightarrow \sfE$, 这是译文 $\phi$ 这是译文 $\psi$ 这是译文 $\phi$ 这是译文.
        
        \paragraph{这是译文.} 这是译文 $\caG=(\sfH\overset{\sft}{\longrightarrow}\sfG,\acton)$, 这是译文 \uline{这是译文} $\sfLie(\caG)\coloneqq(\frh\overset{\sft}{\longrightarrow}\frg,\acton)$ 这是译文 $\frg$ 这是译文 $\frh$ 这是译文 $\sfG$ 这是译文 $\sfH$, 这是译文.\footnote{这是译文 $\acton$ 这是译文 $\sft$ 这是译文.} 这是译文~\eqref{eq:crossedModuleConditions} 这是译文 
        \begin{equation}\label{eq:crossedModuleLieAlgebras}
            \sft(V\acton W_1)\ =\ [V,\sft(W_1)]
            \eand
            \sft(W_1)\acton W_2\ =\ [W_1,W_2]
        \end{equation}  这是译文 $V\in\frg$ 这是译文 $W_{1,2}\in\frh$. 
        
 这是译文 (这是译文.~这是译文) 这是译文 $\sfLie(\caG_{1,2})=(\frh_{1,2}\xrightarrow{~t_{1,2}~}\frg_{1,2})$,         \begin{equation}
            \begin{tikzcd}
                \frh_1 \arrow[dd,"\sft_1",swap] \arrow[dr,"\lambda_1"]& & \frh_2 \arrow[dl,"\lambda_2",swap] \arrow[dd,"\sft_2"]
                \\
                & \fre \arrow[dl,"\gamma_1",swap] \arrow[dr,"\gamma_2"]& 
                \\
                \frg_1 & & \frg_2
            \end{tikzcd}
        \end{equation}  这是译文 $\sigma:\frg_1\rightarrow \fre$ 这是译文 $\gamma_1$ 这是译文 $L_\infty$-这是译文 $\sfLie(\caG_{1})$ 这是译文 $\sfLie(\caG_{2})$, 这是译文 $L_\infty$-这是译文~\cite[Prop.~3.4]{Noohi:0910.1818}. 这是译文 $\frg_1$ 这是译文 $\tau:\fre\rightarrow \frh_2$, 这是译文 
        \begin{equation}
            \begin{aligned}
                \phi_\frh&:\frh_1\rightarrow \frh_2~,~~~&\phi_\frh(W)&=-\tau(\lambda_1(W))~,
                \\
                \phi_\frg&:\frg_1\rightarrow \frg_2~,~~~&\phi_\frg(V)&=\gamma_2(\sigma(V))~,
                \\
                \phi_2&:\frg_1\wedge \frg_1\rightarrow \frh_2~,~~~&\phi_2(V_1,V_2)&=-\tau([\sigma(V_1),\sigma(V_2)]_\fre)
            \end{aligned}
        \end{equation}  这是译文 $V,V_{1,2}\in \frg_1$ 这是译文 $W\in \frh_1$ 这是译文 $L_\infty$-这是译文 $\sfLie(\caG_{1})$ 这是译文 $\sfLie(\caG_{2})$. 这是译文, 
        \begin{equation}\label{eq:L_infty_alg}
			\begin{aligned}
				\sft_2(\phi_\frh(W))\ &=\ \phi_\frg (\sft_1(W))~,
                \\
                \phi_\frg([V_1,V_2]_{\frg_1})\ &=\ [\phi_\frg(V_1),\phi_\frg(V_2)]_{\frg_2}+\sft_2(\phi_2(V_1,V_2))~,
                \\
                \phi_\frh(V_1\acton_1 W)\ &=\ \phi_\frg(V_1)\acton_2 \phi_\frh(W)+\phi_2(V_1,\sft_1(W))~,
                \\
                0\ &=\ \phi_2(V_1,[V_2,V_3])+\phi_\frg(V_1)\acton_2\phi_2(V_2,V_3)+\mbox{cycl.}
			\end{aligned}
		\end{equation}  这是译文 $V_{1,2}\in \frg$ 这是译文 $W\in \frh$; 这是译文.~\cite{Jurco:2018sby,Jurco:2019bvp} 这是译文 $L_\infty$-这是译文. 


        \subsection{这是译文}\label{sec:descentData}
        
这是译文. 
        
        \paragraph{这是译文.}
 这是译文 \uline{这是译文} 这是译文 $\sigma:Y\rightarrow X$ 这是译文         \begin{equation}\label{eq:fibreProducts}
            Y^{[n]}\ \coloneqq\ \underbrace{Y\times_X\cdots\times_XY}_{\text{$n$ factors}}\ \coloneqq\ \{(y_1,\ldots,y_n)\in Y^n\,|\,\sigma(y_1)=\cdots=\sigma(y_n)\}
        \end{equation}  这是译文 $n\in\IN$. 这是译文 \emph{这是译文} 这是译文 $Y^{[n]}$ 这是译文.

        \paragraph{这是译文.} 这是译文 $\caG\coloneqq(\sfH\overset{\sft}{\longrightarrow}\sfG,\acton)$ 这是译文 $\sfLie(\caG)\coloneqq(\frh\overset{\sft}{\longrightarrow}\frg,\acton)$ 这是译文. 
        \begin{definition}
            这是译文 \uline{这是译文 $\caG$-这是译文 (这是译文) 这是译文} 这是译文 $X$ 这是译文 $Y\rightarrow X$ 这是译文             \begin{subequations}\label{eq:unadjustedCocycleConditions}
                \begin{equation}
                    \begin{aligned}
                        h\ &\in\ \scC^\infty(Y^{[3]},\sfH)~,
                        \\
                        (g,\Lambda)\ &\in\ \scC^\infty(Y^{[2]},\sfG)\oplus\Omega^1(Y^{[2]},\frh)~,
                        \\
                        (A,B)\ &\in\ \Omega^1(Y^{[1]},\frg)\oplus\Omega^2(Y^{[1]},\frh)~,
                    \end{aligned}
                \end{equation}
                such that
                \begin{equation}\label{eq:unadjustedCocycleConditionsB}
                    \begin{aligned}
                        h_{ikl}h_{ijk}\ &=\ h_{ijl}(g_{ij}\acton h_{jkl})~,
                        \\
                        g_{ik}\ &=\ \sft(h_{ijk})g_{ij}g_{jk}~,
                        \\
                        \Lambda_{ik}\ &=\ \Lambda_{jk}+g_{jk}^{-1}\acton \Lambda_{ij}-g_{ik}^{-1}\acton(h_{ijk}\nabla_i h_{ijk}^{-1})~,
                        \\
                        A_j\ &=\ g^{-1}_{ij}A_ig_{ij}+g^{-1}_{ij}\rmd g_{ij}-\sft(\Lambda_{ij})~,
                        \\
                        B_j\ &=\ g^{-1}_{ij}\acton B_i+\rmd\Lambda_{ij}+ A_j\acton\Lambda_{ij}+\tfrac12[\Lambda_{ij},\Lambda_{ij}]~,
                        \\
                        0\ &=\ \rmd A_i+\tfrac12[A_i,A_i]+\sft(B_i)
                    \end{aligned}
                \end{equation}
            \end{subequations}  这是译文 $(i,j,\ldots)\in Y^{[n]}$; 这是译文 $\nabla_i\coloneqq\rmd+A_i\acton$.
        \end{definition}
 这是译文~\cite{Breen:math0106083,Aschieri:2003mw} 这是译文         \begin{equation}
            \delta\ \in\ \Omega^2(Y^{[2]},\frh)
		\end{equation}  这是译文\footnote{这是译文.~\cite{Schreiber:2008aa,Saemann:2012uq}} 这是译文.\footnote{这是译文 $\delta$ 这是译文 (这是译文) 这是译文 $\Lambda$.}

 这是译文 $(A,B,\Lambda)$ 这是译文 \v 这是译文 $\caG$-这是译文~\eqref{eq:unadjustedCocycleConditionsB}, 这是译文 \uline{这是译文}. 这是译文 \uline{这是译文} 这是译文         \begin{equation}\label{eq:unadjustedCurvatures}
            F_i\ \coloneqq\ \rmd A_i+\tfrac12[A_i,A_i]+\sft(B_i)
            \eand
            H_i\ \coloneqq\ \rmd B_i+A_i\acton B_i
        \end{equation}  这是译文 $(i)\in Y^{[1]}$. 这是译文~\eqref{eq:unadjustedCocycleConditionsB} 这是译文 $F_i=0$, 这是译文 \uline{这是译文}. 这是译文         \begin{equation}\label{eq:unadjustedBianchiIdentities}
            \nabla_i F_i\ =\ \sft(H_i)\ =\ 0
            \eand
            \nabla_i H_i\ =\ F_i\acton B_i\ =\ 0
        \end{equation}  这是译文.
        \begin{definition}
            这是译文 $\caG$-这是译文 \uline{这是译文} 这是译文 $(h,g,\Lambda,A,B)$ 这是译文 $(\tilde h,\tilde g,\tilde\Lambda,\tilde A,\tilde B)$ 这是译文             \begin{subequations}\label{eq:unadjustedCoboundaryConditions}
                \begin{equation}
                    \begin{aligned}
                        b\ &\in\ \scC^\infty(Y^{[2]},\sfH)~,
                        \\
                        (a,\lambda)\ &\in\ \scC^\infty(Y^{[1]},\sfG)\oplus\Omega^1(Y^{[1]},\frh)
                    \end{aligned}
                \end{equation}
                by means of
                \begin{equation}
                    \begin{aligned}
                        \tilde h_{ijk}\ &=\ a_i^{-1}\acton(b_{ik}h_{ijk}(g_{ij}\acton b_{jk}^{-1})b_{ij}^{-1})~,
                        \\
                        \tilde g_{ij}\ &=\ a_i^{-1}\sft(b_{ij})g_{ij}a_j~,
                        \\
                        \tilde\Lambda_{ij}\ &=\ a^{-1}_j\acton\Lambda_{ij}+\lambda_j-\tilde{g}^{-1}_{ij}\acton\lambda_i+(a_j^{-1}g_{ij}^{-1})\acton(b_{ij}^{-1}\nabla_ib_{ij})~,                    
                        \\
                        \tilde A_i\ &=\ a_i^{-1}A_ia_i+a_i^{-1}\rmd a_i-\sft(\lambda_i)~,
                        \\
                        \tilde B_i\ &=\ a_i^{-1}\acton B_i+\rmd\lambda_i+\tilde{A}_i\acton\lambda_i+\tfrac12[\lambda_i,\lambda_i]
                    \end{aligned}
                \end{equation}
            \end{subequations}  这是译文 $(i,j,\ldots)\in Y^{[n]}$. 这是译文 \uline{这是译文}.
        \end{definition}

\noindent 这是译文~\eqref{eq:unadjustedCurvatures} 这是译文~\eqref{eq:unadjustedCoboundaryConditions} 这是译文         \begin{equation}\label{eq:unadjustedCurvatureGaugeTransformation}
            \tilde F_i\ =\ a_i^{-1}F_ia_i
            \eand
            \tilde H_i\ =\ a_i^{-1}\acton H_i+\tilde F_i\acton\lambda_i
        \end{equation}  这是译文 $(i)\in Y^{[1]}$. 

这是译文,         \begin{equation}
            \begin{tikzcd}
                (h,g,\Lambda,A,B) 
                \arrow[r,bend left=50,"{(a,b,\lambda)}"{name=U}]
                \arrow[r,bend right=50,"{(\tilde a,\tilde b,\tilde\lambda)}"{name=D},swap]
                \arrow[Rightarrow,"\,m", from=U, to=D,start anchor={[yshift=-1ex]},end anchor={[yshift=1ex]}]
                & (\tilde h,\tilde g,\tilde\Lambda,\tilde A,\tilde B)~,
            \end{tikzcd}
        \end{equation}

        \begin{definition} 
            这是译文 $(a,b,\lambda)$ 这是译文 $(\tilde a,\tilde b,\tilde\lambda)$ 这是译文 \uline{这是译文} 这是译文             \begin{subequations}\label{eq:unadjustedHigherGaugeTransformations}
                \begin{equation}
                    m\ \in\ \scC^\infty(Y^{[1]},\sfH)
                \end{equation}
                by means of 
                \begin{equation}\label{eq:unadjustedHigherGaugeTransformationsB}
                    \begin{aligned}
                        \tilde b_{ij}\ &=\ m_ib_{ij}(g_{ij}\acton m_j^{-1})~,
                        \\
                        \tilde a_i\ &=\ \sft(m_i)a_i~,
                        \\
                        \tilde \lambda_i\ &=\ \lambda_i+a_i^{-1}\acton(m_i^{-1}\nabla_i m_i)
                    \end{aligned}
                \end{equation}
            \end{subequations}  这是译文 $(i,j,\ldots)\in Y^{[n]}$. 这是译文 \uline{这是译文}.
        \end{definition}   

        \paragraph{这是译文.}
 这是译文 $F_i=0$ 这是译文~\eqref{eq:unadjustedCocycleConditions}, 这是译文~\cite{Schreiber:0705.0452,Schreiber:2008aa}, 这是译文.~这是译文~\cite{Kim:2019owc}. 这是译文~\eqref{eq:unadjustedCoboundaryConditions} 这是译文,         \begin{subequations}
            \begin{equation}\label{eq:tripleGaugeTransformations}
                \begin{tikzcd}
                    & (A_2,B_2)\arrow[dr,"{(a_{23},\lambda_{23})}"]\arrow[d,Rightarrow,"{m_{123}}"]
                    \\
                    (A_1,B_1) \arrow[ur,"{(a_{12},\lambda_{12})}"] \arrow[rr,"{(a_{13},\lambda_{13})}",swap] & \phantom{a} & (A_3,B_3) 
                \end{tikzcd}
            \end{equation}
            Then, we have the following result.

            \begin{proposition}\label{prop:unadjustedGluingOfGaugeTransformations}
                The gauge transformations~\eqref{eq:tripleGaugeTransformations} are related by a higher coboundary, parametrised by an $m_{123}\in\scC^\infty(Y,\sfH)$, in the following way\footnote{cf.~\eqref{eq:unadjustedCocycleConditions}}
                \begin{equation}\label{eq:unadjustedCoboundaryComposition}
                    \begin{aligned}
                        a_{13}\ &=\ \sft(m_{123})a_{12}a_{23}~,
                        \\
                        \lambda_{13}\ &=\ \lambda_{23}+a_{23}^{-1}\acton\lambda_{12}-a_{13}^{-1}\acton(m_{123}\nabla_1m_{123}^{-1})~.
                    \end{aligned}
                \end{equation}
                The commutativity of the diagram~\eqref{eq:tripleGaugeTransformations} amounts to
                \begin{equation}\label{eq:fakeCurvatureConditionFromTriangle}
                    (a_{23}^{-1}a_{12}^{-1})\acton (m_{123}^{-1}(F_1\acton m_{123}))\ =\ 0~.
                \end{equation}
            \end{proposition}
        \end{subequations}
        \begin{proof}
            这是译文~\eqref{eq:unadjustedCoboundaryComposition} 这是译文 $A$ 这是译文~\eqref{eq:fakeCurvatureConditionFromTriangle} 这是译文 $B$, 这是译文 \cref{app:cocycleConsistency} 这是译文.
        \end{proof}
        
 这是译文~\eqref{eq:unadjustedCocycleConditions} 这是译文~$\acton$ 这是译文.

这是译文. 

        \begin{proposition}[{\cite[Theorem 4.1]{Gastel:2018joi}, \cite[Theorem 4.3]{Saemann:2019dsl}}]
这是译文.
        \end{proposition}

        \begin{remark}
            这是译文 $(\sfG\overset{\sft}{\longrightarrow}\unit,\id)$, 这是译文 $\sfG=\ker(\sft)$ 这是译文 $(\sfG\overset{\sft}{\longrightarrow}\unit,\id)$-这是译文~\cite{Hitchin:1999fh,Chatterjee:1998}, 这是译文~\cite{Murray:9407015} 这是译文~\cite{Murray:2007ps} 这是译文.
        \end{remark}

        \subsection{这是译文}\label{ssec:adjustedCrossedModules}

 这是译文 $L_\infty$-这是译文 $L_\infty$-这是译文 (这是译文) 这是译文~\cite{Sati:2008eg,Fiorenza:2010mh}.\footnote{这是译文 $\rmd H=\tfrac12 p_1$, 这是译文 $H$ 这是译文 $p_1$ 这是译文.} 这是译文 $L_\infty$-这是译文 \uline{这是译文} 这是译文~\cite{Saemann:2019dsl}. 这是译文.~这是译文~\cite{Sati:2008eg,Saemann:2019dsl,Borsten:2021ljb}, 这是译文 $L_\infty$-这是译文~\cite{Sati:2008eg,Saemann:2019dsl}. 

        \paragraph{这是译文.} 这是译文.

        \begin{definition}\label{def:generalAdjustment}  这是译文 \uline{这是译文} $\caG_\kappa=(\sfH\overset{\sft}{\longrightarrow}\sfG,\acton,\kappa)$ 这是译文 $\caG=(\sfH\overset{\sft}{\longrightarrow}\sfG,\acton)$ 这是译文 
            \begin{subequations}\label{eq:alternativeAdjustmentCondition}
                \begin{equation}\label{eq:general_adjustment}
                    \kappa\,:\,\sfG\times\frg\ \rightarrow\ \frh~,
                \end{equation}
                called the \uline{adjustment}, which is linear in $\frg$ and satisfies
                \begin{align}
                    \kappa(\sft(h),V)\ &=\ h(V\acton h^{-1})~,\label{eq:alternativeAdjustmentCondition_a}
                    \\
                    \kappa(g_2g_1,V)\ &=\ g_2\acton\kappa(g_1,V)+\kappa\big(g_2,g_1Vg^{-1}_1-\sft(\kappa(g_1,V))\big)\label{eq:alternativeAdjustmentCondition_b}
                \end{align}
                for all $g_{1,2}\in\sfG$, $h\in\sfH$, and $V\in\frg$.
            \end{subequations}
        \end{definition}

        \begin{proposition}
            这是译文~\eqref{eq:alternativeAdjustmentCondition_a} 这是译文~\eqref{eq:alternativeAdjustmentCondition_b} 这是译文 
            \begin{subequations}
            	\begin{equation}\label{eq:adjustmentNormalisation}
            		\kappa(\unit,V)\ =\ 0~,
            	\end{equation}
            	and
            	\begin{equation}\label{eq:adjustmentCondition}
                \begin{aligned}
                    &(g_2g_1)\acton(h(V\acton h^{-1}))+g_2\acton\kappa(g_1,V)
                    \\
                    &\kern1cm+\kappa(g_2,g_1Vg^{-1}_1-\sft(\kappa(g_1,V)))-\kappa(g_2g_1\sft(h),V)\ =\ 0
                \end{aligned}
            \end{equation}
            for all $g_{1,2}\in\sfG$, for all $h\in\sfH$, and for all $V\in\frg$.
        	\end{subequations}
        \end{proposition}

        \begin{proof}
            这是译文 $g_1=g_2=\unit$ 这是译文~\eqref{eq:adjustmentCondition} 这是译文~\eqref{eq:adjustmentNormalisation}, 这是译文~\eqref{eq:alternativeAdjustmentCondition_a}. 这是译文,~\eqref{eq:alternativeAdjustmentCondition_b} 这是译文~\eqref{eq:adjustmentCondition} 这是译文 $h=\unit$.

这是译文 $g_2=\tilde g_2\tilde g_1$ 这是译文 $g_1=\sft(h)$ 这是译文~\eqref{eq:alternativeAdjustmentCondition_b} 这是译文 $\tilde g_{1,2}\in\sfG$ 这是译文 $h\in\sfH$ 这是译文~\eqref{eq:alternativeAdjustmentCondition_a} 这是译文 $\sft(\kappa(\sft(h),V))=\sft(h(V\acton h^{-1}))=\sft(h)V\sft(h^{-1})-V$, 这是译文~\eqref{eq:adjustmentCondition}. 这是译文 $h=\unit$ 这是译文~\eqref{eq:alternativeAdjustmentCondition_a}, 这是译文~\eqref{eq:adjustmentNormalisation}. 
        \end{proof}
        
        \paragraph{这是译文.} 这是译文 $(\sfH\overset{\sft}{\longrightarrow}\sfG,\acton)$ 这是译文 $\sft$. 这是译文.
        
        \begin{proposition}\label{prop:special_adjustment}  这是译文 $\caG\coloneqq(\sfH\overset{\sft}{\longrightarrow}\sfG,\acton)$ 这是译文 $\sfLie(\caG)\coloneqq(\frh\overset{\sft}{\longrightarrow}\frg,\acton)$. 这是译文             \begin{equation}
				P\,:\,\frg\ \rightarrow\ \frh
			\end{equation}  这是译文             \begin{subequations}
                \begin{equation}\label{eq:P_rels_1}
                    \sft\circ P\circ\sft\ =\ \sft~,
                \end{equation}
                and\footnote{Here, we use the shorthand notation $g[V,g^{-1}]\coloneqq gVg^{-1}-V$.}
                \begin{equation}\label{eq:P_rels_2}
                    \sft(P(g_2\sft(P(g_1[V,g_1^{-1}]))g_2^{-1}))\ =\ g_2\sft(P(g_1[V,g_1^{-1}]))g_2^{-1}
                \end{equation}
            \end{subequations}  这是译文 $g_{1,2}\in\sfG$ 这是译文 $V\in\frg$. 这是译文 
            \begin{equation}
                \kappa_0(g,V)\ =\ P(g[V,g^{-1}])
			\end{equation}  这是译文 \eqref{eq:alternativeAdjustmentCondition} 这是译文 $\sft$ 这是译文 $\kappa$ 这是译文 $\sft\circ \kappa=\sft\circ \kappa_0$ 这是译文 $P$ 这是译文 \uline{这是译文}.
        \end{proposition}
        \begin{proof}
            这是译文             \begin{equation}
				\begin{aligned}
                    \sft(\kappa_0(\sft(h),V))\ &=\ \sft(P(\sft(h)[V,\sft(h)^{-1}]))
                    \\ 
                    \ &=\ \sft(P(\sft(h(V\acton h^{-1}))))
                    \\
                    \ &=\ \sft(h(V\acton h^{-1}))
				\end{aligned}
			\end{equation}  这是译文 $h\in\sfH$ 这是译文 $V\in\frg$, 这是译文 	
            \begin{equation}
				\begin{aligned}
                    \sft(\kappa_0(g_2g_1,V))\ &=\ \sft(P(g_2g_1[V,g_1^{-1}g_2^{-1}]))
                    \\
                    \ &=\ g_2\sft(P(g_1[V,g_1^{-1}]))g_2^{-1}+\sft\Big(P(g_2[g_1Vg^{-1}_1,g^{-1}_2])\Big)
                    \\
                    &\hspace{4cm}-\sft\Big(P(g_2[\sft(P(g_1[V,g_1^{-1}])),g^{-1}_2])\Big)
                    \\
                    \ &=\ \sft\Big(g_2\acton P(g_1[V,g_1^{-1}])+P(g_2[g_1Vg^{-1}_1-\sft(P(g_1[V,g_1^{-1}])),g^{-1}_2])\Big)
                    \\
					\ &=\ \sft\Big( g_2\acton\kappa_0(g_1,V)+\kappa\big(g_2,g_1Vg^{-1}_1-\sft(\kappa_0(g_1,V))\big)\Big)~.
				\end{aligned}
			\end{equation}
        \end{proof}
这是译文 $P$ 这是译文 $\sft\circ P$ 这是译文 $\im(\sft)\subseteq\frg$, 这是译文 $\kappa_0$, 这是译文 $P$, 这是译文.

        \paragraph{这是译文.}
这是译文 (这是译文) \uline{这是译文} 这是译文 \uline{这是译文} 这是译文 
        \begin{equation}\label{eq:basedPathLoopGroups}
            P_0\sfG\ \coloneqq\ \{p\in\scC^\infty([0,1],\sfG)\,|\,p(0)=\unit\}
            \eand
            L_0\sfG\ \coloneqq\ \{p \in P_0\sfG\,|\,p(1)=\unit\}~,
        \end{equation}  这是译文.\footnote{这是译文~\cite{Baez:2005sn}, 这是译文.~\cite{Schreiber:0705.0452}. 这是译文.} 这是译文 $P_0\sfG$ 这是译文 $L_0\sfG$ 这是译文{\'这是译文}这是译文         \begin{equation}\label{eq:endpointEvaluation}
            \begin{aligned}
                \boundary\,:\,P_0\sfG\ &\rightarrow\ \sfG~,
                \\
                p\ &\mapsto\ p(1)
            \end{aligned}
        \end{equation}  这是译文{\'这是译文}这是译文 \uline{这是译文} 这是译文 \uline{这是译文} 这是译文 $\frg$         \begin{equation}
            P_0\frg\ \coloneqq\ \{\gamma\in\scC^\infty([0,1],\frg)\,|\,\gamma(0)=0\}\eand
            L_0\frg\ \coloneqq\ \{\gamma\in P_0\frg\,|\,\gamma(1)=0\}
        \end{equation}  这是译文. 
        
这是译文 $P_0\sfG$ 这是译文 $L_0\sfG$ 这是译文         \begin{equation}\label{eq:definitionLoopLieCrossed}
            \caL\sfG\ \coloneqq\ (L_0\sfG\hookrightarrow P_0\sfG,\Ad)~,
        \end{equation}  这是译文 $\Ad$ 这是译文 $L_0\sfG$ 这是译文 $P_0\sfG$. 这是译文 $\sfLie(\caL\sfG)$ 这是译文 
        \begin{equation}\label{eq:crossedModuleLg}
            \caL\frg\ \coloneqq\ (L_0\frg\hookrightarrow P_0\frg,\ad)~.
        \end{equation}
        
这是译文 $\caL\sfG$ 这是译文 $\sfG$, 这是译文\footnote{`这是译文}
        \begin{equation}\label{eq:definitionLieGroupLieCrossed}
            \sfG_{\rm cm}\ \coloneqq\ (\unit\hookrightarrow\sfG,\id)~,
        \end{equation}  这是译文 $\caL \sfG$ 这是译文 $\sfG_{\rm cm}$ 这是译文\footnote{这是译文 $L_0\sfG$ 这是译文 $P_0\sfG$, 这是译文
        \begin{equation}\label{eq:identification}
            \begin{tikzcd}[ampersand replacement=\&]
                P_0\sfG \arrow[r,->>] \arrow[rr,bend right=30,"\boundary"] \& P_0\sfG/L_0\sfG \arrow[r,"\cong"] \& \sfG~.
            \end{tikzcd}
        \end{equation}
        这是译文 $\caL\sfG$ 这是译文 $\sfG_{\rm cm}$ 这是译文 $P_0\sfG/L_0\sfG$ 这是译文 $\sfG$.}
        \begin{equation}\label{eq:loopGroupButterfly}
            \begin{tikzcd}
                \unit \arrow[dd] \arrow[dr]& & L_0\sfG \arrow[dl] \arrow[dd]
                \\
                & P_0\sfG \arrow[dl,"\boundary",swap] \arrow[dr,"\id"]& 
                \\
                \sfG & & P_0\sfG
            \end{tikzcd}
        \end{equation}

        \begin{proposition}\label{prop:adjustmentPathLoopGroups}  这是译文 $\wp\in\scC^\infty([0,1],\IR)$ 这是译文 $\wp(0)=0$ 这是译文 $\wp(1)=1$. 这是译文 $\caL\sfG$ 这是译文             \begin{subequations}\label{eq:adjustmentPathLoopGroups}
                \begin{equation}
                    \begin{aligned}
                        \kappa\,:\,P_0\sfG\times P_0\frg\ &\rightarrow\ L_0\frg~,
                        \\
                        (g,V)\ &\mapsto\ P(g[V,g^{-1}])
                    \end{aligned}
                \end{equation}
                for all $g\in P_0\sfG$ and for all $V\in P_0\frg$ with
                \begin{equation}\label{eq:Pmap}
                    P\ \coloneqq\ \id-\wp\cdot\boundary~,
                \end{equation}
            \end{subequations}  这是译文 $\id$ 这是译文.
        \end{proposition}

        \begin{proof}
            这是译文 $\sft$ 这是译文 \cref{prop:special_adjustment}. 这是译文 $P=\id-\wp\cdot\boundary$, 这是译文 $\sft\circ P\circ\sft=\sft$ 这是译文 $\boundary\circ\sft=0$. 这是译文 $g\in P_0\sfG$ 这是译文 $W\in P_0\frh$             \begin{equation}
                P(g\sft(W)g^{-1})\ =\ P(\sft(g\acton W))\ =\ \sft(g(\acton W))~,
			\end{equation}  这是译文 $\boundary\circ\sft=0$.
        \end{proof}
        
这是译文 $P$ 这是译文\footnote{这是译文.} 这是译文 $\frg$ 这是译文 $P_0\frg$, 这是译文         \begin{equation}
			\begin{aligned}
                P_0\frg\ &\rightarrow\ \frg\oplus L_0\frg~,
                \\
                V\ &\mapsto\ (P-\sfid)(V)+P(V)\ =\ \wp\cdot\boundary(V)+(\sfid-\wp\cdot\boundary)(V)
			\end{aligned}
		\end{equation}  这是译文 $V\in P_0\frg$.

 这是译文 $\caL\frg$ 这是译文 $\frg$, 这是译文 $\frg_{\rm cm}\coloneqq(0\hookrightarrow\frg,\sfid)$, 这是译文 $P$ 这是译文,         \begin{equation}\label{eq:quasi-iso_Lie_algebras}
            \frg_{\rm cm}\ \xrightarrow{~\cong~}\ \caL\frg~,
		\end{equation}  这是译文.~\cite{Saemann:2019dsl} 这是译文.

        \paragraph{这是译文.} 这是译文. 
        
        \begin{definition}
            这是译文 \uline{这是译文} $\phi$ 这是译文 $\caG_{1,2}$ 这是译文 $\kappa_{1,2}$ 这是译文 $\phi:\caG_1\rightarrow\caG_2$ 这是译文 
            \begin{equation}
                \phi_\frh(\kappa_1(g,V))\ =\ \kappa_2(\phi_\sfG(g),\phi_\frg(V))
            \end{equation}  这是译文 $g\in\sfG$ 这是译文 $V\in\frg$, 这是译文 $\phi_\frg$ 这是译文 $\phi_\frg$ 这是译文 $\phi:\sfLie(\caG_1)\rightarrow\sfLie(\caG_2)$ 这是译文 $\phi$. \uline{这是译文} 这是译文.
        \end{definition}
        
这是译文. 

        \begin{definition}
            这是译文 \uline{(这是译文) 这是译文} $\phi$ 这是译文 $\caG_{1,2}$ 这是译文 $\kappa_{1,2}$ 这是译文 
            \begin{equation}
                \begin{tikzcd}
                    & \caK\arrow[dl,"\phi"'] \arrow[dr,"\psi"]
                    \\
                    \caG_1 & & \caG_2
                \end{tikzcd}
            \end{equation}  这是译文 $\caK$ 这是译文 $\phi$ 这是译文 $\psi$ 这是译文 $\phi$ 这是译文.
        \end{definition}
        
        \begin{example} 
            这是译文~\eqref{eq:loopGroupButterfly} 这是译文 $\sfG_{\rm cm}$ 这是译文~\eqref{eq:definitionLieGroupLieCrossed}, 这是译文 $\caL\sfG$ 这是译文~\eqref{eq:crossedModuleLg} 这是译文~\eqref{eq:adjustmentPathLoopGroups}. 这是译文             \begin{equation}
                \begin{tikzcd}
                    & \caL\sfG \arrow[dl,"\phi"'] \arrow[dr,"\psi"]
                    \\
                    \sfG_{\rm cm} & & \caL\sfG
                \end{tikzcd}
    		\end{equation}  这是译文 $\phi$ 这是译文 $\psi$ 这是译文 $\phi$ 这是译文 $\phi$ 这是译文 $\caL\sfG$ (这是译文 $P$) 这是译文 $\sfG_{\rm cm}$.
        \end{example}
        
        \paragraph{这是译文.}  这是译文 (这是译文) 这是译文 $\kappa$ 这是译文 $\kappa$, 这是译文         \begin{equation}\label{eq:linearisedKappa}
            \kappa(V_1,V_2)\ \coloneqq\ \lim_{\eps\to0}\tfrac1\eps\kappa(\exp(\eps V_1),V_2)
        \end{equation}  这是译文 $V_{1,2}\in\frg$.
        
这是译文.
        \begin{proposition}
            这是译文 $\caG$, 这是译文 $\kappa:\frg\times \frg\rightarrow \frh$ 这是译文 $\sfLie(\caG)\coloneqq(\frh\overset{\sft}{\longrightarrow}\frg,\acton)$ 这是译文 
            \begin{equation}
                \begin{aligned}            
                    \kappa(\sft(W),V)\ &=\ -V\acton W~,
                    \\
                    \kappa([V_1,V_2],V_3)\ &=\ V_1\acton \kappa(V_2,V_3)-V_2\acton \kappa(V_1,V_3)+\kappa(V_1,[V_2,V_3])-\kappa(V_2,[V_1,V_3])
                    \\
                    &\hspace{1cm}-\kappa(V_1,\sft(\kappa(V_2,V_3)))+\kappa(V_2,\sft(\kappa(V_1,V_3)))~.
                \end{aligned}
			\end{equation}  这是译文 $V_{1,2}\in\frg$ 这是译文 $W\in\frh$. 这是译文 $\kappa$, 这是译文 \uline{这是译文}.
        \end{proposition}

        \begin{proof}
            这是译文.
        \end{proof}
        
        \begin{example}
            这是译文~\eqref{eq:adjustmentPathLoopGroups} 这是译文             \begin{equation}\label{eq:adjustmentPathLoopAlgebras}
                \begin{aligned}
                    \kappa\,:\,P_0\frg\times P_0\frg\ &\rightarrow\ L_0\frg~,
                    \\
                    (V_1,V_2)\ &\mapsto\ P([V_1,V_2])
                \end{aligned}
            \end{equation}  这是译文 $V_{1,2}\in P_0\frg$ 这是译文.
        \end{example}

 这是译文 $\sfLie(\caG_{1,2})\coloneqq(\frh_{1,2}\overset{\sft_{1,2}}{\longrightarrow}\frg_{1,2},\acton_{1,2},\kappa_{1,2})$ 这是译文 $\phi:\sfLie(\caG_{1})\rightarrow\sfLie(\caG_2)$ 这是译文 $L_\infty$-这是译文~\eqref{eq:L_infty_alg}. 这是译文~\cite[Proposition 3.4]{Noohi:0910.1818}, 这是译文 $\phi$ 这是译文         \begin{equation}
            \begin{tikzcd}
                \frh_1 \arrow[dd,"\sft_1",swap] \arrow[dr,"\lambda_1"]& & \frh_2 \arrow[dl,"\lambda_2",swap] \arrow[dd,"\sft_2"]
                \\
                & \frh_2\oplus \frg_1 \arrow[dl,"\gamma_1",swap] \arrow[dr,"\gamma_2"]& 
                \\
                \frg_1 & & \frg_2
            \end{tikzcd}
        \end{equation}  这是译文         \begin{equation}
            \gamma_2(W_2,V_1)\ =\ \sft_2(W_2)+\phi_\frg(V_1)
        \end{equation}  这是译文 $W_2\in\frh_2$ 这是译文 $V_1\in\frg_1$. 这是译文~\cite[Remark 8.5]{Noohi:0506313}, 这是译文 
        \begin{equation}
            \begin{tikzcd}
                & \frh_1\oplus\frh_2\rightarrow \frh_2\oplus \frg_1 \arrow[dl,"{\pr_1\rightarrow \pr_2}"'] \arrow[dr,"{\pr_2\rightarrow \gamma_2}"]
                \\
                \frh_1\rightarrow \frg_1 & & \frh_2\rightarrow \frg_2
            \end{tikzcd}
        \end{equation}  这是译文 $\kappa_{1,2}$ 这是译文 $\hat\kappa$ 这是译文 $\frh_{1,2}\rightarrow\frg_{1,2}$ 这是译文 $\frh_1\oplus\frh_2\rightarrow\frh_2\oplus\frg_1$, 这是译文         \begin{equation}
            \begin{aligned}
                \pr_1(\hat\kappa(W_2,V_1;W_2',V_1'))\ &=\ \kappa_1(V_1,V_1')~,
                \\
                \pr_2(\hat\kappa(W_2,V_1;W_2',V_1'))\ &=\ \kappa_2(\sft(W_2)+\phi_\frg(V_1),\sft(W'_2)+\phi_\frg(V'_1))
            \end{aligned}
        \end{equation}  这是译文 $W_1,W'_1\in\frh_1$ 这是译文 $W_2,W'_2\in\frh_2$ 这是译文 $V_1,V'_1\in\frg_1$ 这是译文 $V_2,V'_2\in\frg_2$. 这是译文 $\hat\kappa$ 这是译文 $\kappa_{1,2}$. 
        
        \subsection{这是译文}\label{sec:adjustedDescentData}
        
 这是译文 \uline{这是译文}, 这是译文 \uline{这是译文}, 这是译文~\eqref{eq:unadjustedCocycleConditions} 这是译文~\eqref{eq:unadjustedCoboundaryConditions}, 这是译文 \uline{这是译文}.

        \paragraph{这是译文.}
这是译文 $\caG\coloneqq(\sfH\overset{\sft}{\longrightarrow}\sfG,\acton,\kappa)$ 这是译文 $\sfLie(\caG)\coloneqq(\frh\overset{\sft}{\longrightarrow}\frg,\acton,\kappa)$ 这是译文\footnote{这是译文 $\caG$ 这是译文 $\sfLie(\caG)$ 这是译文.}, 这是译文.
        \begin{definition}\label{def:adjustedDescentData}  这是译文 \uline{这是译文 $\caG$-这是译文} 这是译文 $X$ 这是译文 $Y\rightarrow X$ 这是译文             \begin{subequations}\label{eq:adjustedCocycleConditions}
                \begin{equation}
                    \begin{aligned}
                        h\ &\in\ \scC^\infty(Y^{[3]},\sfH)~,
                        \\
                        (g,\Lambda)\ &\in\ \scC^\infty(Y^{[2]},\sfG)\oplus\Omega^1(Y^{[2]},\frh)~,
                        \\
                        (A,B)\ &\in\ \Omega^1(Y^{[1]},\frg)\oplus\Omega^2(Y^{[1]},\frh)~,
                    \end{aligned}
                \end{equation}
                such that
                \begin{equation}\label{eq:adjustedCocycleConditionsB}
                    \begin{aligned}
                        h_{ikl}h_{ijk}\ &=\ h_{ijl}(g_{ij}\acton h_{jkl})~,
                        \\
                        g_{ik}\ &=\ \sft(h_{ijk})g_{ij}g_{jk}~,
                        \\
                        \Lambda_{ik}\ &=\ \Lambda_{jk}+g_{jk}^{-1}\acton \Lambda_{ij}-g_{ik}^{-1}\acton(h_{ijk}\nabla_i h_{ijk}^{-1})~,
                        \\
                        A_j\ &=\ g^{-1}_{ij}A_ig_{ij}+g^{-1}_{ij}\rmd g_{ij}-\sft(\Lambda_{ij})~,
                        \\
                        B_j\ &=\ g^{-1}_{ij}\acton B_i+\rmd\Lambda_{ij}+ A_j\acton\Lambda_{ij}+\tfrac12[\Lambda_{ij},\Lambda_{ij}]-\kappa\big(g_{ij}^{-1},\rmd A_i+\tfrac12[A_i,A_i]+\sft(B_i)\big)\,,
                    \end{aligned}
                \end{equation}
            \end{subequations}  这是译文 $(i,j,\ldots)\in Y^{[n]}$; 这是译文 $\nabla_i\coloneqq\rmd+A_i\acton$. 
        \end{definition}

这是译文 \uline{这是译文} 这是译文 
        \begin{equation}\label{eq:adjustedCurvatures}
            \begin{aligned}
                F_i\ &\coloneqq\ \rmd A_i+\tfrac12[A_i,A_i]+\sft(B_i)\ \in\ \Omega^2(Y^{[1]},\frg)~,
                \\
                H_i\ &\coloneqq\ \rmd B_i+A_i\acton B_i-\kappa(A_i,F_i)\ \in\ \Omega^3(Y^{[1]},\frh)~,
            \end{aligned}
        \end{equation}  这是译文 $(i)\in Y^{[1]}$, 这是译文.

这是译文~\eqref{eq:adjustedCocycleConditions}, 这是译文 $F=0$, 这是译文~\eqref{eq:unadjustedCocycleConditions}.

        \begin{proposition}
            这是译文~\eqref{eq:adjustedCurvatures} 这是译文 \uline{这是译文}             \begin{equation}\label{eq:adjustedBianchiIdentities}
                \nabla_i F_i\ =\ \sft(H_i+\kappa(A_i,F_i))
                \eand
                \nabla_i H_i\ =\ \kappa(A_i,\sft(H_i))-\kappa(F_i,F_i)~.
            \end{equation}
        \end{proposition}

        \begin{proof}
            这是译文~\eqref{eq:adjustedBianchiIdentities} 这是译文~\eqref{eq:kappat_linear} 这是译文~\eqref{eq:kappa_commutator}, 这是译文             \begin{subequations}
                \begin{equation}
                    \kappa(\sft(B_i),F_i)\ =\ -F_i\acton B_i     
                \end{equation} 
                and
                \begin{equation}
                    A_i\acton\kappa(A_i,F_i)\ =\ \kappa\big(\tfrac12[A_i,A_i],F_i\big)-\kappa\big(A_i,[A_i,F_i]-\sft(\kappa(A_i,F_i))\big)\,,
                \end{equation}
            \end{subequations}  这是译文 $F_i$, 这是译文 $\rmd$ 这是译文 $H_i$ 这是译文~\eqref{eq:adjustedBianchiIdentities}. 
        \end{proof}

        \begin{definition}
            这是译文 $\caG$-这是译文 \uline{这是译文} 这是译文 $(h,g,\Lambda,A,B)$ 这是译文 $(\tilde h,\tilde g,\tilde\Lambda,\tilde A,\tilde B)$ 这是译文             \begin{subequations}\label{eq:adjustedCoboundaryConditions}
                \begin{equation}
                    \begin{aligned}
                        b\ &\in\ \scC^\infty(Y^{[2]},\sfH)~,
                        \\
                        (a,\lambda)\ &\in\ \scC^\infty(Y^{[1]},\sfG)\oplus\Omega^1(Y^{[1]},\frh)
                    \end{aligned}
                \end{equation}
                by means of
                \begin{equation}\label{eq:adjustedCoboundaryConditionsB}
                    \begin{aligned}
                        \tilde h_{ijk}\ &=\ a_i^{-1}\acton(b_{ik}h_{ijk}(g_{ij}\acton b_{jk}^{-1})b_{ij}^{-1})~,
                        \\
                        \tilde g_{ij}\ &=\ a_i^{-1}\sft(b_{ij})g_{ij}a_j~,
                        \\
                        \tilde\Lambda_{ij}\ &=\ a^{-1}_j\acton\Lambda_{ij}+\lambda_j-\tilde{g}^{-1}_{ij}\acton\lambda_i+(a_j^{-1}g_{ij}^{-1})\acton(b_{ij}^{-1}\nabla_ib_{ij})~,                    
                        \\
                        \tilde A_i\ &=\ a_i^{-1}A_ia_i+a_i^{-1}\rmd a_i-\sft(\lambda_i)~,
                        \\
                        \tilde B_i\ &=\ a_i^{-1}\acton B_i+\rmd\lambda_i+\tilde{A}_i\acton\lambda_i+\tfrac12[\lambda_i,\lambda_i]-\kappa\big(a_i^{-1},F_i\big)
                    \end{aligned}
                \end{equation}
            \end{subequations}  这是译文 $(i,j,\ldots)\in Y^{[n]}$. 这是译文 \uline{这是译文}.
        \end{definition}

这是译文~\eqref{eq:unadjustedCurvatureGaugeTransformation} 这是译文.

        \begin{proposition}\label{prop:adjustedCurvatureGaugeTransformation}  这是译文~\eqref{eq:adjustedCurvatures} 这是译文~\eqref{eq:adjustedCoboundaryConditions} 这是译文             \begin{equation}\label{eq:adjustedCurvatureGaugeTransformation}
                \tilde F_i\ =\ a_i^{-1}F_ia_i-\sft(\kappa(a_i^{-1},F_i))
                \eand
                \tilde H_i\ =\ a_i^{-1}\acton H_i-\kappa(a_i^{-1},\sft(H_i))
            \end{equation}  这是译文 $(i)\in Y^{[1]}$. 这是译文~\eqref{eq:adjustmentPathLoopGroups}, 这是译文             \begin{equation}
				\sft(\tilde H_i)\ =\ \sft(H_i)
			\end{equation}  这是译文 $(i)\in Y^{[1]}$.
        \end{proposition}

        \begin{proof}
            这是译文~\eqref{eq:adjustedCoboundaryConditionsB} 这是译文 $\tilde A_i$ 这是译文 $\tilde B_i$ 这是译文~\eqref{eq:adjustedCurvatures}, 这是译文 \cref{app:useful,app:kappaProperties}. 这是译文 $\tilde H_i$ 这是译文             \begin{equation}
            	\tilde H_i+\kappa(\tilde A_i, \tilde F_i)\ =\ a^{-1}_i\acton(H_i+\kappa(A_i,F_i))+(a_i^{-1}F_i a_i)\acton\lambda_i-\tilde\nabla_i\kappa(a_i^{-1},F_i)~,
            \end{equation}  这是译文 $\kappa$ 这是译文~\eqref{eq:dofkappa} 这是译文 $F_i$ 这是译文~\eqref{eq:adjustedBianchiIdentities} 这是译文~\eqref{eq:kappat_linear} 这是译文~\eqref{eq:kappa_galfaX}. 这是译文 
            \begin{equation}
                \sft(\kappa(a_i^{-1},\sft(H_i))=\sft(P(\sft(a_i^{-1}[\sft(H_i),a_i])))=\sft(a_i^{-1}[\sft(H_i),a_i]))~.
			\end{equation}
        \end{proof}

这是译文.

        \begin{definition}
            这是译文 $(a,b,\lambda)$ 这是译文 $(\tilde a,\tilde b,\tilde\lambda)$ 这是译文 \uline{这是译文} 这是译文             \begin{subequations}\label{eq:adjustedHigherGaugeTransformations}
                \begin{equation}
                    m\ \in\ \scC^\infty(Y^{[1]},\sfH)
                \end{equation}
                by means of
                \begin{equation}\label{eq:adjustedHigherGaugeTransformationsB}
                    \begin{aligned}
                        \tilde b_{ij}\ &=\ m_ib_{ij}(g_{ij}\acton m_j^{-1})~,
                        \\
                        \tilde a_i\ &=\ \sft(m_i)a_i~,
                        \\
                        \tilde \lambda_i\ &=\ \lambda_i+a_i^{-1}\acton(m_i^{-1}\nabla_i m_i)
                    \end{aligned}
                \end{equation}
            \end{subequations}  这是译文 $(i,j,\ldots)\in Y^{[n]}$. 这是译文 \uline{这是译文}.
        \end{definition}
        
 这是译文~\eqref{eq:adjustedCocycleConditions} 这是译文~\eqref{eq:unadjustedCocycleConditions} 这是译文~\eqref{eq:adjustedCoboundaryConditions} 这是译文~\eqref{eq:unadjustedCoboundaryConditions} 这是译文~\eqref{eq:unadjustedHigherGaugeTransformations} 这是译文~\eqref{eq:adjustedHigherGaugeTransformations}, 这是译文 $B$ 这是译文.

这是译文~\eqref{eq:adjustedCoboundaryConditions} 这是译文,         \begin{subequations}
            \begin{equation}\label{eq:adjustedTripleGaugeTransformations}
                \begin{tikzcd}
                    & (A_2,B_2)\arrow[dr,"{(a_{23},\lambda_{23})}"]\arrow[d,Rightarrow,"{m_{123}}"]
                    \\
                    (A_1,B_1) \arrow[ur,"{(a_{12},\lambda_{12})}"] \arrow[rr,"{(a_{13},\lambda_{13})}",swap] & \phantom{a} & (A_3,B_3) 
                \end{tikzcd}
            \end{equation}
            Then, \cref{prop:unadjustedGluingOfGaugeTransformations} changes as follows.

            \begin{proposition}
                The gauge transformations~\eqref{eq:adjustedTripleGaugeTransformations} are related by a higher coboundary, parametrised by an $ m_{123}\in\scC^\infty(Y,\sfH) $, in the following way
                \begin{equation}\label{eq:adjustedCoboundaryComposition}
                    \begin{aligned}
                        a_{13}\ &=\ \sft(m_{123})a_{12}a_{23}~,
                        \\
                        \lambda_{13}\ &=\ \lambda_{23}+a_{23}^{-1}\acton\lambda_{12}-a_{13}^{-1}\acton(m_{123}\nabla_1m_{123}^{-1})~.
                    \end{aligned}
                \end{equation}
                The diagram~\eqref{eq:adjustedTripleGaugeTransformations} is commutative if and only if
                \begin{equation}\label{eq:adjustmentConditionFromTriangle}
                    (a_{23}^{-1}a_{12}^{-1})\acton (m_{123}^{-1}(F_1\acton m_{123}))\ =\ \kappa(a^{-1}_{13},F_1)-a_{23}^{-1}\acton\kappa(a^{-1}_{12},F_1)-\kappa(a^{-1}_{23},F_2)~,
                \end{equation}
                which, in turn, is equivalent to the adjustment condition~\eqref{eq:adjustmentCondition}.
            \end{proposition}
        \end{subequations}

        \begin{proof}
            这是译文~\eqref{eq:unadjustedCoboundaryComposition} 这是译文 $A$ 这是译文~\eqref{eq:adjustmentConditionFromTriangle} 这是译文 $B$, 这是译文 \cref{app:cocycleConsistency} 这是译文.
        \end{proof}

\noindent 这是译文 $\kappa$ 这是译文 $F$. 这是译文.
        
        \begin{proposition}
            这是译文 $Y$, 这是译文\footnote{这是译文.~这是译文}, 这是译文~\cite[Definition 4.2]{Saemann:2019dsl}.
        \end{proposition}

        \begin{proof}
            这是译文.
        \end{proof}
        
        \paragraph{这是译文.}  这是译文.~这是译文 $L_\infty$-这是译文.\footnote{这是译文 $\frg\xrightarrow{~\sfid~}\frg$, 这是译文 $*\xrightarrow{~~}*$. 这是译文 $L_\infty$-这是译文 $\Omega^\bullet(Y)$, 这是译文 $Y$, 这是译文 $\geq 2$. 这是译文 $L_\infty$-这是译文 $L_\infty$-这是译文.} 这是译文.

        \begin{definition}
            这是译文 $(\kappa,P)$ 这是译文~\eqref{eq:adjustmentPathLoopGroups} 这是译文 \uline{这是译文} 这是译文             \begin{equation}\label{eq:adjustedFakeCurvatureCondition}
				P(F_i)\ =\ 0
			\end{equation}  这是译文 $i\in Y^{[1]}$.
        \end{definition}
        
 这是译文~\eqref{eq:adjustmentPathLoopGroups} 这是译文 $(\frh\xrightarrow{~\sft~}\frg)$, 这是译文         \begin{equation}
			\frh\ \cong\ \ker(\sft)\oplus\im(P\circ\sft)
            \eand
            \frg\ \cong\ \coker(\sft)\oplus\im(\sft\circ P)~.
		\end{equation}  这是译文 $\ker(\sft)$ 这是译文 $\coker(\sft)$, 这是译文 $B_i$ 这是译文 $\im(P\circ \sft)$ 这是译文 $A_i$ 这是译文 $\im(\sft\circ P)$ 这是译文 $(\sft\circ P)(A_i)$ 这是译文 $A_i$. 这是译文 $P(F_i)=0$ 这是译文 
        \begin{equation}
			(P\circ \sft)(B_i)\ =\ P(\rmd A_i+\tfrac12[A_i,A_i])~,
		\end{equation}  这是译文 $(P\circ\sft)(B_i)$ 这是译文 $P((\id-\sft\circ P)(A_i))$.
        
        \section{这是译文}

这是译文 $\sfG$-这是译文 $\sfG$ 这是译文 $\caL\sfG\cong\sfG_{\rm cm}$, 这是译文~\eqref{eq:definitionLieGroupLieCrossed} 这是译文~\eqref{eq:crossedModuleLg}.\footnote{这是译文 $\sfG_{\rm cm}$, 这是译文~\eqref{eq:fakeCurvatureConditionFromTriangle} 这是译文.}

        \subsection{这是译文}\label{sec:principalAdjustedHigherPrincipal}

 这是译文 $(\caL\sfG,\kappa)$-这是译文 $\kappa$ 这是译文~\eqref{eq:adjustmentPathLoopGroups} 这是译文 \cref{def:adjustedDescentData}. 这是译文~\eqref{eq:adjustedCurvatures} 这是译文         \begin{equation}\label{eq:curvature_LG_kompressed}
            F_i\ =\ \rmd A_i+\tfrac12[A_i,A_i]+B_i
            \eand
            H_i\ =\ (\id-\wp\cdot\boundary)(\rmd F_i)
        \end{equation}  这是译文 $(i)\in Y^{[1]}$ 这是译文~\eqref{eq:adjustedBianchiIdentities} 这是译文         \begin{equation}
            \nabla_i F_i\ =\ H_i+\kappa(A_i,F_i)
            \eand
            \rmd H_i\ =\ 0
        \end{equation}  这是译文 $(i)\in Y^{[1]}$. 这是译文~\eqref{eq:adjustedFakeCurvatureCondition} 这是译文         \begin{equation}\label{eq:adjustedFakeCurvatureConditionPathLoopGroups}
            (\id-\wp\cdot\boundary)(F_i)\ =\ 0
            \quad\Rightarrow\quad
            H_i\ = 0
        \end{equation}  这是译文 $(i)\in Y^{[1]}$. 这是译文.

        \begin{proposition}
            这是译文 $\sfG$-这是译文 $\caL\sfG$-这是译文~\eqref{eq:adjustedFakeCurvatureConditionPathLoopGroups} 这是译文.
        \end{proposition}

        \begin{proof}
            这是译文. 

            \vspace{5pt}
 \noindent \uline{这是译文 $\Rightarrow$ 这是译文:} 这是译文 $\sfG$-这是译文 $X$ 这是译文 $Y\rightarrow X$,             \begin{subequations}\label{eq:descentDataOrdinaryPrincipalBundle}
                \begin{equation}
                    g\ \in\ \scC^\infty(Y^{[2]},\sfG)
                    \eand
                    A\ \in\ \Omega^1(Y^{[1]},\frg)
                \end{equation}
                with
                \begin{equation}
                    g_{ik}\ =\ g_{ij}g_{jk}
                    \eand
                    A_j\ =\ g_{ij}^{-1}A_i g_{ij}+g_{ij}^{-1}\rmd g_{ij}
                \end{equation}
            \end{subequations}  这是译文 $(i,j,\ldots)\in Y^{[n]}$. 这是译文 $\boundary:P_0\sfG\rightarrow \sfG$ 这是译文~\eqref{eq:endpointEvaluation} 这是译文 $\boundary$ 这是译文 $Y^{[2]}$ 这是译文 $X$ 这是译文 $g\in\scC^\infty(Y^{[2]},\sfG)$, 这是译文.~这是译文 $g^\circ\in\scC^\infty(Y^{[2]},P_0\sfG)$ 这是译文 
            \begin{equation}
                \begin{tikzcd}
                    & P_0\sfG \arrow[d,"\boundary"]
                    \\
                    Y^{[2]} \arrow[ru,dotted,"g^\circ"] \arrow[r,"g"] & \sfG
                \end{tikzcd}
            \end{equation}  这是译文.\footnote{这是译文 $Y$ 这是译文, $Y^{[2]}$ 这是译文 (这是译文.~\cite[Lemma 79.1]{Munkres:2014aa}) 这是译文 $Y^{[2]}$.}
                
 这是译文 $g^\circ$ 这是译文 $P_0\sfG$-这是译文 $L_0\sfG$. 这是译文 $\caL\sfG$-这是译文 $h^\circ\in\scC^\infty(Y^{[3]},L_0\sfG)$ 这是译文 
            \begin{subequations}\label{eq:lifted_cocycles_to_LSpin}
                \begin{equation}
                    h^\circ_{ijk}\ \coloneqq\ g^\circ_{ik}(g^\circ_{jk})^{-1}(g^\circ_{ij})^{-1}
                    \eforall(i,j,k)\ \in\ Y^{[3]}~.
                \end{equation}
                This indeed defines a loop in $\sfG$, that is, $\boundary(h^\circ_{ijk})=\unit$, because of the gluing condition in~\eqref{eq:descentDataOrdinaryPrincipalBundle}. It is possible and convenient (but not necessary) to define the lift $g^\circ$ in such a way that
                \begin{equation}
                    g^\circ_{ii}\ =\ \unit
                    \eand
                    g^\circ_{ij}\ =\ (g^\circ_{ji})^{-1}
                \end{equation}
                with $\unit\in P_0\sfG$ the constant path, which also implies the normalisations
                \begin{equation}
                    h^\circ_{iij}\ =\ h^\circ_{ijj}\ =\ h^\circ_{iji}\ =\ \unit~.
                \end{equation}
                For the connection forms, we generalise the usual prescription for transporting Maurer--Cartan elements in an $L_\infty$-algebra along quasi-isomorphisms, cf.~\cite{Saemann:2017rjm}:
                \begin{equation}\label{eq:connection_loop_embedded}
                    A^\circ_i\ \coloneqq\ \wp\cdot A_i
                    \eand
                    B^\circ_i\ \coloneqq\ \tfrac12(\wp-\wp^2)\cdot[A_i,A_i]
                    \eforall
                    (i)\ \in\ Y^{[1]}
                \end{equation}
                so that
                \begin{equation}\label{eq:Lambda-lift}
                    \Lambda_{ij}^\circ\ \coloneqq\ (g_{ij}^\circ)^{-1}A_i^\circ g_{ij}^\circ+(g_{ij}^\circ)^{-1}\rmd g_{ij}^\circ-A_j^\circ
                    \eforall
                    (i,j)\ \in\ Y^{[2]}~.
                \end{equation}
            \end{subequations}  这是译文 $\frg$ 这是译文~\eqref{eq:descentDataOrdinaryPrincipalBundle} 这是译文 $A_i$ 这是译文 $A_j$. 这是译文~\eqref{eq:adjustedCocycleConditions} 这是译文~\eqref{eq:adjustedFakeCurvatureConditionPathLoopGroups} 这是译文 $(h^\circ,g^\circ,\Lambda^\circ,A^\circ,B^\circ)$ 这是译文 $\caL\sfG$-这是译文.

            \vspace{5pt}
 \noindent \uline{这是译文 $\Leftarrow$ 这是译文:}  这是译文~\eqref{eq:endpointEvaluation} 这是译文 $(h^\circ,g^\circ,\Lambda^\circ,A^\circ,B^\circ)$ 这是译文 $\caL\sfG$-这是译文 $X$ 这是译文~\eqref{eq:adjustedFakeCurvatureConditionPathLoopGroups} 这是译文 $(g,A)$ 这是译文 $\sfG$-这是译文             \begin{equation}\label{eq:recover_ordinary_cocycles}
                g_{ij}\ \coloneqq\ \boundary(g_{ij}^\circ)
                \eand
                A_i\ \coloneqq\ \boundary(A_i^\circ)
            \end{equation}  这是译文 $(i,j,\ldots)\in Y^{[n]}$. 这是译文~\eqref{eq:descentDataOrdinaryPrincipalBundle}.

            \vspace{5pt}
 \noindent \uline{这是译文:}  这是译文 $\phi:(g,A)\mapsto(h^\circ,g^\circ,\Lambda^\circ,A^\circ,B^\circ)$ 这是译文 $\psi:(h^\circ,g^\circ,\Lambda^\circ,A^\circ,B^\circ)\mapsto(g,A)$. 这是译文, $\psi\circ\phi=\id$, 这是译文 $\phi\circ\psi\Rightarrow\id$. 这是译文 $(h^\circ,g^\circ,\Lambda^\circ,A^\circ,B^\circ)$ 这是译文             \begin{equation}
                (\tilde h^\circ,\tilde g^\circ,\tilde\Lambda^\circ,\tilde A^\circ,\tilde B^\circ)\ \coloneqq\ (\phi\circ\psi)(h^\circ,g^\circ,\Lambda^\circ,A^\circ,B^\circ)~.
            \end{equation}  这是译文 $(a^\circ,b^\circ,\lambda^\circ)$ 这是译文             \begin{equation}\label{eq:coboundaryAlpha}
                a^\circ_i\ \coloneqq\ \unit~,
                \quad
                b^\circ_{ij}\ \coloneqq\ \tilde g^\circ_{ij}(g^\circ_{ij})^{-1}~,
                \eand
                \lambda^\circ_i\ \coloneqq\ (\id-\wp\cdot\boundary)(A^\circ_i)
            \end{equation}  这是译文 $(i,j,\ldots)\in Y^{[n]}$ 这是译文 $(h^\circ,g^\circ,\Lambda^\circ,A^\circ,B^\circ)$ 这是译文 $(\tilde h^\circ,\tilde g^\circ,\tilde\Lambda^\circ,\tilde A^\circ,\tilde B^\circ)$. 这是译文 $g^\circ_{ij}$ 这是译文 $\tilde g^\circ_{ij}$ 这是译文 $b^\circ_{ij}$. 这是译文 $\tilde h^\circ_{ijk}$ 这是译文 $\tilde g^\circ_{ij}$ 这是译文 $\sft$ 这是译文 $A^\circ_i$ 这是译文 $\tilde A^\circ_i$ 这是译文 $\lambda^\circ_i$, 这是译文 $\tilde \Lambda^\circ_{ij}$. 这是译文~\eqref{eq:adjustedFakeCurvatureConditionPathLoopGroups} 这是译文 $\tilde B^\circ_i$ 这是译文 $\tilde A^\circ_i$ 这是译文 $\sft$ 这是译文~\eqref{eq:adjustedFakeCurvatureConditionPathLoopGroups}, 这是译文~\eqref{eq:coboundaryAlpha} 这是译文 $\phi\circ\psi\Rightarrow\id$. 这是译文:             \begin{subequations}
                \begin{equation}
                    \begin{tikzcd}[row sep=2cm,column sep=3cm]
                        (g,A)\arrow[r,bend left=20,"\phi"]\arrow[d,"a"] & (h^\circ,g^\circ,\Lambda^\circ,A^\circ,B^\circ)\arrow[l,bend left=20,"\psi"{below}]\arrow[d,bend left=30,"{(\tilde a^\circ,\tilde b^\circ,\tilde\lambda^\circ)}"{name=U,right}]
                        \arrow[d,bend right=30,"{(a^\circ,b^\circ,\lambda^\circ)}"{name=D,left}]
                        \\
                        (\tilde g,\tilde A)\arrow[r,bend left=20,"\phi"] & (\tilde h^\circ,\tilde g^\circ,\tilde\Lambda^\circ,\tilde A^\circ,\tilde B^\circ)\arrow[l,bend left=20,"\psi"{below}]\arrow[Rightarrow,"m^\circ", from=D, to=U,start anchor={[xshift=1ex]},end anchor={[xshift=-1ex]}]
                    \end{tikzcd}
                \end{equation}
                Explicitly, $a^\circ$ and $\tilde a^\circ$ are lifts of $a$ such that 
                \begin{equation}\label{eq:boundary_kond_1}
                    \boundary(a^\circ_i)\ =\ a_i\ =\ \boundary(\tilde a_i^\circ)
                    \eforall
                    (i)\ \in\ Y^{[1]}~.
                \end{equation}
                The maps $b^\circ$ and $\tilde b^\circ$ are then determined by the explicit form of $\tilde g^\circ$,
                \begin{equation}
                    b^\circ_{ij}\ \coloneqq\ a^\circ_i\tilde g^\circ_{ij}a^{\circ-1}_jg_{ij}^{\circ-1}
                    \eand
                    \tilde b^\circ_{ij}\ \coloneqq\ \tilde a^\circ_i\tilde g^\circ_{ij}\tilde a^{\circ-1}_j(g_{ij}^\circ)^{-1}
                    \eforall
                    (i,j)\ \in\ Y^{[2]}~,
                \end{equation}
                where the expressions on the right-hand sides are indeed loops due to~\eqref{eq:boundary_kond_1}. Furthermore, comparing $A^\circ$ with $\tilde A^\circ$ fixes the maps $\lambda^\circ$ and $\tilde \lambda^\circ$,
                \begin{equation}
                    \begin{aligned}
                        \lambda^\circ_i\ &\coloneqq\ (a_i^\circ)^{-1}A^\circ_ia^\circ_i+(a_i^\circ)^{-1}\rmd a^\circ_i-\tilde A^\circ_i~,
                        \\
                        \tilde \lambda^\circ_i\ &\coloneqq\ (\tilde a^\circ_i)^{-1} A^\circ_i\tilde a^\circ_i+(\tilde a^\circ_i)^{-1}\rmd\tilde a^\circ_i-\tilde A^\circ_i
                    \end{aligned}
                \end{equation}
            \end{subequations}  这是译文 $(i)\in Y^{[1]}$. 这是译文 $(h^\circ,g^\circ,\Lambda^\circ,A^\circ,B^\circ)$ 这是译文 $(\tilde h^\circ,\tilde g^\circ,\tilde\Lambda^\circ,\tilde A^\circ,\tilde B^\circ)$ 这是译文.
            
 这是译文~\eqref{eq:adjustedHigherGaugeTransformations}. 这是译文 
            \begin{equation}
                m^\circ_i\ \coloneqq\ \tilde a^\circ_i(a^\circ_i)^{-1}
                \eforall
                (i)\ \in\ Y^{[1]}
            \end{equation}  这是译文 $a^\circ$ 这是译文 $\tilde a^\circ$ 这是译文 $b^\circ$ 这是译文 $\tilde b^\circ$ 这是译文 $\lambda^\circ$ 这是译文 $\tilde\lambda^\circ$ 这是译文.
        \end{proof}
        
        \subsection{这是译文 \texorpdfstring{$\caL\sfSpin(4)$}{LSpin(4)}-这是译文}\label{sec:Spin4Bundle}
        
        \paragraph{这是译文.} 这是译文 $\sfG$-这是译文 $\caL\sfG$-这是译文 $\sfSU(2)$-这是译文 $S^4$ 这是译文 $S^4$. 这是译文 $\sfSU(2)$-这是译文 $S^3\hookrightarrow S^7\rightarrow S^4$. 这是译文~这是译文. 
        
这是译文 $\sfSpin(4)\cong \sfSU(2)\times \sfSU(2)$, 这是译文 $\sfSpin(4)\cong \sfSU(2)\times \sfSU(2)$-这是译文 
        \begin{equation}\label{eq:Spin4-bundle}
            \sfSU(2)\times\sfSU(2)\ \hookrightarrow\ P\ \coloneqq\ \sfSpin(5)\ \rightarrow\ \sfSpin(5)/\sfSpin(4)\ \cong\ \sfSO(5)/\sfSO(4)\ \cong\ S^4~.
        \end{equation}  这是译文 $S^4$, 这是译文.~\cite{Dabrowski:1986en}.
        
 这是译文 \cref{app:YMCoset} 这是译文~\eqref{eq:Spin4-bundle} 这是译文 $\sfSU(2)$-这是译文 $\sfSU(2)$-这是译文.
        
 这是译文~\cite{Saemann:2017rjm}, 这是译文.~这是译文.~\cite{Fiorenza:2019usl}, 这是译文.
        
        
        \paragraph{$S^4$ 这是译文 $\sfSpin(5)/\sfSpin(4)$.}
 这是译文~\eqref{eq:Spin4-bundle} 这是译文.~这是译文~\cite[Diagram 24.4]{Porteous:1995eh}, 这是译文 
        \begin{equation}\label{eq:def_Sp2}
            \sfSpin(5)\ \cong\ \sfSp(2)\ =\ \sfU(2,\IH)\ \coloneqq\ \left\{g\in\sfMat(2,\IH)\,|\,g^\dagger g=gg^\dagger=\unit_2\right\},
        \end{equation}  这是译文 $\sfSpin(5)$ 这是译文         \begin{equation}\label{eq:Spin5-elements}
            g\ =\ \begin{pmatrix}
                a & b
                \\
                c & d
            \end{pmatrix},
            \quad
            a,b,c,d\in\IH~,
            \quad
            \begin{array}{ll}
                |a|^2+|b|^2\ =\ 1~,&a\bar c+b\bar d\ =\ 0~,
                \\
                |c|^2+|d|^2\ =\ 1~,&c\bar a+d\bar b\ =\ 0~,
                \\
                |a|^2+|c|^2\ =\ 1~,&\bar a b+\bar c d\ =\ 0~,
                \\
                |b|^2+|d|^2\ =\ 1~,&\bar b a+\bar d c\ =\ 0~.
            \end{array}
        \end{equation}  这是译文 $\sfSpin(5)$ 这是译文 $\sfSU(2)$ 这是译文 $\sfSp(1)\cong\sfU(1,\IH)$, 这是译文         \begin{equation}
            \begin{aligned}
                \sfSU(2)\times\sfSU(2)\ &\hookrightarrow\ \sfSpin(5)~,
                \\
                \left(\frac{e}{|e|},\frac{f}{|f|}\right)\ &\mapsto\ 
                \begin{pmatrix}
                    \frac{e}{|e|} & 0 
                    \\
                    0 & \frac{f}{|f|}
                \end{pmatrix}.
            \end{aligned}
        \end{equation}  这是译文 $\sfSpin(5)\rightarrow\sfSpin(5)/\sfSpin(4)$, 这是译文         \begin{equation}\label{eq:Spin4Action}
            \begin{pmatrix}
                a & b
                \\
                c & d
            \end{pmatrix}\ \sim\ 
            \begin{pmatrix}
                a & b
                \\
                c & d
            \end{pmatrix}
            \begin{pmatrix}
                \frac{e}{|e|} & 0 
                \\
                0 & \frac{f}{|f|}
            \end{pmatrix}
            \eforall
            e,f\ \in\ \IH~.
        \end{equation}  这是译文 $g\in\sfSpin(5)$ 这是译文~\eqref{eq:Spin5-elements}, 这是译文         \begin{equation}
            \begin{pmatrix}
                a & b
                \\
                c & d
            \end{pmatrix}
            \begin{pmatrix}
                1 & 0
                \\
                0 & -1
            \end{pmatrix}
            \begin{pmatrix}
                a & b
                \\
                c & d
            \end{pmatrix}^\dagger
            \ = \
            \begin{pmatrix}
                2|a|^2-1 & 2a\bar c
                \\
                2 c\bar a & 2|a|^2-1
            \end{pmatrix}
        \end{equation}  这是译文~\eqref{eq:Spin4Action} 这是译文 $\sfSpin(4)$ 这是译文 $g$. 这是译文 
        \begin{equation}
            x^1+\rmi x^2+\rmj x^3+\rmk x^4\ \coloneqq\ 2 c\bar a~\eand
            x^5\ \coloneqq\ 2|a|^2-1~,
        \end{equation}  这是译文 $\rmi$, $\rmj$, 这是译文 $\rmk$ 这是译文 $x^1,\ldots,x^5\in\IR$. 这是译文~\eqref{eq:Spin5-elements}, 这是译文 $\sum_{i=1}^5(x^i)^2=1$, 这是译文, $S^4\hookrightarrow\IR^5$. 这是译文 
        \begin{equation}
            \pi\,:\,\sfSpin(5)\ \rightarrow\ S^4~,
        \end{equation}  这是译文. 
        
这是译文 $\sfSpin(5)\rightarrow S^7\cong\sfSpin(5)/\sfSU(2)$ 这是译文 
        \begin{equation}
            \begin{pmatrix}
                a & b
                \\
                c & d
            \end{pmatrix}
            \ \mapsto\ (a,c)~,
        \end{equation}  这是译文 $S^7$ 这是译文 $\IH^2\cong \IR^8$ 这是译文 $\sfSU(2)$ 这是译文 $\sfSpin(4)\cong \sfSU(2)\times \sfSU(2)$.
        
        \paragraph{这是译文 $S^4$.}
 这是译文 $\frspin(5)$ 这是译文         \begin{equation}
            X\ =\ 
            \begin{pmatrix}
                x & z
                \\
                -\bar z & y
            \end{pmatrix}
            \eforall
            x,y,z\ \in\ \IH
            \ewith
            \Re(x)\ =\ \Re(y)\ =\ 0~.
        \end{equation}  这是译文, $\frspin(4)\cong\frsu(2)\oplus\frsu(2)$ 这是译文 $x$ 这是译文 $y$, 这是译文 
        \begin{subequations}
            \begin{equation}\label{eq:split_spin5}
                \frspin(5)\ \cong\ \frspin(4)\oplus\frm~,
                \quad
                \begin{pmatrix}
                    x & z
                    \\
                    -\bar z & y
                \end{pmatrix}
                \ =\
                \begin{pmatrix}
                    x & 0
                    \\
                    0 & y
                \end{pmatrix}
                +
                \begin{pmatrix}
                    0 & z
                    \\
                    -\bar z & 0
                \end{pmatrix}
            \end{equation}
            with the evident relations
            \begin{equation}\label{eq:algebraic structure}
                [\frspin(4),\frspin(4)]\ \subseteq\ \frspin(4)~,
                \quad
                [\frspin(4),\frm]\ \subseteq\ \frm~,
                \eand
                [\frm,\frm]\ \subseteq\ \frspin(4)~.
            \end{equation}
        \end{subequations}  这是译文 \cref{app:YMCoset} 这是译文 $\sfSpin(5)\rightarrow S^4$. 这是译文.
        
        \paragraph{这是译文.}
这是译文 \cref{app:YMCoset}, 这是译文 $\sfSpin(4)$-这是译文 $\sfSpin(5)\rightarrow S^4$ 这是译文 
        \begin{equation}\label{eq:transition_function_1}
            \begin{aligned}
                g\,:\,(\sfSpin(5))^{[2]}\ &\rightarrow\ \sfSU(2)\times\sfSU(2)~,
                \\
                (g_1,g_2)\ &\mapsto\ g_1^{-1}g_2~,
            \end{aligned}
        \end{equation}  这是译文~\eqref{eq:fibreProducts}.
        
 这是译文 $S^4$, 这是译文 $S^4$ 这是译文 $\IR^5$, 这是译文 $U_\pm\coloneqq S^4\setminus\{x^5=\pm1\}$, 这是译文 $U_+\sqcup U_-\to S^4$ 这是译文         \begin{subequations}\label{eq:stereographicProjections}
            \begin{equation}
                \begin{aligned}
                    \pi_\pm\,:\,U_\pm\ &\rightarrow\ \pi_\pm(U_\pm)\ \cong \IH~,
                    \\
                    (x^1,\ldots,x^5)\ &\mapsto\ q_\pm\ \coloneqq\ \frac{1}{1\mp x^5}(x^1\pm\rmi x^2\pm\rmj x^3\pm\rmk x^4)~,
                \end{aligned}
            \end{equation}
            where, as before, $\rmi$, $\rmj$, and $\rmk$ are the standard quaternionic units. Consequently, $q_+q_-=1$ on $\pi_+(U_+\cap U_-)\cong\pi_+(U_+\cap U_-)\cong\IH\setminus\{0\}$.
        \end{subequations}  这是译文         \begin{equation}
            \begin{tikzcd}
                U_+\sqcup U_-\arrow[rr,"\phi"]\arrow[rd] & &\sfSpin(5)\arrow[ld] 
                \\
                & S^4
            \end{tikzcd}
        \end{equation}  这是译文 $\phi$ 这是译文         \begin{equation}\label{eq:comp1}
            \begin{gathered}
                (\pi_\pm^{-1})^*\phi\,:\,
                \pi_\pm(U_\pm)\ \rightarrow\ \sfSpin(5)~,
                \\
                q_+\ \mapsto\ \frac{1}{\sqrt{1+|q_+|^2}}
                \begin{pmatrix}
                    \bar q_+ & 1 
                    \\
                    1 & -q_+
                \end{pmatrix}
                \eand
                q_-\ \mapsto\ \frac{1}{\sqrt{1+|q_-|^2}}
                \begin{pmatrix}
                    1 & q_-
                    \\
                    \bar q_- & -1
                \end{pmatrix}.
            \end{gathered}
        \end{equation}  这是译文 $g_{+-}\coloneqq\phi^*g$ 这是译文~\eqref{eq:transition_function_1} 这是译文 $(U_+\sqcup U_-)^{[2]}$ 这是译文 
        \begin{subequations}\label{eq:instanton_bundle}
            \begin{equation}
                \begin{aligned}
                    (\pi_\pm^{-1})^*g_{+-}|_{\pi_\pm(U_+\cap U_-)}\,:\,\pi_\pm(U_+\cap U_-)\ &\rightarrow\ \sfSU(2)\times\sfSU(2)~,
                    \\
                    q_\pm\ &\mapsto\ \left(\frac{q_+}{|q_+|},\frac{\bar q_+}{|q_+|}\right)\ =\ \left(\frac{\bar q_-}{|q_-|},\frac{q_-}{|q_-|}\right).
                \end{aligned}
            \end{equation}
            This principal bundle can be endowed with a connection given by the gauge potentials $A_\pm$ on $U_\pm$ with
            \begin{equation}
                A_-\ =\ g_{+-}^{-1}A_+g_{+-}+g_{+-}^{-1}\rmd g_{+-}
                \eon
                U_+\cap U_-~.
            \end{equation}
            This connection describes a pair of an elementary instanton with an elementary anti-instanton, which, in quaternionic notation, is given by
            \begin{equation}
                (\pi_\pm^{-1})^*A_\pm\ \coloneqq\ \frac12\left(\frac{q_\pm\rmd\bar q_\pm-\rmd q_\pm\bar q_\pm}{1+|q_\pm|^2},\frac{\bar q_\pm\rmd q_\pm-\rmd\bar q_\pm q_\pm}{1+|q_\pm|^2}\right)
                \eon
                \pi_\pm(U_\pm)~,
            \end{equation}
            cf.~\cite{Atiyah:1979iu}. Consequently, the curvature is
            \begin{equation}
                (\pi_\pm^{-1})^*F_\pm\ =\ \left(\frac{\rmd q_\pm\wedge\rmd\bar q_\pm}{(1+|q_\pm|^2)^2},\frac{\rmd\bar q_\pm\wedge\rmd q_\pm}{(1+|q_\pm|^2)^2}\right)
                \eon
                \pi_\pm(U_\pm)
            \end{equation}
        \end{subequations}  这是译文 (这是译文~这是译文) 这是译文.
        
        \paragraph{这是译文.} 这是译文~\eqref{eq:instanton_bundle} 这是译文 $\sfSpin(4)$-这是译文 $P$ 这是译文 $\caL\sfSpin(4)$-这是译文.\footnote{这是译文~\cite{Popov:2015wsa}.}
 这是译文 $U_+\sqcup U_-$ 这是译文 $P$ 这是译文 $\caL\sfSpin(4)$-这是译文 $g_{+-}$ 这是译文~\eqref{eq:instanton_bundle} 这是译文 $P_0\sfSpin(4)$. 这是译文: $U_+\sqcup U_-\cong S^3\times \IR$, 这是译文 $\pi_3(\sfSpin(3)\times \sfSpin(3))\cong \IZ\times \IZ$, 这是译文 $P_0\sfSpin(4)$ 这是译文 $\pi_3$ 这是译文 $L_0\sfSpin(4)$-这是译文 $U_+\cap U_-\cong S^3\times\IR$, 这是译文 $\pi_3(L_0\sfSpin(4))\cong\pi_4(\sfSpin(4))\cong\pi_4(S^3\times S^3)\cong\IZ_2\oplus\IZ_2$, 这是译文.\footnote{这是译文 $\pi_{n}(\sfG)\cong\pi_{n-1}(L_0\sfG)$ 这是译文 $X$ 这是译文 $p_0\in X$ 这是译文 \uline{这是译文} 这是译文 $(X,p_0)$ 这是译文 $\Sigma X\coloneqq(X\times[0,1])/\!\!\sim$ 这是译文 $(p,0)\sim(p,1)\sim(p_0,t)$ 这是译文 $p\in X$ 这是译文 $t\in[0,1]$. 这是译文 $Y$ 这是译文 $\scC_0(\Sigma X,Y)\cong\scC_0(X,L_0Y)$ 这是译文 $\Sigma S^{n-1}$ 这是译文 $S^n$. 这是译文 $x_0\in S^{n-1}$ 这是译文 $y_0\in S^n$ 这是译文 $x_0$ 这是译文 $S^{n-1}\times\{0,1\}$ 这是译文 $y_0$. 这是译文 $(0,1)$ 这是译文 $\IR$ 这是译文 $S^{n-1}\setminus\{x_0\}$ 这是译文 $\IR^{n-1}$ 这是译文 $x_0$. 这是译文 $\IR^{n-1}\times\IR$ 这是译文 $S^n\setminus\{y_0\}$ 这是译文 $y_0$.\label{fn:reduced_suspension}}
                
 这是译文 $U_+\sqcup U_-$ 这是译文 $\pi_3(\sfSpin(4))\cong\pi_2(L_0\sfSpin(4))$ 这是译文 $S^2\times\IR^2$. 这是译文 $\pi_+(U_+\cap U_-)\cong\pi_-(U_+\cap U_-) \cong\IH\setminus\{0\}$ 这是译文 $\IH$ 这是译文, $\IH\setminus\IR\cong S^2\times\IR^2$. 这是译文         \begin{equation}
            \begin{aligned}
                L_+\ &\coloneqq\ \{q\in\IH\,|\,\Re(q)>0\text{ and }\Im(q)=0\}~,
                \\
                L_-\ &\coloneqq\ \{q\in\IH\,|\,\Re(q)<0\text{ and }\Im(q)=0\}~,
            \end{aligned}
        \end{equation}  这是译文 $L\coloneqq L_-\cup\{0\}\cup L_+$; 这是译文 $L\cong\IR$. 这是译文 $U_1\sqcup U_2\sqcup U_3\rightarrow S^4$ 这是译文\footnote{这是译文.}
        \begin{equation}\label{eq:3patchCover}
            \begin{gathered}
                (U_1,\pi_1)
                \ewith
                U_1\ \coloneqq\ \pi_+^{-1}(\underbrace{\pi_+(U_+)\setminus L_+}_{\cong\,\IH\setminus L_+})
                \eand
                \pi_1\ \coloneqq\ \pi_+|_{U_1}~,
                \\
                (U_2,\pi_2)
                \ewith
                U_2\ \coloneqq\ \pi_+^{-1}(\underbrace{\pi_+(U_+)\setminus L_-}_{\cong\,\IH\setminus L_-})
                \eand
                \pi_2\ \coloneqq\ \pi_+|_{U_2}~,
                \\
                (U_3,\pi_3)
                \ewith
                U_3\ \coloneqq\ U_-
                \eand
                \pi_3\ \coloneqq\ \pi_-~.
            \end{gathered}
        \end{equation}  这是译文 $U_1\cup U_2=U_+$ 这是译文 $U_1\cap U_2\cap U_3\cong S^2\times\IR^2$. 这是译文 $\sfSpin(5)\rightarrow S^4$ 这是译文         \begin{subequations}\label{eq:cocycle_Y_tilde}
            \begin{equation}
                g_{12}\ \coloneqq\ -\unit~,
                \quad
                g_{13}\ \coloneqq\ -g_{+-}~,
                \eand
                g_{23}\ \coloneqq\ g_{+-}~,
            \end{equation}
            where we inserted a constant change of frame between patches $U_1$ and $U_2$. This does not affect the gauge potential, and we have
            \begin{equation}
                \begin{aligned}
                    (\pi_{1,2}^{-1})^*A_{1,2}\ &\coloneqq\ \frac12\left(\frac{q_{1,2}\rmd\bar q_{1,2}-\rmd q_{1,2}\bar q_{1,2}}{1+|q_{1,2}|^2},\frac{\bar q_{1,2}\rmd q_{1,2}-\rmd\bar q_{1,2} q_{1,2}}{1+|q_{1,2}|^2}\right),
                    \\
                    (\pi_3^{-1})^*A_3\ &\coloneqq\ \frac12\left(\frac{q_3\rmd\bar q_3-\rmd q_3\bar q_3}{1+|q_3|^2},\frac{\bar q_3\rmd q_3-\rmd\bar q_3 q_3}{1+|q_3|^2}\right),
                \end{aligned}
            \end{equation}
        \end{subequations}  这是译文 $q_{1,2,3}\in\pi_{1,2,3}(U_{1,2,3})$, 这是译文 $q_{1,2}$ 这是译文 $q_+$ 这是译文 $q_3=q_-$.
        
        \paragraph{这是译文.}  这是译文~\eqref{eq:cocycle_Y_tilde} 这是译文 \cref{sec:adjustedDescentData}. 这是译文~\eqref{eq:3patchCover}, 这是译文:         \begin{subequations}
            \begin{equation}
                g^\circ_{12}\ \coloneqq\ \big(t\mapsto(\rme^{\rmi\pi \wp(t)},\rme^{-\rmi\pi \wp(t)})\big)
                \eand
                g^\circ_{i3}\ \coloneqq\ \big(t\mapsto(u_i(t),\bar u_i(t))\big)
            \end{equation}
            for all $i=1,2$ with
            \begin{equation}\label{eq:cocycles_LSpin_bundle_gui}
                u_i(t)\ \coloneqq\ \left(q_i \mapsto \frac{1-\wp(t)+(-1)^i\wp(t)q_i}{|1-\wp(t)+(-1)^i\wp(t)q_i|}\right)
            \end{equation}
            and for all $t\in[0,1]$, which induces
            \begin{equation}
                h^\circ_{123}\ \coloneqq\ g_{13}^\circ(g_{23}^\circ)^{-1}(g_{12}^\circ)^{-1}~.
            \end{equation}
            Furthermore, the lift of the connection reads as
            \begin{equation}
                \begin{gathered}
                    A^\circ_1\ \coloneqq\ \wp\cdot A_1~,
                    \quad
                    A^\circ_2\ \coloneqq\ \wp\cdot A_2~,
                    \quad
                    A^\circ_3\ \coloneqq\ \wp\cdot A_3~,
                    \\
                    B^\circ_1\ \coloneqq\ \tfrac12(\wp-\wp^2)\cdot[A_1,A_1]~,
                    \quad
                    B^\circ_2\ \coloneqq\ \tfrac12(\wp-\wp^2)\cdot[A_2,A_2]~,
                    \\
                    B^\circ_3\ \coloneqq\ \tfrac12(\wp-\wp^2)\cdot[A_3,A_3]~,
                \end{gathered}
            \end{equation}
            and
            \begin{equation}
                \begin{aligned}
                    \Lambda_{12}^\circ\ &\coloneqq\ \bigg(t\mapsto\wp(t)\cdot(\rme^{-\rmi\pi \wp(t)},\rme^{\rmi\pi \wp(t)}) \left[A_+,(\rme^{\rmi\pi \wp(t)},\rme^{-\rmi\pi \wp(t)})\right]\bigg)\,,
                    \\
                    \Lambda_{i3}^\circ\ &\coloneqq\ \bigg(t\mapsto\frac{1}{1+|q_+|^2}\bigg(\wp(t)|q_+|^2\bar u_i(t)u_i(1)\rmd\big(\bar u_i(1)u_i(t)\big)+|q_+|^2(1-\wp(t))\bar u_i(t)\rmd u_i(t)
                    \\
                    &\kern4.5cm+\bar u_i(t)\rmd u_i(t)-\wp(t)\bar u_i(1)\rmd u_i(1),u_i(t)\leftrightarrow\bar u_i(t)\bigg)\bigg)
                \end{aligned}
            \end{equation}
        \end{subequations}  这是译文 $i=1,2$ 这是译文 $u_i$ 这是译文~\eqref{eq:cocycles_LSpin_bundle_gui} 这是译文 $t\in[0,1]$. 这是译文, $u_i(t)\leftrightarrow\bar u_i(t)$ 这是译文 $u_i(t)$ 这是译文 $\bar u_i(t)$ 这是译文.
        
这是译文~\eqref{eq:recover_ordinary_cocycles}.
        
        \paragraph{这是译文.}
 这是译文 $U_1\sqcup U_2\sqcup U_3$ 这是译文 $\bar U_1\sqcup \bar U_2\rightarrow S^4$ 这是译文 
        \begin{equation}\label{eq:2patchCover}
            \begin{gathered}
                (\bar U_1,\bar\pi_1)
                \ewith
                \bar U_1\ \coloneqq\ \pi_+^{-1}(\pi_+(U_+)\setminus(L_+\cup L_-))
                \eand
                \bar\pi_1\ \coloneqq\ \pi_+|_{\bar U_1}
                \\
                (\bar U_2,\bar\pi_2)
                \ewith
                \bar U_2\ \coloneqq\ U_-
                \eand
                \bar\pi_2\ \coloneqq\ \pi_-~.
            \end{gathered}
        \end{equation}  这是译文 $\bar U_1\cap\bar U_2\cong S^2\times\IR^2$. 这是译文 $\sfSpin(5)\rightarrow S^4$  这是译文         \begin{subequations}\label{eq:cocycle_Y_tilde2}
            \begin{equation}
                g_{12}\ \coloneqq\ g_{+-}\big|_{(\bar U_1\sqcup\bar U_2)^{[2]}}~.
            \end{equation}
            The local connection 1-form reads as
            \begin{equation}
                (\pi_{1,2}^{-1})^*A_{1,2}\ \coloneqq\ \left(\frac{\Im{(q_{1,2}\rmd\bar q_{1,2})}}{1+|q_{1,2}|^2},\frac{\Im{(\bar q_{1,2}\rmd q_{1,2})}}{1+|q_{1,2}|^2}\right),
            \end{equation}
        \end{subequations}  这是译文 $q_{1,2}\in\bar\pi_{1,2}(\bar U_{1,2})$, 这是译文 $q_1$ 这是译文 $q_+$ 这是译文 $q_2=q_-$.
        
 这是译文~\eqref{eq:cocycle_Y_tilde2} 这是译文 \cref{sec:adjustedDescentData}. 这是译文~\eqref{eq:3patchCover}, 这是译文         \begin{subequations}\label{eq:cocycles_LSpin_bundle}
            \begin{equation}
                g^\circ_{12}\ \coloneqq\ (u_{q},u_{\bar q})
                \eand
                g^\circ_{21}\ \coloneqq\ (v_{q},v_{\bar q})~,
            \end{equation}
            where $ q=q_+ $ and for all $ t \in [0,1] $
            \begin{equation}
                \begin{aligned}
                    u_q(t)\ &\coloneqq\ \cos(\theta_q\wp(t))+\frac{\Im(q)}{|\Im(q)|}\sin(\theta_q \wp(t))~,
                    \\
                    v_q(t)\ &\coloneqq\ \cos(\theta'_q\wp(t))+\frac{\Im(q)}{|\Im(q)|}\sin(\theta'_q \wp(t))
                \end{aligned}
            \end{equation}
            with
            \begin{equation}\label{eq:cocycles_LSpin_bundle_theta}
                \theta_q\ \coloneqq\ \frac{\pi}{2}-\arctan\left(\frac{\Re(q)}{|\Im(q)|}\right) \eand \theta'_q\ \coloneqq\ 2\pi - \theta_q ~. 
            \end{equation}
            This induces
            \begin{equation}
                \begin{aligned}
                    h^\circ_{121}\ &\coloneqq\ (g_{12}^\circ g_{21}^\circ)^{-1}
                    \\
                    &\kern2.5pt=\ \left(t\mapsto \left(\cos(2\pi \wp(t))+\frac{\Im(\bar q)}{|\Im(q)|}\sin(2\pi \wp(t)),\cos(2\pi \wp(t))+\frac{\Im(q)}{|\Im(q)|}\sin(2\pi \wp(t))\right)\right),
                    \\
                    h^\circ_{212}\ &\coloneqq\ (g_{21}^\circ g_{12}^\circ)^{-1}\ =\ h^\circ_{121}~,
                \end{aligned}
            \end{equation}
            where the last equality follows because $g^\circ_{12}$ and $g^\circ_{21}$ commute. We note that the above two expressions only depend on $\frac{\Im(q)}{|\Im(q)|}$, which describes the embeddings $S^2\hookrightarrow\IR^3\hookrightarrow\IR^4$, as expected from the abstract discussion involving reduced suspension mentioned previously.
            
            The lift of the connection reads as
            \begin{equation}
                \begin{gathered}\label{eq:cocycles_LSpin_bundle_AB2}
                    A^\circ_1\ \coloneqq\ \wp\cdot A_1~,
                    \quad
                    A^\circ_2\ \coloneqq\ \wp\cdot A_2~,
                    \\
                    B^\circ_1\ \coloneqq\ \tfrac12(\wp-\wp^2)\cdot[A_1,A_1]~,
                    \quad
                    B^\circ_2\ \coloneqq\ \tfrac12(\wp-\wp^2)\cdot[A_2,A_2]~,
                \end{gathered}
            \end{equation}
            and, for all $ t \in [0,1] $, 
            \begin{equation}
                \begin{aligned}
                    \Lambda_{12}^\circ\ &\coloneqq\ \bigg(t\mapsto\frac{1}{1+|q|^2}\bigg(|q|^2\wp(t)\bar u_q(t)u_q(1)\rmd\big(\bar u_q(1)u_q(t)\big)+|q|^2(1-\wp(t))\bar u_q(t)\rmd u_q(t)
                    \\
                    &\kern4.5cm+\bar u_q(t)\rmd u_q(t)-\wp(t)\bar u_q(1)\rmd u_q(1),u_q(t)\leftrightarrow\bar u_q(t)\bigg)\bigg)~.
                \end{aligned}
            \end{equation}
            Explicitly,
            \begin{equation}
                \begin{aligned}
                    \Lambda_{12}^\circ(t)\ &=\ \frac{1}{1+|q|^2}\left(Q_q(t) \frac{\Im\big(\Im(q)\rmd \Im(\bar q)\big)}{2|\Im(q)|^2},Q_{\bar q}(t)\frac{\Im\big(\Im(\bar q)\rmd \Im(q)\big)}{2|\Im(q)|^2}\right)
                    \\
                    &=\ \big(Q_q(t),Q_{\bar q}(t)\big)\frac{\Im\big(\Im(q)\rmd \Im(q)\big)}{2|\Im(q)|^2(1+|q|^2)}~,
                \end{aligned}                    
            \end{equation}   
            where
            \begin{equation}
                Q_q(t)\ \coloneqq\ \big(1+|q|^2+\wp(t)(q^2-|q|^2)\big)\bar u_q^2(t)-\left(1+|q|^2-\wp(t)\left(1-\frac{\bar q^2}{|q|^2}\right)\right).
            \end{equation}  
            Similarly, we find
            \begin{equation}
                \begin{aligned}
                    \Lambda_{21}^\circ\ &\coloneqq\ \bigg(t\mapsto\frac{1}{1+|q|^2}\bigg(\wp(t)\bar{v}_{q}(t)v_q(1)\rmd\big(\bar{v}_{q}(1)v_q(t)\big)+(1-\wp(t))\bar{v}_{q}(t)\rmd v_q(t)
                    \\
                    &\kern4.5cm+|q|^2(\bar{v}_{q}(t)\rmd v_q(t)-\wp(t)\bar{v}_{q}(1)\rmd v_q(1)),~v_q(t)\leftrightarrow\bar{v}_q(t)\bigg)\bigg)
                    \\
                    &\kern2.5pt=\ \bigg(t\mapsto\bigg(Q'_q(t),Q'_{\bar q}(t)\bigg)\frac{\Im\big(\Im(q)\rmd\Im(q)\big)}{2|\Im(q)|^2(1+|q|^2)}\bigg)\,,          
                \end{aligned}
            \end{equation}  
            where
            \begin{equation}
                Q'_q(t)\ \coloneqq\ \left(1+|q|^2-\wp(t)\left(1-\frac{\bar q^2}{|q|^2}\right)\right)\bar{v}_{q}^2(t)-\big(1+|q|^2+\wp(t)(q^2-|q|^2)\big)\,.
            \end{equation}  
        \end{subequations}  这是译文 $A$ 这是译文 $B$ 这是译文~\eqref{eq:cocycles_LSpin_bundle_AB2} 这是译文~\eqref{eq:lifted_cocycles_to_LSpin}. 
        
        \begin{remark}\label{rem:better_cover}  这是译文 $S^4$. 这是译文.
        \end{remark}
        
        \section{这是译文}\label{sec:stringBundlesWithConnection}
        
 这是译文 \cref{app:group_structures} 这是译文 $\sfSpin(4)$-这是译文 $S^4$ 这是译文 \cref{sec:Spin4Bundle} 这是译文.
        
 这是译文~\cite{Bergshoeff:1981um,Chapline:1982ww} 这是译文~\cite{Sati:2008eg,Waldorf:2009uf,Sati:2009ic,Fiorenza:2010mh} 这是译文 $P$ 这是译文 $P$~\cite{Waldorf:2009uf}, 这是译文 $H^3_{\rm D}(X,\IZ)$. 这是译文~\cite{Tellez-Dominguez:2023wwr}, 这是译文.
        
        \subsection{这是译文}\label{sec:Lie2GroupModel}
        
 这是译文 \cref{app:group_structures}, 这是译文 $A_\infty$-这是译文~\cite{Baez:2005sn}, 这是译文~\cite{Carey:1989ck,Murray:2001xq} 这是译文~\cite[Section 4]{Pressley:1988qk} 这是译文~\cite[Section 4]{Mickelsson:1989hp} 这是译文 $\sfG$ 这是译文~\cite{Nikolaus:2011zg} 这是译文 $\caG=(G_1\multirightarrow{2}G_0)$ 这是译文 $\pi:\caG\rightarrow\sfG$ 这是译文 $\sfG$ 这是译文 $\unit\in G_0$ 这是译文 $\sfU(1)$.

        \paragraph{这是译文.}
 这是译文~\eqref{eq:basedPathLoopGroups}. 这是译文 $\sfString(\sfG)$ 这是译文 $\sfG$ 这是译文         \begin{equation}\label{eq:centralExtension2Groups}
            \unit\ \longrightarrow\ \sfB\sfU(1)\ \longrightarrow\ \sfString(\sfG)\ \longrightarrow\ \caL\sfG\ \longrightarrow\ \unit~,
        \end{equation}  这是译文 $\sfB\sfU(1)$ 这是译文 $(\sfU(1)\longrightarrow*,\id)$ 这是译文 $\caL\sfG$ 这是译文~\eqref{eq:definitionLoopLieCrossed}. 这是译文 $\sfString(\sfG)$ 这是译文~\cite{Baez:2005sn}, 这是译文.

        \paragraph{这是译文.}
这是译文 \uline{这是译文}\footnote{这是译文 \cref{app:centralExtensions} 这是译文.}
        \begin{equation}\label{eq:Kac-Moody-group-extension}
            \unit\ \longrightarrow\ \sfU(1)\ \overset{\iota}{\longrightarrow}\ \widehat{L_0\sfG}\ \overset{\pi}{\longrightarrow}\ L_0\sfG\ \longrightarrow\ \unit
        \end{equation}  这是译文 $L_0\sfG$. 这是译文 $\widehat{L_0\sfG}$, 这是译文~\cite{Murray:1987ua}, 这是译文~\cite{Mickelsson:1987:173-183,Murray:2001eu}. 这是译文 $\sfU(1)$-这是译文 $\widehat{L_0\sfG}$ 这是译文\footnote{这是译文 \cref{sec:higherPrincipalBundles}.} $P_0L_0\sfG$. 这是译文 $L_0G\cong P_0L_0\sfG/L_0L_0\sfG$, 这是译文~\eqref{eq:identification}, 这是译文 $\widehat{L_0\sfG}\cong \big(P_0L_0\sfG\times\sfU(1)\big)/\sfN$, 这是译文 $N$ 这是译文: 
        
        \begin{proposition}
            这是译文 $\sfN\subseteq L_0L_0\sfG\times\sfU(1)$ 这是译文             \begin{equation}\label{eq:subgroupN}
                \sfN\ \coloneqq\ \left\{(f,z)\in L_0L_0\sfG\times\sfU(1)\,\middle|\,z=\sfhol^{-1}(f)=\exp\left(\int_{D_f}\omega\right)\right\},
            \end{equation}  这是译文 $D_f$ 这是译文 $L_0\sfG$ 这是译文 $f$, 这是译文 $P_0L_0\sfG\times\sfU(1)$.
        \end{proposition}

        \begin{proof}
            这是译文             \begin{equation}
                \left(f_1,\exp\left(\int_{D_{f_1}}\omega\right)\right)\left(f_2,\exp\left(\int_{D_{f_2}}\omega\right)\right)\ =\ \left(f_1f_2,\exp\left(\int_{D_{f_1f_2}}\omega\right)\right)
            \end{equation}  这是译文 
            \begin{equation}
                c(f_1,f_2)\ =\ \exp\left(-\int_{D_{f_1}}\omega-\int_{D_{f_2}}\omega+\int_{D_{f_1f_2}}\omega\right),
            \end{equation}  这是译文~\eqref{eq:cocyclePLGU} 这是译文 $f_{1,2}\in L_0L_0\sfG$. 

 这是译文 $\sfN$ 这是译文             \begin{equation}
                \begin{aligned}
                    (f_1,z_1)(f_2,\sfhol^{-1}(f_2))(f_1,z_1)^{-1}\ &=\ (f_1f_2 f_1^{-1},\sfhol^{-1}(f_2)\,c(f_1,f_2)\,c(f_1f_2,f_1^{-1})\,c^{-1}(f_1,f_1^{-1}))
                    \\
                    &=\ (f_1f_2 f_1^{-1},\sfhol^{-1}(f_1f_2 f_1^{-1}))
                \end{aligned}
            \end{equation}  这是译文 $f_1\in P_0L_0\sfG$, 这是译文 $f_2\in L_0L_0\sfG$, 这是译文 $z_1\in\sfU(1)$. 这是译文 $\sfhol^{-1}(f_2)$ 这是译文~\eqref{eq:holonomy_shift} 这是译文.
        \end{proof}

 这是译文         \begin{equation}\label{eq:KacMoodyExtensionLoopGroup}
            \begin{tikzcd}
                P_0L_0\sfG\times\sfU(1)\arrow[r,"\hat \boundary"] & \boundary^*\widehat{L_0\sfG}\arrow[r]\arrow[d] & \big(P_0L_0\sfG\times\sfU(1)\big)/\sfN\arrow[d]\arrow[r,"\cong"] & \widehat{L_0\sfG}\arrow[d,"\pi"]
                \\
                & P_0L_0\sfG\arrow[r]\arrow[rr,bend right=20,"\boundary"] & P_0L_0\sfG/L_0L_0\sfG\arrow[r,"\cong"] & L_0\sfG
            \end{tikzcd}
        \end{equation}  这是译文 $\boundary^*\widehat{L_0\sfG}$ 这是译文 $\boundary$ 这是译文         \begin{equation}
            \boundary^*\widehat{L_0\sfG}\ \coloneqq \ P_0L_0\sfG~\times_{L_0\sfG}\widehat{L_0\sfG}~,
        \end{equation}  这是译文 $\hat \boundary$ 这是译文         \begin{equation}
            \hat \boundary (f,z)\ \coloneqq\ (f,\hat f(1)z)
        \end{equation}  这是译文 $(f,z)\in P_0L_0\sfG\times\sfU(1)$, 这是译文 $\hat f$ 这是译文\footnote{这是译文 \cref{app:proofs}} of $f\in P_0L_0\sfG$ to $\widehat{L_0\sfG}$ with $\hat f(0)=\unit\in \widehat{L_0\sfG}$.
        
        \paragraph{这是译文 $\widehat{L_0\sfG}$.}
 这是译文 $\widehat{L_0\sfG}$, 这是译文 $c$, 这是译文.~\cref{app:centralExtensions}. 这是译文, $c$ 这是译文         \begin{equation}\label{eq:group_cocycle}
            c\,:\,P_0L_0\sfG\times P_0L_0\sfG\ \rightarrow\ \sfU(1)
            \ewith 
            c(f_1,f_2)c(f_1f_2,f_3)\ =\ c(f_1,f_2f_3)c(f_2,f_3)
        \end{equation}  这是译文 $f_{1,2,3}\in P_0L_0\sfG$, 这是译文 
        \begin{equation}\label{eq:productOnP0L0G}
            (f_1,z_1)(f_2,z_2)\ \coloneqq\ \big(f_1f_2,z_1z_2c(f_1,f_2)\big)
        \end{equation}  这是译文 $(f_{1,2},z_{1,2})\in P_0L_0\sfG\times\sfU(1)$. 这是译文 $c$ 这是译文 $L_0\frg$ 这是译文 $\frg$ 这是译文 $\sfG$. 
        \begin{equation}\label{eq:Lie_algebra_cocycle}
            \begin{aligned}
                \omega\,:\,L_0\frg\times L_0\frg\ &\rightarrow\ \fru(1)~,
                \\
                (\beta_1,\beta_2)\ &\mapsto\ \frac{\rmi}{2\pi}\int_0^1 \rmd r\,\innerLarge{\beta_1(r)}{\parder[\beta_2(r)]{r}}.
            \end{aligned}
        \end{equation}  这是译文, $\inner{-}{-}$ 这是译文 $\frg$ 这是译文~\cite{Pressley:1988qk}.\footnote{这是译文, $\inner{h_\alpha}{h_\alpha}=2$ 这是译文 $h_\alpha$ 这是译文 $\inner{U}{V}=-\tr(UV)$ 这是译文 $\fru(n)$ 这是译文 $\inner{U}{V}=-\tfrac12\tr(UV)$ 这是译文 $\frso(2n)$ 这是译文.} 这是译文~\eqref{eq:Lie_algebra_cocycle} 这是译文 $\omega\in\Omega^2(L_0\sfG,\fru(1))$ 这是译文 $\widehat{L_0\sfG}\rightarrow L_0\sfG$~\cite{Pressley:1988qk},         \begin{equation}
            \omega_g(g\beta_1,g\beta_2)\ \coloneqq\ \omega(\beta_1,\beta_2)
        \end{equation}  这是译文 $g\in L_0\sfG$ 这是译文 $g\beta_{1,2}\in T_gL_0\sfG$, 这是译文 
        \begin{equation}\label{eq:2FormCurvatureKacMoody}
            \omega_g \ =\ \frac{\rmi}{4\pi}\int_0^1\rmd r\,\innerLarge{\theta_{g(r)}}{\parder[\theta_{g(r)}]{r}},
        \end{equation}  这是译文 $\theta$ 这是译文 $L_0\sfG$. 这是译文 $\omega$. 这是译文 $f_{1,2}\in P_0L_0\sfG$, 这是译文         \begin{equation}\label{eq:triangle}
            \ell(f_1,f_2)\ \coloneqq\ \overline{f_1f_2}\circ f_1(1)f_2\circ f_1\ \in\ L_0L_0\sfG~,
        \end{equation}  这是译文 $\overline{f}$ 这是译文 $f\in P_0L_0\sfG$ 这是译文, $\overline{f}(t,r)\coloneqq f(1-t,r)$. 这是译文.

        \begin{proposition} (\cite{Murray:1987ua})  这是译文 $f_{1,2}\in P_0L_0\sfG$ 这是译文 $D_{\ell(f_1,f_2)}$ 这是译文 $L_0\sfG$ 这是译文 $\ell(f_1,f_2)$ 这是译文~\eqref{eq:triangle}. 这是译文,             \begin{equation}\label{eq:cocyclePLGU}
                \begin{aligned}
                    c(f_1,f_2)\ &\coloneqq\ \sfhol(\ell(f_1,f_2))
                    \\
                    &\coloneqq\ \exp\left(-\int_{D_{\ell(f_1,f_2)}}\omega\right)
                    \\
                    &\,=\ \exp\left(-\frac{\rmi}{2\pi}\int_0^1\rmd r\int_0^1\rmd s\int_0^s \rmd t\,\innerLarge{f_1^{-1}\parder[f_1]{s}}{f_2\left\{\parder{r}\left(f_2^{-1}\parder[f_2]{t}\right)\right\}f_2^{-1}}\right)
                \end{aligned}
            \end{equation}  这是译文 $f_{1,2}\in P_0L_0\sfG$ 这是译文 $f_{1,2}=f_{1,2}(s,r)$ 这是译文 $s,t$ 这是译文 $r$ 这是译文. 
        \end{proposition}
        
 \noindent 这是译文         \begin{equation}\label{eq:helpful_identity}
            \parder{t}\left(\parder[f_2(t,r)]{r}f_2^{-1}(t,r)\right)\ =\ f_2(t,r)\left\{\parder{r}\left(f_2^{-1}(t,r)\parder[f_2(t,r)]{t}\right)\right\}f_2^{-1}(t,r)~,
        \end{equation}  这是译文~\eqref{eq:cocyclePLGU} 这是译文~\cite{Baez:2005sn}         \begin{equation}\label{eq:cocyclePLGUBaez}
            c(f_1,f_2)\ =\ \exp\left(-\frac{\rmi}{2\pi}\int_0^1\rmd r\int_0^1\rmd s\,\innerLarge{f_1^{-1}(s,r)\parder[f_1(s,r)]{s}}{\parder[f_2(s,r)]{r}f_2^{-1}(s,r)}\right),
        \end{equation}  这是译文, $s$ 这是译文 $r$ 这是译文~\eqref{eq:group_cocycle} 这是译文 $c$ 这是译文~\eqref{eq:cocyclePLGUBaez}.
        
        \begin{remark}
            这是译文 $\omega$ 这是译文             \begin{equation}\label{eq:holonomy_shift}
                \sfhol(\ell(f_1,f_2))\ =\ \sfhol(g\ell(f_1,f_2))
            \end{equation}  这是译文 $g\in L_0\sfG$. 这是译文~\eqref{eq:group_cocycle} 这是译文 $\omega$ 这是译文             \begin{equation}
                \begin{aligned}
                    &(\overline{f_1f_2},f_1(1)f_2,f_1)~,~
                    &&(\overline{f_1f_2f_3},f_1(1)f_2f_3,f_1)~,
                    \\
                    &(\overline{f_1f_2f_3},(f_1f_2)(1)f_3,f_1f_2)~,~
                    &&(\overline{f_1(1)f_2f_3},(f_1f_2)(1)f_3,f_1(1)f_2)~,
                \end{aligned}
            \end{equation}  这是译文.~\cite{Murray:1987ua}.
        \end{remark}
        
这是译文~\eqref{eq:productOnP0L0G} 这是译文~\eqref{eq:cocyclePLGUBaez}, 这是译文 $P_0L_0\sfG\times\sfU(1)$ 这是译文 $\big(P_0L_0\sfG\times\sfU(1)\big)/\sfN$. 这是译文         \begin{subequations}\label{eq:tMomorphismString2Group}
            \begin{equation}
                \sft\,:\,P_0L_0\sfG\times\sfU(1)\big)/\sfN\ \rightarrow\ P_0\sfG
            \end{equation}
            via the composition of the evident projection $\big(P_0L_0\sfG\times\sfU(1)\big)/\sfN\rightarrow P_0L_0\sfG/L_0L_0\sfG\rightarrow L_0\sfG$ with the embedding $L_0\sfG\hookrightarrow P_0\sfG$. Explicitly, for all $[(f,z)]\in\big(P_0L_0\sfG\times\sfU(1)\big)/\sfN$ we have
            \begin{equation}
                \sft\big([(f,z)]\big)\ =\ \boundary(f)\ \in\ L_0\sfG~,
            \end{equation}
        \end{subequations}  这是译文. 
        
        \paragraph{这是译文.} 这是译文~\cite{Murray:1987ua}.
        
这是译文~\eqref{eq:cocyclePLGU} 这是译文\footnote{这是译文.} 
        \begin{equation}
            c(f_1,f_2)\ =\ \exp\left(-\frac{\rmi}{2\pi}\int_0^1\rmd r\int_0^1\rmd s\int_0^s \rmd t\,\innerLarge{f_1^{-1}\parder[f_1]{s}}{f_2\left\{\parder{r}\left(f_2^{-1}\parder[f_2]{t}\right)\right\}f_2^{-1}}\right)
        \end{equation}  这是译文~\cite{Murray:1987ua} 这是译文 $f_{1,2}\in P_0L_0\sfG$ 这是译文 $f_1=f_1(s,r)$ 这是译文 $f_2=f_2(t,r)$ 这是译文 $s,t$ 这是译文 $r$ 这是译文~\eqref{eq:helpful_identity}, 这是译文 
        \begin{equation}\label{eq:MurrayBaez}
            c(f_1,f_2)\ =\ \exp\left(-\frac{\rmi}{2\pi}\int_0^1\rmd r\int_0^1\rmd s\,\innerLarge{f_1^{-1}(s,r)\parder[f_1(s,r)]{s}}{\parder[f_2(s,r)]{r}f_2^{-1}(s,r)}\right),
        \end{equation}  这是译文~\cite{Baez:2005sn}. 
        
 这是译文~\cite{Mickelsson:1987:173-183}         \begin{equation}
            \begin{aligned}
                \tilde c(f_1,f_2)\ &\coloneqq\ \exp\left(-\frac{\rmi}{4\pi}\int_0^1\rmd r\int_0^1\rmd s\left\{\innerLarge{f_1^{-1}(s,r)\parder[f_1(s,r)]{s}}{\parder[f_2(s,r)]{r}f_2^{-1}(s,r)}\right.\right.
                \\
                &\kern5cm\left.\left.-\innerLarge{f_1^{-1}(s,r)\parder[f_1(s,r)]{r}}{\parder[f_2(s,r)]{s}f_2^{-1}(s,r)}\right\}\right).
            \end{aligned}
        \end{equation}  这是译文 $c$ 这是译文 $\tilde c$ 这是译文         \begin{equation}\label{eq:coboundaryDifferentGroupCocycles}
            \tilde c(f_1,f_2)\ =\ c(f_1,f_2)(d(f_1f_2))^{-1}d(f_1)d(f_2)~,
        \end{equation}  这是译文.~\eqref{eq:coboundaryCondition}, 这是译文 
        \begin{equation}
            \begin{aligned}
                d\,:\,P_0L_0\sfG\ &\rightarrow\ \sfU(1)~,
                \\
                f\ &\mapsto\ \exp\left(-\frac{\rmi}{4\pi}\int_{0}^{1}\rmd r\int_0^1\rmd s\,\innerLarge{f^{-1}(s,r)\parder[f(s,r)]{s}}{f^{-1}(s,r)\parder[f(s,r)]{r}}\right).
            \end{aligned}
        \end{equation}  这是译文~\cite{Murray:1987ua} 这是译文 $c$ 这是译文 $\tilde c$ 这是译文.

        \paragraph{这是译文 $\caL\sfG$.}  这是译文 $P_0\sfG$ 这是译文 $L_0\sfG$ 这是译文 $\caL\sfG$ 这是译文 $P_0\sfG$ 这是译文 $P_0L_0\sfG\times\sfU(1)\big)/\sfN$. 这是译文 $\xi_g\in\Omega^1(L_0\sfG,\fru(1))$ 这是译文 (这是译文.~\cite{Baez:2005sn} 这是译文~\cite[Section 4]{Pressley:1988qk})         \begin{equation}\label{eq:Baez1Form}
            (\xi_g)_\ell\ \coloneqq\ \frac{\rmi}{2\pi}\int_0^{1}\rmd r\,\innerLarge{\theta_{\ell(r)}}{g(r)^{-1}\parder[g(r)]{r}}
        \end{equation}  这是译文 $g\in P_0\sfG$ 这是译文 $\ell\in L_0\sfG$. 这是译文 $(\Ad_g)^*\omega=\omega+\rmd \xi_g$, 这是译文 \cref{app:proofs}, 这是译文 $h\in L_0L_0\sfG$,         \begin{equation}\label{eq:simp_loops}
            \exp\left(-\int_0^1\rmd s\,\xi_g\left(h^{-1}(s)\parder[h(s)]{s}\right)\right)\ =\ \sfhol(ghg^{-1})\sfhol^{-1}(h)~.
        \end{equation}  这是译文, $h(s)\in L_0\sfG$ 这是译文 $s\in[0,1]$. 

        \begin{proposition}
            这是译文~\cite{Baez:2005sn}             \begin{equation}\label{eq:actionOfP0G}
                \begin{aligned}
                    g\acton(f,z)\ &\coloneqq\ \left(gfg^{-1},z\exp\left(\int_0^1\rmd s\,\xi_g\left(f^{-1}(s)\parder[f(s)]{s}\right)\right)\right)
                    \\
                    &\,=\ \left(gfg^{-1},z\exp\left(\frac{\rmi}{2\pi}\int_0^1\rmd r\int_0^1\rmd s\,\innerLarge{f^{-1}(s,r)\parder[f(s,r)]{s}}{g^{-1}(r)\parder[g(r)]{r}}\right)\right)
                \end{aligned}
            \end{equation}  这是译文 $g\in P_0\sfG$ 这是译文 $(f,z)\in P_0L_0\sfG\times\sfU(1)$ 这是译文 $P_0L_0\sfG\times\sfU(1)$. 这是译文 $\sfN$ 这是译文~\eqref{eq:subgroupN} 这是译文 $P_0\sfG$ 这是译文 $\big(P_0L_0\sfG\times\sfU(1)\big)/\sfN$. 这是译文~\eqref{eq:tMomorphismString2Group} 这是译文~\eqref{eq:actionOfP0G} 这是译文~\eqref{eq:crossedModuleConditions}.
        \end{proposition}

        \begin{proof}
            这是译文~\eqref{eq:cocyclePLGUBaez} 这是译文~\eqref{eq:Baez1Form} 这是译文 $c$ 这是译文 $\xi_g$, 这是译文~\eqref{eq:simp_loops}, 这是译文 $\sfN$. 

这是译文 $\sft(g\acton\hat h)=g\sft(\hat h)g^{-1}$ 这是译文 $g\in P_0\sfG$ 这是译文 $\hat h\in P_0L_0\sfG\times\sfU(1)$. 这是译文 $(\sft(\hat h_1)\acton\hat h_2)\hat h_1\hat h_2^{-1}\hat h_1^{-1}\in\sfN$ 这是译文 $\hat h_{1,2}\in P_0L_0\sfG\times\sfU(1)$ 这是译文 \cref{app:proofs}.
        \end{proof}

        \paragraph{这是译文.} 这是译文.

        \begin{definition} (\cite{Baez:2005sn})  这是译文 \uline{这是译文} $\sfString(\sfG)$ 这是译文             \begin{equation}\label{eq:string2GroupDefinition}
                \sfString(\sfG)\ \coloneqq\ \big(\underbrace{\big(P_0L_0\sfG\times\sfU(1)\big)/\sfN}_{\cong\,\widehat{L_0\sfG}}\overset{\sft}{\longrightarrow}P_0\sfG,\acton\big)~.
            \end{equation}  这是译文 $\sfString(n)$ 这是译文 $\sfString(\sfSpin(n))$.
        \end{definition}

        \paragraph{这是译文.}
 这是译文 $P_0L_0\sfG\times\sfU(1)$ 这是译文 $P_0L_0\frg\oplus\fru(1)$. 这是译文~\eqref{eq:cocyclePLGUBaez},         \begin{equation}\label{eq:P0L0G+R_commutator}
            \begin{aligned}
            [(\gamma_1,q_1),(\gamma_2,q_2)]\ &=\ \bigg([\gamma_1,\gamma_2],-\frac{\rmi}{2\pi}\int_0^1\rmd r\int_0^1\rmd s\,\bigg\{\innerLarge{\parder[\gamma_1(s,r)]{s}}{\parder[\gamma_2(s,r)]{r}}
            \\
            &\kern6cm -\innerLarge{\parder[\gamma_2(s,r)]{s}}{\parder[\gamma_1(s,r)]{r}}\bigg\}\bigg)
            \\
            &=\ \left([\gamma_1,\gamma_2],-\frac{\rmi}{2\pi}\int_0^1\rmd    r\,\innerLarge{\boundary(\gamma_1(s,r))}{\parder[\,\boundary(\gamma_2(s,r))]{r}}\right),
        \end{aligned}
        \end{equation}  这是译文 $\boundary(\gamma_{1,2}(s,r))=\gamma_{1,2}(1,r)$. 这是译文~\eqref{eq:Lie_algebra_cocycle}, 这是译文.\footnote{这是译文 $\frg$-这是译文 $\mu$ 这是译文 $\sfG$-这是译文 $\omega$, 这是译文 $V_{1,2}$ 这是译文 $\omega(V_1^h,V_2^h)=-\mu([V_1^h,V_2^h])$, 这是译文 ${}^h$ 这是译文 $\frg$ 这是译文 $\omega(V_1^h,V_2^h)=[V_1,V_2]^h-[V_1^h,V_2^h]$, 这是译文 $[V_1^h,V_2^h]_H=[V_1,V_2]^h$, 这是译文 ${}_H$ 这是译文 (这是译文). 这是译文 $\hat\sfK$ 这是译文 $\sfK$ 这是译文 $\sfG=\sfU(1)$, 这是译文 $[(V_1,0),(V_2,0)]=([V_1,V_2],0)-(0,\omega(V_1^h,V_2^h))=([V_1,V_2],-\omega(V_1,V_2))$, 这是译文 $\hat\frk=\frk\oplus\fru(1)$ 这是译文.}
        
 这是译文         \begin{equation}\label{eq:derhol}
            \dder{t}\bigg|_0\sfhol(\exp(t\gamma))\ =\ 0
        \end{equation}  这是译文 $\gamma\in L_0L_0\frg$, 这是译文 \cref{app:proofs}, 这是译文 $\frn$ 这是译文 $\sfN$ 这是译文         \begin{equation}
            \frn\ =\ \{(\gamma,0)\in L_0L_0\frg\oplus\fru(1)\}\ \cong\ L_0L_0\frg~.
        \end{equation}  这是译文~\eqref{eq:P0L0G+R_commutator} 这是译文 $ \frn $ 这是译文,         \begin{equation}
            (P_0L_0\frg\oplus\fru(1))/\frn\ \cong\ L_0\frg\oplus\fru(1)~.
        \end{equation}  这是译文~\eqref{eq:Lie_algebra_cocycle} 这是译文~\eqref{eq:actionOfP0G} 这是译文         \begin{equation}
            \alpha\,\acton\,(\gamma,q)\ =\ \left([\alpha,\gamma],\frac{\rmi}{2\pi}\int_0^1\rmd r\,\innerLarge{\parder[\alpha]{r}}{\gamma}\right)
        \end{equation}  这是译文 $\alpha\in P_0\frg$ 这是译文 $(\gamma,q)\in L_0\frg\oplus\fru(1)$. 这是译文~\cite{Baez:2005sn}         \begin{equation}\label{eq:strict_string_Lie_2_algebra}
            \frstring(\frg)\ \coloneqq\ \big(L_0\frg\oplus\fru(1)\overset{\sft}{\longrightarrow}P_0\frg,\acton\big)
        \end{equation}  这是译文, $\sft$ 这是译文 $L_0 \frg$ 这是译文 $P_0\frg$. 
        
        \paragraph{这是译文.}
 这是译文 $L_\infty$-这是译文 
        \begin{equation}
            \frstring^\circ(\frg)\ \coloneqq\ \fru(1)\oplus\frg~,
        \end{equation}  这是译文 $\frstring^\circ(\frg)_{-1}\coloneqq\fru(1)$ 这是译文 $\frstring^\circ(\frg)_0\coloneqq\frg$ 这是译文         \begin{equation}
            \begin{aligned}
                \mu_2\,:\,\frg\times\frg\ &\rightarrow\ \frg~~,
                \\
                (V_1,V_2)\ &\mapsto\ [V_1,V_2]~,
                \\
                \mu_3\,:\,\frg\times\frg\times\frg\ &\rightarrow\ \fru(1)~,
                \\
                (V_1,V_2,V_3)\ &\mapsto\ \rmi\inner{V_1}{[V_2,V_3]}
            \end{aligned}
        \end{equation}  这是译文 $V_{1,2,3}\in\frg$. 这是译文~\eqref{eq:quasi-iso_Lie_algebras} 这是译文 $\frg$, 这是译文 $\caL\frg$ 这是译文 
        \begin{subequations}
            \begin{equation}
                \frstring^\circ(\frg)\ \xrightarrow{~\phi~}\ \frstring(\frg)\ \xrightarrow{~\psi~}\ \frstring^\circ(\frg)~.
            \end{equation}
            Explicitly, we have the chain maps
            \begin{equation}
                \begin{gathered}
                    \begin{tikzcd}
                        \fru(1)\arrow[d,"0"]\arrow[r,"\phi_1"] & L_0\frg\oplus\fru(1)\arrow[d,"\sft"]\arrow[r,"\psi_1"] & \fru(1)\arrow[d,"0"]
                        \\
                        \frg \arrow[r,"\phi_1\coloneqq\cdot_\wp"] & P_0\frg\arrow[r,"\psi_1\coloneqq\boundary"]  & \frg
                    \end{tikzcd}
                \end{gathered}
            \end{equation}
            where in the top row $\phi_1$ and $\psi_1$ are the evident embedding and projection maps. In this case, both $\phi_2$ and $\psi_2$ are non-trivial,
            \begin{equation}
                \begin{aligned}
                    \phi_2\,:\,\frg\times\frg\ &\rightarrow\ L_0\frg~,
                    \\
                    (V_1,V_2)\ &\mapsto\ (\wp-\wp^2)\cdot[V_1,V_2]
                \end{aligned}
            \end{equation}
            for all $V_{1,2}\in\frg$ and
            \begin{equation}
                \begin{aligned}
                    \psi_2\,:\,P_0\frg\times P_0\frg\ &\rightarrow\ \fru(1)~,
                    \\
                    (V_1,V_2)\ &\mapsto\ -\rmi\int_0^1\rmd r\left(\innerLarge{\parder[V_1]{r}}{V_2}-\innerLarge{V_1}{\parder[V_2]{r}}\right)
                \end{aligned}
            \end{equation}
            for all $V_{1,2}\in P_0\frg$.
        \end{subequations}
        
        \begin{remark}
            这是译文, $\sfString(\sfG)$ 这是译文 $\sfG$. 这是译文 $\caL\sfG$ 这是译文 $\sfB\sfU(1)$, 这是译文 $(\sfU(1)\rightarrow *,\sfid)$. 这是译文 $H^3(\sfG,\IZ)$, 这是译文 $H^3(\sfG,\IZ)\cong \IZ$ 这是译文 $+1$. 这是译文 $\sfSpin(n)$ 这是译文 $n=3$ 这是译文 $n\geq 5$ 这是译文 $\sfString(n)$. 这是译文 $n=4$, 这是译文 $H^3(\sfSpin(4),\IZ)\cong \IZ\times \IZ$, 这是译文 $(\pm 1,\mp 1)$ 这是译文 $\sfString(4)$.
        \end{remark}

        \subsection{这是译文}\label{sec:adjustedDescentDataStringGroup}
        
这是译文 $\sfString(\sfG)$ 这是译文 $\sfG$. 这是译文 $\frstring(\frg)$ 这是译文~\cite{Saemann:2017rjm,Saemann:2019dsl}.

        \paragraph{这是译文 $\sfString(\sfG)$.}
这是译文 $\kappa$ 这是译文~\eqref{eq:adjustmentPathLoopGroups}. 这是译文 $\caL\sfG$ 这是译文 $\sfString(\sfG)$ 这是译文.

        \begin{proposition}
            这是译文 
            \begin{equation}\label{eq:adjustmentStringGroup}
                \begin{aligned}
                    \hat\kappa\,:\,P_0\sfG\times P_0\frg\ &\rightarrow\ L_0\frg\oplus\fru(1)~,
                    \\
                    (g,V)\ &\mapsto\ \left(\kappa(g,V),\frac{\rmi}{2\pi}\int_0^1\rmd r\,\innerLarge{g^{-1}\parder[g]{r}}{V}\right).
                \end{aligned}
            \end{equation}  这是译文 $g\in P\sfG$ 这是译文 $V\in P_0\frg$ 这是译文 $\sfString(\sfG)$. 这是译文 
            \begin{equation}\label{eq:adjustmentStringAlgebra}
                \begin{aligned}
                    \hat\kappa\,:\,P_0\frg\times P_0\frg\ &\rightarrow L_0\frg\oplus\fru(1)~,
                    \\
                    (V_1,V_2)\ &\mapsto\ \left(\kappa(V_1,V_2),\frac{\rmi}{2\pi}\int_0^1\rmd r\,\innerLarge{\parder[V_1]{r}}{V_2}\right)
                \end{aligned}
            \end{equation}  这是译文 $V_{1,2}\in P_0\frg$ 这是译文 $\kappa(V_1,V_2)$ 这是译文~\eqref{eq:adjustmentPathLoopAlgebras}. 这是译文~\cite{Saemann:2017rjm,Saemann:2019dsl}. 
        \end{proposition}

        \begin{proof}
            这是译文 \cref{app:kappaProperties,app:proofs}. 这是译文,~\eqref{eq:alternativeAdjustmentCondition_a} 这是译文~\eqref{eq:hAhinv}. 这是译文~\eqref{eq:StringkappaIdentities} 这是译文~\eqref{eq:StringkappaIdentity_b}, 这是译文~\eqref{eq:alternativeAdjustmentCondition_b}. 这是译文.
        \end{proof}

 这是译文 \cref{def:adjustedDescentData}, 这是译文 $\sfString(\sfG)$ 这是译文~\eqref{eq:adjustmentStringGroup}, 这是译文         \begin{subequations}
            \begin{equation}
                \begin{aligned}
                    \hat h=[(h,h')]\ &\in\ \scC^\infty\Big(Y^{[3]},\big(P_0L_0\sfG\times\sfU(1)\big)/\sfN\Big)~,
                    \\
                    (g,\hat\Lambda=(\Lambda,\Lambda'))\ &\in\ \scC^\infty(Y^{[2]},P_0\sfG)~\oplus~\Omega^1(Y^{[2]},L_0\frg\oplus\fru(1))~,
                    \\
                    (A,\hat B=(B,B'))\ &\in\ \Omega^1(Y^{[1]},P_0\frg)~\oplus~\Omega^2(Y^{[1]},L_0\frg\oplus\fru(1))~,
                \end{aligned}
            \end{equation}
            such that for all appropriate $(i,j,\ldots)\in Y^{[n]}$, with $Y^{[n]}$ given in~\eqref{eq:fibreProducts}, we have
            \begin{equation}
                \begin{gathered}
                    g_{ik}\ =\ \sft(h_{ijk})g_{ij}g_{jk}~,
                    \\
                    (h_{ikl}h_{ijk})^{-1}h_{ijl}(g_{ij}\acton h_{jkl})\ \in\ L_0L_0\sfG~,
                    \\
                    h'_{ikl}h'_{ijk}\ =\ h'_{ijl}h'_{jkl} \exp\left(\frac{\rmi}{2\pi}\int_0^1\rmd s\int_0^1\rmd r\,\innerLarge{g_{ij}^{-1}\parder[g_{ij}]{r}}{h_{jkl}^{-1}\parder[h_{jkl}]{s}}\right)\exp\left(-\int_\Box\omega\right),
                \end{gathered}
            \end{equation}
            where $\Box$ denotes the loop $\overline{h_{ikl}}\circ\big(\boundary(h_{ikl})\overline{h_{ijk}}\big)\circ\big(\boundary(h_{ijl})(g_{ij}\acton h_{jkl})\big)\circ h_{ijl}$. The cocycle conditions for the differential refinement specialise to 
            \begin{equation}
                \begin{aligned}
                    \Lambda_{ik}\ &=\ \Lambda_{jk}+g_{jk}^{-1}\acton\Lambda_{ij}-g_{ik}^{-1}\acton\big(\boundary(h_{ijk})\nabla_i\boundary(h_{ijk}^{-1})\big)\,,
                    \\
                    \Lambda'_{ik}\ &=\ \Lambda'_{jk}+\Lambda'_{ij}-h'_{ijk}\rmd h'^{-1}_{ijk}
                    \\
                    &\kern1cm-\frac{\rmi}{2\pi}\int_0^1\rmd r\,\left\{\innerLarge{\parder[g_{jk}]{r}g_{jk}^{-1}}{\Lambda_{ij}}-\innerLarge{\parder[g_{ik}]{r}g_{ik}^{-1}}{\boundary(h_{ijk})\nabla_i\boundary(h_{ijk}^{-1})}\right.
                    \\
                    &\kern2cm+\left.\innerLarge{\boundary(h_{ijk}^{-1})\parder[\,\boundary(h_{ijk})]{r}}{A_i}+\int_{0}^{1}\rmd s\,\innerLarge{\parder[h_{ijk}]{r}h_{ijk}^{-1}}{\parder[(h_{ijk}\rmd h_{ijk}^{-1})]{s}}\right\},\\
                    A_j\ &=\ g^{-1}_{ij}A_ig_{ij}+g^{-1}_{ij}\rmd g_{ij}- \Lambda_{ij}~,\\
                    B_j\ &=\ B_i+\rmd\Lambda_{ij}+A_j\acton \Lambda_{ij}+\tfrac12[\Lambda_{ij},\Lambda_{ij}]-\kappa(g^{-1}_{ij},F_i-B_i)~,\\
                    B'_j\ &=\ B'_i+\rmd\Lambda'_{ij} +\frac{\rmi}{2\pi}\int_0^1\rmd r\left\{\innerLarge{\parder[A_j]{r}}{\Lambda_{ij}}+\frac12\innerLarge{\parder[\Lambda_{ij}]{r}}{\Lambda_{ij}}+\innerLarge{\parder[g_{ij}]{r}g_{ij}^{-1}}{F_i-B_i}\right\}.
                \end{aligned}
            \end{equation}
        \end{subequations}  这是译文~\eqref{eq:adjustmentStringAlgebra}, 这是译文~\eqref{eq:adjustedCurvatures} 这是译文         \begin{subequations}\label{eq:curvature_LG_compressed_extended}
            \begin{equation}
                \begin{aligned}
                    F_i\ &=\ \rmd A_i+\tfrac12[A_i,A_i]+\sft(\hat B_i)~,
                    \\
                    H_i\ &=\ \rmd B_i+A_i\acton B_i-\kappa(A_i,F_i)\ =\ (\id-\wp\cdot\boundary)(\rmd F_i)~,
                    \\
                    H'_i\ &=\ \rmd B'_i-\frac{\rmi}{2\pi}\int_0^1\rmd r\,\innerLarge{\parder[A_i]{r}}{F_i-\sft(B_i)}\ =\ \rmd\big(B'_i-\tfrac12\omega(A_i,A_i)\big)-\tfrac{\rmi}{4\pi}{\rm cs}(\boundary(A_i))
                \end{aligned}
            \end{equation}
            for all $(i)\in Y^{[1]}$. Here, in analogy to writing $\hat B=(B,B')$ we did write for the 3-form curvature $\hat H=(H,H')$, and we have made us of the Chern--Simons 3-form
            \begin{equation}\label{eq:chern-simons}
                {\rm cs}(A_i)\ \coloneqq\ \inner{A_i}{\rmd A_i}+\tfrac13\inner{A_i}{[A_i,A_i]}~
            \end{equation}
        \end{subequations}  这是译文 $(i)\in Y^{[1]}$.  这是译文~\eqref{eq:adjustedBianchiIdentities} 这是译文         \begin{equation}\label{eq:Bianchi_identities_extended}
            \begin{aligned}
                \nabla_iF_i\ &=\ \sft(H_i+\kappa(A_i,F_i))~,
                \\
                \rmd H_i\ &=\ 0~,
                \\
                \rmd H'_i\ &=\ -\tfrac{\rmi}{4\pi}\inner{\boundary(F_i-\sft(B_i))}{\boundary(F_i-\sft(B_i))}\ =\ -\tfrac{\rmi}{4\pi}\inner{\boundary(F_i)}{\boundary(F_i)}
            \end{aligned}
        \end{equation}  这是译文 $(i)\in Y^{[1]}$. 
        
        \subsection{这是译文 \texorpdfstring{$S^4$}{S4}}\label{sec:instantonAntiInstantonStringLift}
        
 这是译文 $\sfSpin(n)$ 这是译文 $\sfString(n)$, 这是译文         \begin{equation}
            \sfSpin(5)\ \rightarrow\ \sfSpin(5)/\sfSpin(4)~~~\mbox{to}~~~\sfString(5)\ \rightarrow\ \sfString(5)/\sfString(4)~.
		\end{equation}  这是译文~\cite{Roberts:2022wwl}. 这是译文 $S^4$, 这是译文 $S^4$ 这是译文 $(S^4\multirightarrow{2}S^4)$, 这是译文.
        
        \paragraph{这是译文.}
 这是译文.~\cite{Brown:1976:296-302,Baez:0307200,Porst:0812.1464}. 这是译文 $\caG=(\sfH\overset{\sft}{\longrightarrow}\sfG,\acton)$, 这是译文 $\caG$ 这是译文 $\grp{\caG}$ 这是译文         \begin{equation}\label{eq:crossedModuleGroupoid}
            \begin{gathered}
                \grp{\caG}\ \coloneqq\ (\sfG\ltimes\sfH\multirightarrow{2}\sfG)
                \ewith
                \kern-.8cm
                \begin{tikzcd}[column sep=2.0cm,row sep=large]
                    \phantom{\sft(h)}g & \sft(h^{-1})\,g\arrow[l,bend left,swap,out=-20,in=200]{}{(g,h)}
                \end{tikzcd}\!,
                \\
                (g_1,h_1)\circ(\sft(h^{-1}_1)g_1,h_2)\ =\ (g_1,h_1h_2)~,
                \\
                (g_1,h_1)\otimes (g_2,h_2)\ \coloneqq\ (g_1g_2,(g_1\acton h_2)h_1)
            \end{gathered}
        \end{equation}  这是译文 $g,g_{1,2}\in\sfG$ 这是译文 $h,h_{1,2}\in\sfH$. 这是译文~\cite{Jurco:2014mva}, 这是译文~\cite{Baez:0307200}.
        
 这是译文 \uline{这是译文} 这是译文 $\caG_1\coloneqq(\sfH_1\overset{\sft_1}{\longrightarrow}\sfG_1,\acton_1)$ 这是译文 $\caG_2\coloneqq(\sfH_2\overset{\sft_2}{\longrightarrow}\sfG_2,\acton_2)$, 这是译文 (这是译文) 这是译文 $\sfe:\caG_1\hookrightarrow\caG_2$ 这是译文 $\sfe_0:\sfG_1\hookrightarrow\sfG_2$ 这是译文 $\sfe_1:\sfH_1\hookrightarrow\sfH_2$ 这是译文         \begin{equation}\label{eq:consistency_2_group_embedding}
            \sfe_0(\sft_1(h_1))\ =\ \sft_2(\sfe_1(h_1))
            \eand 
            \sfe_0(g_1)\acton_2\sfe_1(h_1)\ =\ \sfe_1(g_1\acton_1 h_1)~.
        \end{equation}  这是译文 $h_1\in\sfH_1$ 这是译文 $g_1\in\sfG$. 这是译文 $\grp{\caG_1}$ 这是译文 $\grp{\caG_2}$ 这是译文         \begin{equation}
                \caR_{(g_1,h_1)}(g_2,h_2)\ \coloneqq\ (g_2,h_2)\otimes (\sfe_0(g_1),\sfe_1(h_1))\ =\ \big(g_2 \sfe_0(g_1),(g_2\acton_2\sfe_1(h_1))h_2\big)
        \end{equation}  这是译文 $(g_1,h_1)\in\sfG_1\ltimes\sfH_1$ 这是译文 $(g_2,h_2)\in\sfG_2\ltimes\sfH_2$. 这是译文 \uline{这是译文}\footnote{这是译文.}
        \begin{equation}\label{eq:quotionLieGroupoid}
            \caG_2/\caG_1\ \coloneqq\ \grp{\caG_2}/\grp{\caG_1}\ \coloneqq\ \Big(\big(\sfG_2\ltimes\sfH_2)/(\sfG_1\ltimes\sfH_1)\multirightarrow{2}\sfG_2/\sfG_1\Big)~,
        \end{equation}  这是译文         \begin{equation}
            g_2\ \sim\ g_2\sfe_0(g_1)
            \eand
            (g_2,h_2)\ \sim\ (g_2 \sfe_0(g_1),(g_2\acton_2\sfe_1(h_1))h_2)
        \end{equation}  这是译文 $(g_1,h_1)\in\sfG_1\ltimes\sfH_1$ 这是译文 $(g_2,h_2)\in\sfG_2\ltimes\sfH_2$. 这是译文 $\circ$ 这是译文~\eqref{eq:crossedModuleGroupoid}.
        
        \begin{example}\label{ex:bu1_quotient}  这是译文 $\sfB\sfU(1)=(\sfU(1)\longrightarrow*,\id)$ 这是译文 $\sfString(\sfG)$ 这是译文~\eqref{eq:string2GroupDefinition}. 这是译文 $\sfB\sfU(1)$ 这是译文 $\sfString(\sfG)$ 这是译文:             \begin{subequations}\label{eq:first_coset_example}
                \begin{equation}
                    \begin{tikzcd}
                        \sfU(1)\arrow[r,"\sfe_1"]\arrow[d] & \widehat{L_0\sfG}\arrow[d]
                        \\
                        *\arrow[r,"\sfe_0"] & P_0\sfG
                    \end{tikzcd}
                    \quad
                    \begin{array}{c}
                        {}
                        \\[-5pt]
                        \sfe_1\,:\,\sfU(1)\ \ni\ z\ \mapsto\ [(\unit,z)]\ \in\ \big(P_0L_0\sfG\times\sfU(1)\big)/\sfN\ \cong\ \widehat{L_0\sfG}~,
                        \\
                        \kern-4pt\sfe_0\,:\,*\ \ni\ \unit\ \mapsto\ \unit\ \in\ P_0\sfG~.    
                    \end{array}
                \end{equation}
                The resulting quotient is
                \begin{equation}\label{eq:cosetStringBU}
                    \sfString(\sfG)/\sfB\sfU(1)\ \cong\ \caL\sfG
                \end{equation}
            \end{subequations}  这是译文 $\caL\sfG$ 这是译文~\eqref{eq:definitionLoopLieCrossed}. 这是译文~\eqref{eq:centralExtension2Groups}. 这是译文 $\sfG$ 这是译文 $(*\hookrightarrow\sfG,\id)$ 这是译文 $\caL\sfG$ 这是译文~\eqref{eq:loopGroupButterfly}, 这是译文 $\sfG$ 这是译文 $\sfString(\sfG)/\sfB\sfU(1)$.
        \end{example}
        
这是译文 $\sfString(1)$ 这是译文 $\sfB\sfU(1)$, 这是译文 \cref{ex:bu1_quotient} 这是译文 $\sfString(\sfG)/\sfString(1)$. 这是译文.
        
这是译文 \cref{ex:bu1_quotient} 这是译文 $S^1\hookrightarrow S^3\rightarrow S^2$ 这是译文         \begin{equation}\label{eq:categorified_Hopf}
			\sfString(1)\ \hookrightarrow\ \sfString(3)\ \rightarrow\ \caL\sfSpin(3) \ \cong\ \sfSpin(3)~.
		\end{equation}  这是译文.~\cite{Saemann:2017rjm} 这是译文.

        \paragraph{这是译文.} 这是译文 $\sfString(\sfH)$ 这是译文 $\sfString(\sfG)$ 这是译文~\eqref{eq:string2GroupDefinition} 这是译文 $\sfH$ 这是译文 $\sfG$ 这是译文 $\sfH\hookrightarrow\sfG$. 这是译文 $\sfString(\sfH)\hookrightarrow \sfString(\sfG)$, 这是译文         \begin{equation}
            P_0\sfH\ \hookrightarrow\ P_0\sfG~,
            \quad
            P_0L_0\sfH\ \hookrightarrow\ P_0L_0\sfG~,
            \quad
            L_0L_0\sfH\ \hookrightarrow\ L_0L_0\sfG~.
        \end{equation}  这是译文~\eqref{eq:subgroupN}, 这是译文 $\sfN_\sfH$ 这是译文 $\sfN_\sfG$. 这是译文         \begin{equation}
			\begin{aligned}
				\sfString(\sfG)/\sfString(\sfH)\ &= \ \left(\frac{P_0\sfG\ltimes((P_0L_0\sfG\times\sfU(1))/\sfN_\sfG)}{P_0\sfH\ltimes((P_0L_0\sfH\times \sfU(1))/\sfN_\sfH)}\multirightarrow{2}P_0\sfG/P_0\sfH\right)
                \\
                \ &\cong \ \left(\frac{P_0\sfG\ltimes L_0\sfG}{P_0\sfH\ltimes L_0\sfH}\multirightarrow{2}P_0\sfG/P_0\sfH\right)
                \\
                \ &\cong \ \caL\sfG/\caL\sfH
                \\
                \ &\cong \ \sfG/\sfH~.
			\end{aligned}
		\end{equation}  这是译文 $\caL\sfG\cong\sfG_{\rm cm}$, 这是译文~\eqref{eq:crossedModuleLg} 这是译文~\eqref{eq:definitionLieGroupLieCrossed}, 这是译文 $\caL\sfG\rightarrow \sfG_{\rm cm}$ 这是译文 $\caL\sfH\rightarrow\sfH_{\rm cm}$ 这是译文.

        \begin{corollary}
            这是译文 $\sfString(5)/\sfString(4)$ 这是译文 $S^4$, 这是译文 $(S^4\multirightarrow{2}S^4)$.
        \end{corollary}
        
        \paragraph{这是译文.}  这是译文 $S^4$, 这是译文 
        \begin{equation}\label{eq:nonAbelian_coset_bundle}
            \sfString(5)\ \rightarrow\ S^4\ \cong\ \sfString(5)/\sfString(4)~.
        \end{equation}  这是译文         \begin{equation}
            (S^4\multirightarrow{2}S^4)\ \cong\ \big(L_0S^4\multirightarrow{2}P_0S^4\big)\,,
        \end{equation}  这是译文         \begin{equation}
            \begin{tikzcd}
                \widehat{L_0\sfSpin(5)} \arrow[d]&
                \\
                L_0\sfSpin(5)\arrow[r,shift left]\arrow[r,shift right] \arrow[d] & P_0\sfSpin(5)\arrow[d]
                \\
                L_0S^4\arrow[r,shift left]\arrow[r,shift right] & P_0S^4\arrow[d,"\boundary"]
                \\
                & S^4
            \end{tikzcd}
        \end{equation}  这是译文 $\sfString(4)$-这是译文.~\cite{Aschieri:2003mw,Jurco:2005qj} 这是译文~\cite{Nikolaus:2011ag}, 这是译文.
        
        \paragraph{这是译文.}  这是译文 $\widehat{L_0\sfSpin(5)}$ 这是译文 $L_0S^4$. 这是译文 $H^3(S^4,\IZ)$. 这是译文 $\bar U_1\sqcup\bar U_2\rightarrow S^4$ 这是译文~\eqref{eq:2patchCover}, 这是译文 $\caL\sfSpin(5)/\caL\sfSpin(4)\cong\sfSpin(5)/\sfSpin(4)\cong S^4$ 这是译文 \cref{sec:Spin4Bundle}. 这是译文~\eqref{eq:cocycles_LSpin_bundle}.
        
这是译文.

        \begin{theorem}
            这是译文 $\bar U_1\sqcup\bar U_2$, 这是译文 $S^4$ 这是译文 
            \begin{subequations}
                \begin{equation}
                    \begin{aligned}
                        \hat h_{ijk}\,:\,\bar Y^{[3]}\ &\rightarrow\ \big(P_0L_0\sfSpin(4)\times\sfU(1)\big)/\sfN~,
                        \\
                        q_+\ &\mapsto\ \big[\big((s,t)\mapsto h^\circ_{ijk}(q_+,st),h'_{ijk}\big)\big]\,,
                    \end{aligned}
                \end{equation}
                as well as
                \begin{equation}
                    h'_{111}\ \coloneqq\ h'_{222}\ \coloneqq\ h'_{112}\ \coloneqq\ h'_{122}\ \coloneqq\ h'_{121}\ \coloneqq\ 1~,
                \end{equation}
                and
                \begin{equation}
                    \Lambda'_{11}\ \coloneqq\ \Lambda'_{22}\ \coloneqq\ \Lambda'_{12}\ \coloneqq\ 0 \eand B'_{1}\ \coloneqq\ 0~.
                \end{equation}
                The cocycle relations~\eqref{eq:adjustedCocycleConditions} then reduce to 
                \begin{equation}
                    \begin{aligned}
                        \hat h_{212}\ &=\ g_{21}\acton \hat h_{121}~,
                        \\
                        \hat \Lambda_{21}\ &=\ \hat h_{121}\nabla_1\hat h_{121}^{-1}-g_{21}^{-1}\acton \hat \Lambda_{12}~,
                        \\
                        \hat B_2\ &=\ \hat B_1+\rmd\hat \Lambda_{12}+A_2\acton \hat \Lambda_{12}+\tfrac12[\hat \Lambda_{12},\hat \Lambda_{12}]-\hat \kappa(g^{-1}_{12},F_1-B_1)
                    \end{aligned}
                \end{equation}
                thus providing the missing data. The additional components arising in the lift to a string bundle read as
                \begin{equation}
                    \begin{aligned}
                        h_{212}\ &=\ h_{121}~,
                        \\
                        h'_{212}\ &=\ \exp\left(\frac{\rmi}{2\pi}\int_{0}^{1}\rmd t \int_{0}^{1}\rmd s\innerLarge{g_{21}^{\circ-1}(q,t)\parder[g^\circ_{21}(q,t)]{t}}{h^{\circ-1}_{121}(q,st)\parder[h^{\circ}_{121}(q,st)]{s}}\right) 
                        \\
                        &=\ \exp(\rmi\theta_q)\ =\ \frac{\Re(q)}{|q|}+\rmi\frac{|\Im(q)|}{|q|}~,
                        \\
                        \Lambda'_{21}\ &=\ -\frac{\rmi}{2\pi}\int_0^1\rmd r\left\{\innerLarge{\parder[g^\circ_{12}]{r}g_{12}^{\circ-1}}{A^\circ_1}+\innerLarge{\parder[g^\circ_{21}]{r}g_{21}^{\circ-1}}{A^\circ_2}\right.
                        \\
                        &\kern2cm-\left.\innerLarge{\parder[g^\circ_{21}]{r}g_{21}^{\circ-1}}{g_{12}^{\circ-1} \rmd g^\circ_{12}}-\int_{0}^{1}\rmd s\,\innerLarge{\parder[h^\circ_{121}]{r}h_{121}^{\circ-1}}{\parder[(h^\circ_{121}\rmd h_{121}^{\circ-1})]{s}}\right\},
                        \\
                        B'_{2}\ &=\ \frac{\rmi}{4\pi}\left.\left(\innerLarge{g_{12}^{\circ-1}\rmd g^\circ_{12} + g_{12}^{\circ-1} A^\circ_1 g^\circ_{12}}{A^\circ_2} +\innerLarge{g_{12}^{\circ-1} A^\circ_1 g^\circ_{12}}{g_{12}^{\circ-1}\rmd g^\circ_{12}}\right)\right|_{r=1}
                        \\
                        &\kern1cm+\frac{\rmi}{2\pi}\rmd\int_0^1\rmd r\innerLarge{\parder[g^\circ_{12}]{r}g_{12}^{\circ-1}}{A^\circ_1} 
                        \\
                        &\kern1cm+\frac{\rmi}{4\pi}\int_0^1\rmd r\left\{\innerLarge{\parder[A^\circ_1]{r}}{A^\circ_1}-\innerLarge{\parder[A^\circ_2]{r}}{A^\circ_2}+\innerLarge{\parder[g_{12}^{\circ-1}\rmd g^\circ_{12}]{r}}{g_{12}^{\circ-1}\rmd g^\circ_{12}}\right\}.
                    \end{aligned}
                \end{equation}
            \end{subequations}
        \end{theorem}
        \begin{proof}
            这是译文 $\sfString(4)$-这是译文.
        \end{proof}
        
\noindent 这是译文 \cref{rem:better_cover} 这是译文.
        
        \section{这是译文}
        
 这是译文.~\cite{Borsten:2021ljb} 这是译文 $n$-这是译文 $n=2,\ldots,6$, 这是译文 $n=2$. 这是译文.
        
        \subsection{这是译文}\label{sec:nonAbelianSelfDualString}
        
        \paragraph{这是译文.}
 这是译文 \uline{这是译文} 这是译文 \uline{$\sfU(1)$-这是译文} 这是译文 $A$ 这是译文 $F$ 这是译文 $\IR^3\setminus\{0\}$ 这是译文 $\phi:\IR^3\setminus\{0\}\rightarrow\fru(1)$ 这是译文 $F={\star\rmd\phi}$. 这是译文 (这是译文.~这是译文~1) 这是译文 \uline{这是译文}, 这是译文~\cite{Howe:1997ue}. 这是译文 $B$ 这是译文 $H$ 这是译文 $\IR^4\setminus\{0\}$ 这是译文 $\phi:\IR^4\setminus\{0\}\rightarrow\fru(1)$ 这是译文 $H={\star\rmd\phi}$. 这是译文~\eqref{eq:categorified_Hopf}.
        
        \paragraph{这是译文.}
这是译文 \uline{这是译文}. 这是译文 $\sfG$-这是译文\footnote{这是译文.} $\IR^3$ 这是译文 $\frg$-这是译文 $\phi$ 这是译文 $F={\star\nabla\phi}$. 
        
 这是译文~\cite{Saemann:2017rjm}, 这是译文 $\sfString(\sfG)$-这是译文 $\phi:\IR^4\rightarrow \widehat{L_0\frg}$ 这是译文 
        \begin{equation}\label{eq:non_Abelian_SDS}
            H\ =\ {\star\nabla\phi}~,
            \quad
            \pr_1(\phi)\ =\ 0~,
            \eand
            F\ =\ {\star F}~,
		\end{equation}  这是译文 $F$ 这是译文 $H$ 这是译文, $\nabla=\rmd+A\acton$, 这是译文 $\pr_1:\widehat{L_0\frg}\rightarrow L_0\frg$ 这是译文~\cite{Saemann:2017rjm}, 这是译文~\cite{Saemann:2017rjm}, 这是译文~\cite{Akyol:2012cq}.
        
这是译文~\cite{Saemann:2017rjm}, 这是译文 \cref{sec:stringBundlesWithConnection}, 这是译文.

        \subsection{这是译文}\label{sec:commentsTwistorApproach}
        
 这是译文~\cite{Saemann:2011nb,Saemann:2012uq,Saemann:2013pca,Jurco:2014mva,Jurco:2016qwv,Samann:2017dah}, 这是译文.
        
        \paragraph{这是译文.}  这是译文 $X$. 这是译文 \uline{这是译文} 这是译文 $Z$. 这是译文 \uline{这是译文} $Y$, 这是译文 (这是译文) 这是译文:         \begin{equation}\label{eq:twistorDoubleFibration}
            \begin{tikzcd}
                & Y\arrow[dl,"\pi_2",swap] \arrow[dr,"\pi_1"]
                \\
                Z & & X
            \end{tikzcd}
        \end{equation}  这是译文 $\pi_1$ 这是译文 $\pi_2$ 这是译文 \uline{这是译文},         \begin{equation}
            Z\ \ni\ z\ \leftrightarrow\ \pi_1(\pi_2^{-1}(z))\ \subseteq\ X 
            \eand
            Z\ \supseteq\ \pi_2(\pi_1^{-1}(x))\ \leftrightarrow\ x\ \in\ X~.
        \end{equation}
        
这是译文~\eqref{eq:twistorDoubleFibration},\footnote{这是译文.~\cite{Eastwood:1985aa}} 这是译文 \uline{这是译文} 这是译文 $P$ 这是译文 $Z$. 这是译文~\cite{Ward:1977ta,Hitchin:1982gh} 这是译文~\cite{Witten:2003nn,Popov:2004rb,Popov:2005uv} 这是译文~\cite{Isenberg:1978kk,Witten:1978xx,Khenkin:1980ff,Pool:1981aa,Buchdahl:1985aa,Popov:2021mfl} 这是译文~\cite{Witten:1978xx,Saemann:2005ji,Saemann:2012rr}.\footnote{这是译文~\cite{Penrose:1976js,Atiyah:1978wi} 这是译文~\cite{Merkulov:1992,Wolf:2007tx,Mason:2007ct}, 这是译文~\cite{0264-9381-2-4-020} 这是译文~\cite{Merkulov:1992qa}.} 这是译文 (这是译文) 这是译文 (这是译文) 这是译文 (这是译文) 这是译文 $\caN=3$ 这是译文 $\caN=3$ 这是译文.
        
 这是译文 $P$ 这是译文 $\pi_2$ 这是译文 $\pi_2^*P\rightarrow\tilde P$. 这是译文 $\pi_2$ 这是译文 $\tilde P$ 这是译文 $Y$, 这是译文 $\pi_1$, 这是译文 $\pi_1^!\tilde P$ 这是译文 $X$. 这是译文 $\tilde P$ 这是译文 $\pi_1(\pi_2^{-1}(z))\subseteq X$ 这是译文 $z\in Z$, 这是译文.
        
        \paragraph{这是译文.}
这是译文~\cite{Saemann:2011nb,Mason:2011nw}, 这是译文~\cite{Saemann:2011nb,Saemann:2012uq,Saemann:2013pca,Jurco:2014mva,Jurco:2016qwv,Samann:2017dah}. 这是译文~\eqref{eq:twistorDoubleFibration} 这是译文~\cite{Jurco:2016qwv}. 
        
 这是译文 $\pi_2:Y\rightarrow Z$ 这是译文 $B$ 这是译文 $H$, 这是译文 $\pi_2:Y\rightarrow Z$. 这是译文~\eqref{eq:non_Abelian_SDS}.
        
        \paragraph{这是译文.}  这是译文:         \begin{equation}
            \begin{tikzcd}
                & & \IC^4\times \IC P^1\times \IC P^1 \arrow[dl,"\pi_2",swap] \arrow[ddrr,"\pi_1"]
                \\[10pt]
                & Z_{\rm Ins}\times\IC P^1\arrow[dl,"\pi_{3}",swap] \arrow[dr,"\pi_{4}"] & & 
                \\
                Z_{\rm Hyp} & & Z_{\rm Ins} & & \IC^4
            \end{tikzcd}
        \end{equation}  这是译文, $Z_{\rm Ins}$ 这是译文 $\caO(1)\oplus\caO(1)\rightarrow\IC P^1$ 这是译文 $Z_{\rm Hyp}$ 这是译文~\cite{Saemann:2011nb}, 这是译文 $\caO(1,1)\rightarrow\IC P^1\times\IC P^1$. 这是译文 $\sfSpin(4)\cong\sfSU(2)\times\sfSU(2)$, 这是译文 $\alpha,\beta,\ldots,\dot\alpha,\dot\beta,\ldots=1,2$. 这是译文 $x^{\alpha\dot\beta}$ 这是译文 $\IC^4$, $\lambda_{\dot\alpha}$ 这是译文 $\mu_\alpha$ 这是译文 $\IC P^1$, 这是译文 $y$ 这是译文 $z^\alpha$ 这是译文 $z$ 这是译文 $\caO(1,1)\rightarrow\IC P^1\times\IC P^1$ 这是译文 $\caO(1)\oplus\caO(1)\rightarrow\IC P^1$, 这是译文,         \begin{equation}
            \begin{gathered}
                \pi_1\,:\,(x^{\alpha\dot\beta},\mu_\alpha,\lambda_{\dot\alpha})\ \mapsto\ (x^{\alpha\dot\beta})~,
                \\
                \pi_2\,:\,(x^{\alpha\dot\beta},\mu_\alpha,\lambda_{\dot\alpha})\ \mapsto\ (z^\alpha,\mu_\alpha,\lambda_{\dot\alpha})\ \coloneqq\ (x^{\alpha\dot\beta}\lambda_{\dot\beta},\mu_\alpha,\lambda_{\dot\alpha})~,
                \\
                \pi_3\,:\,(z^\alpha,\mu_\alpha,\lambda_{\dot\alpha})\ \mapsto\ (y,\mu_\alpha,\lambda_{\dot\alpha})\ \coloneqq\ (z^\alpha\mu_\alpha,\mu_\alpha,\lambda_{\dot\alpha})~,
                \\
                \pi_4\,:\,(z^\alpha,\mu_\alpha,\lambda_{\dot\alpha})\ \mapsto\ (z^\alpha,\lambda_{\dot\alpha})~.
            \end{gathered}
        \end{equation}  这是译文.~\cite{Popov:2004rb,Wolf:2010av} 这是译文~\cite{Saemann:2011nb}. 
        
        \paragraph{这是译文.}
 这是译文 $\sfG$-这是译文 $P_0$ 这是译文 $Z_{\rm Ins}$ 这是译文 $\IC P^1\cong(\pi_4\circ\pi_2)(\pi_1^{-1}(x))\subseteq Z_{\rm Ins}$ 这是译文 $x\in\IC^4$ 这是译文~\cite{Ward:1977ta}, 这是译文 $\sfG$ 这是译文 $\pi_4\circ\pi_2$ 这是译文 $\pi_1$. 这是译文 
        \begin{equation}
            P_1\ \coloneqq\ \pi^*_4 P_0~,
        \end{equation}  这是译文 $\sfString(\sfG)$-这是译文 $\scP_1$ 这是译文 $P_1$ 这是译文.
        
 这是译文 $\scG$ 这是译文 $Z_{\rm Hyp}$. 这是译文~\cite{Saemann:2011nb}, 这是译文 $\IC^4$ 这是译文 $\pi_4\circ\pi_3$ 这是译文 $\pi_1$. 这是译文 $\sfString(\sfG)$-这是译文 $\scP_1$ 这是译文 $\pi_3^*\scG$, 这是译文 $\sfString(\sfG)$-这是译文 
        \begin{equation}
            \scP_2\ \coloneqq \scP_1\otimes\pi^*_3\scG~.
        \end{equation}  这是译文 $P_1$ 这是译文 $\scP_1$ 这是译文 $\scG$. 这是译文 $\scG$. 这是译文 $\scP_2$ 这是译文~\cite{Saemann:2012uq}, 这是译文~\eqref{eq:non_Abelian_SDS} 这是译文 $\sfString(\sfG)$ 这是译文 $\IC^4$. 这是译文.~\cite{Saemann:2012uq}.

        \appendix
        \addappheadtotoc
        \appendixpage

        \section{这是译文}
        
这是译文.
        
        \subsection{这是译文}\label{app:useful}
        
这是译文 $(\sfH\overset{\sft}{\longrightarrow}\sfG,\acton)$ 这是译文 $(\frh\overset{\sft}{\longrightarrow}\sfg,\acton)$ 这是译文. 

        \paragraph{这是译文.}  这是译文 $g\acton\unit=\unit$ 这是译文 $\unit\acton h=h$ 这是译文 $g\in\sfG$ 这是译文 $h\in\sfH$. 这是译文:         \begin{subequations}
            \begin{equation}\label{eq:app_identity_t(hA1/h)}
                \begin{gathered}
                    \sft(h^{-1}(g\acton h))\ =\ \sft(h^{-1})g\sft(h)g^{-1}
                    \\
                    \Rightarrow\quad
                    \sft(h^{-1}(V\acton h))\ =\ \sft(h^{-1})V\sft(h)-V~,
                \end{gathered}
            \end{equation}
            \begin{equation}
                \begin{gathered}
                    g_1\acton(g_2\acton h)\ =\ g_2\acton((g_2^{-1}g_1g_2)\acton h)
                    \\
                    \Rightarrow\quad 
                    V\acton(g\acton h)\ =\ g\acton((g^{-1}Vg)\acton h)~,
                \end{gathered}
            \end{equation}
            \begin{equation}
                \begin{gathered}
                    (gVg^{-1})\acton W\ =\ g\acton(V\acton(g^{-1}\acton W))
                    \\
                    \Rightarrow\quad
                    [V_1,V_2]\acton W\ =\ V_1\acton V_2\acton W -  V_2\acton V_1\acton W~,
                \end{gathered}
            \end{equation}
            \begin{equation}
                \begin{gathered}
                    g\acton(\sft(h_1)\acton h_2)\ =\ \sft(g\acton h_1)\acton(g\acton h_2)
                    \\
                    \Rightarrow\quad
                    V\acton[W_1,W_2]\ =\ [V\acton W_1,W_2]+[W_1,V\acton W_2]~,
                \end{gathered}
            \end{equation}
            \begin{equation}\label{eq:app_identity_AhA/h}
                \begin{gathered}
                    g_1\acton(h(g_2\acton h^{-1}))\ =\ \big[(g_1\acton h)h^{-1}\big]h((g_1 g_2 g_1^{-1})\acton h^{-1})\big[(g_1 g_2 g_1^{-1})\acton(h(g_1\acton h^{-1}))\big]
                    \\
                    \Rightarrow\quad V_1\acton(h(V_2\acton h^{-1}))\ =\ V_2\acton(h(V_1\acton h^{-1}))+h([V_1,V_2]\acton h^{-1})
                    \\
                    -[h(V_1\acton h^{-1}),h(V_2\acton h^{-1})]
                \end{gathered}
            \end{equation}
        \end{subequations}  这是译文 $g,g_{1,2}\in\sfG$, 这是译文 $h,h_{1,2}\in\sfH$, 这是译文 $V,V_{1,2}\in\frg$, 这是译文 $W,W_{1,2}\in\frh$. 这是译文 $g=\exp(\eps V)$ 这是译文 $h=\exp(\eps W)$, 这是译文 $\eps$.

这是译文 
        \begin{subequations}
            \begin{align}
                \sft(h_1)\acton(h_2(V\acton h_2^{-1}))\ &=\ h_1h_2(V\acton (h_1h_2)^{-1})-h_1(V\acton h_1^{-1})~,\label{eq:app_identity_ht(B)/h}
                \\
                h(\sft(W)\acton h^{-1})\ &=\ \sft(h)\acton W-W
            \end{align}
        \end{subequations}  这是译文 $h,h_{1,2}\in\sfH$, 这是译文 $V\in\frg$, 这是译文 $W\in\frh$.

        \paragraph{这是译文.}
这是译文 $g\in\scC^\infty(U,\sfG)$ 这是译文 $h\in\scC^\infty(U,\sfH)$ 这是译文 $U\subseteq X$ 这是译文,         \begin{equation}
            (g\acton h)^{-1}\rmd(g\acton h)\ =\ g\acton[ h^{-1}((g^{-1}\rmd g)\acton h)]+g\acton( h^{-1}\rmd h)
        \end{equation}  这是译文 $\rmd(g\acton h)=\rmd g\acton h+g\acton\rmd h$. 这是译文, 
        \begin{subequations}
            \begin{align}
                \rmd(\alpha\acton h)\ &=\ \rmd\alpha\acton h+(-1)^p\alpha\acton\rmd h~,\label{eq:app_identity_d(gLambda)}
                \\
                \rmd(g\acton\beta)\ &=\ g\acton((g^{-1}\rmd g)\acton\beta)+g\acton\rmd\beta~,\label{eq:app_identity_d(hA/h)}
                \\
                \rmd(h^{-1}(\alpha\acton h))\ &=\ h^{-1}(\rmd\alpha\acton h)+(-1)^p[h^{-1}(\alpha\acton h),h^{-1}\rmd h]+(-1)^p\alpha\acton h^{-1}\rmd h 
            \end{align}
        \end{subequations}  这是译文 $\alpha\in\Omega^p(U,\frg)$ 这是译文 $\beta\in\Omega^q(U,\frh)$.
        
        \subsection{这是译文}\label{app:kappaProperties}

这是译文 $(\sfH\overset{\sft}{\longrightarrow}\sfG,\acton)$ 这是译文 $(\frh\overset{\sft}{\longrightarrow}\sfg,\acton)$ 这是译文. 
        
        \paragraph{这是译文.}
 这是译文~\eqref{eq:alternativeAdjustmentCondition_b} 这是译文         \begin{equation}\label{eq:kappa_gginv}
        	g\acton\kappa(g^{-1},V)\ =\ -\kappa(g,g^{-1}Vg-\sft(\kappa(g^{-1},V)))
        \end{equation}  这是译文 $g\in\sfG$ 这是译文 $V\in\frg$. 这是译文 $\kappa$ 这是译文~\eqref{eq:linearisedKappa} 这是译文         \begin{equation}
        	\kappa(V_1,V_2)\ \coloneqq\ \lim_{\eps\to0}\tfrac1\eps\kappa(\exp(\eps V_1),V_2)
        \end{equation}  这是译文 $V_{1,2}\in\frg$. 这是译文 $V\in\frg$ 这是译文 $W\in\frh$ 这是译文~\eqref{eq:alternativeAdjustmentCondition_a} 这是译文 
        \begin{equation}\label{eq:kappat_linear}
        	\kappa(\sft(W),V)\ =\ -V\acton W~.
        \end{equation}  这是译文         \begin{equation}\label{eq:kappa_tt}
        	\kappa(\sft(W_1),\sft(W_2))\ =\ -\kappa(\sft(W_2),\sft(W_1))\ =\ [W_1,W_2]
        \end{equation}  这是译文 $W_{1,2}\in\frh$.
        
这是译文~\cref{def:generalAdjustment} 这是译文 $\kappa$ 这是译文 $\kappa$ 这是译文~\eqref{eq:alternativeAdjustmentCondition_b} 这是译文~\eqref{eq:kappa_gginv}. 这是译文 $\kappa$ 这是译文 $\frg\times\frg\rightarrow\frh$.
        
这是译文~\eqref{eq:alternativeAdjustmentCondition_b} 这是译文 $g_2=g$ 这是译文 $g_1=\exp(\eps V_1)$ 这是译文 $V_1\in\frg$. 这是译文 $g_{1,2}$ 这是译文 $\eps$ 这是译文 $g\exp(\eps V_1)g^{-1}=\exp(\eps gV_1g^{-1})$, 这是译文 
        \begin{equation}\label{eq:kappa_galfaX}
        	\begin{aligned}
        	   &(gV_1g^{-1})\acton\kappa(g,V_2)+\kappa(gV_1g^{-1},gV_2g^{-1}-\sft(\kappa(g,V_2)))
        	   \\ 
        	   &\kern4cm=\ g\acton\kappa(V_1,V_2)+\kappa(g,[V_1,V_2]-\sft(\kappa(V_1,V_2)))
        	\end{aligned}
        \end{equation}  这是译文 $V_2\in\frg$.
        
 这是译文 $\kappa$ 这是译文         \begin{equation}\label{eq:kappa_ggginv}
        	\begin{aligned}
        		\kappa(g_1g_2g^{-1}_1,V)\ &=\ g_1g_2\acton\kappa(g_1^{-1},V)+g_1\acton\kappa(g_2,g_1^{-1}Vg_1-\sft(\kappa(g_1^{-1},V)))
        		\\
        		&\kern1cm+\kappa\big(g_1,g_2g_1^{-1}Vg_1g_2^{-1}-\sft(g_2\acton\kappa(g_1^{-1},V))
        		\\
        		&\kern3cm-\sft(\kappa(g_2,g_1^{-1}Vg_1-\sft(\kappa(g_1^{-1},V))))\big)\,,
        	\end{aligned}
        \end{equation}  这是译文~\eqref{eq:alternativeAdjustmentCondition_b}. 这是译文 $g_{1,2}=\exp(\eps_{1,2}V_{1,2})$ 这是译文 $V_{1,2}\in\frg$ 这是译文 $\eps_1\eps_2$, 这是译文~\eqref{eq:kappa_ggginv} 这是译文 
        \begin{equation}\label{eq:kappa_commutator}
        	\begin{aligned}
        		\kappa([V_1,V_2],V_3)\ &=\ V_1\acton\kappa(V_2,V_2)-V_2\acton\kappa(V_1,V_3) 
        		\\
        		&\kern1cm+\kappa(V_1,[V_2,V_3])-\kappa(V_2,[V_1,V_3])
        		\\
        		&\kern1cm-\kappa(V_1,\sft(\kappa(V_2,V_3)))+\kappa(V_2,\sft(\kappa(V_1,V_3)))
        	\end{aligned}
        \end{equation}  这是译文 $V_3\in\frg$.
        
        \paragraph{这是译文.}
 这是译文 $g\in\scC^\infty(U,\sfG)$ 这是译文 $\alpha\in\Omega^p(U,\frg)$ 这是译文 $U\subseteq X$ 这是译文,         \begin{equation}\label{eq:dofkappa}
        	\rmd\kappa(g,\alpha)\ =\ \kappa(g,\rmd\alpha)+g\acton\kappa(g^{-1}\rmd g,\alpha)+\kappa(g,[g^{-1}\rmd g,\alpha]-\sft(\kappa(g^{-1}\rmd g,\alpha)))~.
        \end{equation}  这是译文 $\rmd\kappa(g,\alpha)=\kappa(\rmd g,\alpha)+\kappa(g,\rmd\alpha)$ 这是译文~\eqref{eq:alternativeAdjustmentCondition_b} 这是译文 $\kappa(g\cdot g^{-1}\rmd g,\alpha)$, 这是译文 $g^{-1}\rmd g$ 这是译文 $\frg$. 这是译文~\eqref{eq:kappa_galfaX} 这是译文~\eqref{eq:dofkappa}, 这是译文         \begin{equation}
        	\begin{aligned}
        		\nabla\kappa(g,\alpha)\ \equiv\ (\rmd+A\acton)\kappa(g,\alpha)\ &=\ \kappa(g,\rmd\alpha)+g\acton\kappa(g^{-1}Ag+g^{-1}\rmd g,\alpha) 
        		\\
        		&\kern0.5cm+\kappa(g,[g^{-1}Ag+g^{-1}\rmd g,\alpha]-\sft(\kappa(g^{-1}Ag+g^{-1}\rmd g,\alpha)))
        		\\
        		&\kern0.5cm-\kappa(A,g\alpha g^{-1}-\sft(\kappa(g,\alpha)))~. 
        	\end{aligned}
        \end{equation}
        
        \paragraph{这是译文 $\sfString(\sfG)$.}
 这是译文 $\hat\kappa$ 这是译文 $\sfString(\sfG)$ 这是译文~\eqref{eq:adjustmentStringGroup} (这是译文 $\kappa$ 这是译文 $\caL\sfG$ 这是译文~\eqref{eq:adjustmentPathLoopGroups}) 这是译文 $g,g_{1,2}\in\scC^\infty(X,P_0\sfG)$, $\alpha,\alpha_{1,2}\in\Omega^p(U,P_0\frg)$, 这是译文 $\beta\in\Omega^q(U,L_0\frg\oplus\fru(1))$. 这是译文 
        \begin{subequations}\label{eq:StringkappaIdentities}
        	\begin{equation}
        		\hat\kappa(g,\sft(\beta))\ =\ g\acton \beta - \beta~,
        	\end{equation}
        	from which it follows that
        	\begin{equation}
        		\hat\kappa(\alpha,\sft(\beta))\ =\ \alpha\acton\beta \ =\ -(-1)^{pq}\hat\kappa(\sft(\beta),\alpha)~.
        	\end{equation}
        \end{subequations}         \begin{subequations}
        	In addition, $\hat\kappa$ also obeys
        	\begin{align}
        			\rmd\hat\kappa(g,\alpha)\ &=\ \hat\kappa(g,\rmd\alpha)+\hat\kappa(\rmd g\,g^{-1},g\alpha g^{-1})~,
        			\\
        			\label{eq:StringkappaIdentity_b}
        			\hat\kappa(g_2g_1,\alpha)\ &=\ \hat\kappa(g_2,g_1\alpha g_1^{-1})+\hat\kappa(g_1,\alpha)~,
        			\\
        			\hat\kappa(g,[\alpha_1,\alpha_2])\ &=\ \hat\kappa(g\alpha_1 g^{-1},g\alpha_2 g^{-1})-\hat\kappa(\alpha_1,\alpha_2)~.
        	\end{align}
        \end{subequations}  这是译文 $\hat\kappa$ 这是译文 $\acton$ 这是译文 $\sfString(\sfG)$.  

        \subsection{这是译文}\label{app:cocycleConsistency}
        
 这是译文 \cref{sec:higherPrincipalBundles}. 这是译文 $(i,j,k)\in Y^{[3]}$, 这是译文 $A_k$ 这是译文 $A_i$ 这是译文         \begin{equation}
            \begin{aligned}
                A_k\ &=\ g^{-1}_{ik}A_ig_{ik}+g^{-1}_{ik}\rmd g_{ik}-\sft(\Lambda_{ik})
                \\
                &\overset{!}{=}\ g^{-1}_{jk}\left(g^{-1}_{ij}A_ig_{ij}+g^{-1}_{ij}\rmd g_{ij}-\sft(\Lambda_{ij})\right)g_{jk}+g^{-1}_{jk}\rmd g_{jk}-\sft(\Lambda_{jk})~.
            \end{aligned}
        \end{equation}  这是译文 
        \begin{equation}\label{eq:cocycleForG}
            g_{ij}g_{jk}\ =\ \sft(h^{-1}_{ijk}) g_{ik}
        \end{equation}  这是译文 $\acton$ 这是译文 $\sft$, 这是译文         \begin{equation}
            g^{-1}_{ik}\sft(h_{ijk})A_i\sft(h^{-1}_{ijk})g_{ik}+\sft\big(g^{-1}_{ik}\acton (h_{ijk} \rmd h^{-1}_{ijk})\big)+g^{-1}_{ik}\rmd g_{ik}-\sft(\Lambda_{jk}+g^{-1}_{jk}\acton\Lambda_{ij})~.
        \end{equation}  这是译文~\eqref{eq:app_identity_t(hA1/h)} 这是译文 $\sft$ 这是译文 $\Lambda$         \begin{equation}\label{eq:cocycleForLambda}
            \Lambda_{ik}\ =\ \Lambda_{jk}+g_{jk}^{-1}\acton\Lambda_{ij}-g_{ik}^{-1}\acton(h_{ijk}\nabla_i h_{ijk}^{-1})~,
        \end{equation}  这是译文 $\nabla_i=\rmd+A_i\acton$. 这是译文 $\Lambda$ 这是译文.
        
        \paragraph{这是译文.}
 这是译文 $ B $, 这是译文 $ A $: 
        \begin{equation}\label{eq:BConsistencyCheck}
            \begin{aligned}
                B_k\ &=\ g^{-1}_{ik}\acton B_i+\nabla_k\Lambda_{ik}+\tfrac12[\Lambda_{ik},\Lambda_{ik}]
                \\
                &\overset{!}{=}\ g^{-1}_{jk}\acton (g^{-1}_{ij}\acton B_i+\nabla_j\Lambda_{ij}+\tfrac12[\Lambda_{ij},\Lambda_{ij}])+\nabla_k\Lambda_{jk}+\tfrac12[\Lambda_{jk},\Lambda_{jk}]~,
            \end{aligned}
        \end{equation}  这是译文 $B_k$ 这是译文$ B_j$ 这是译文 $B_j$ 这是译文 $B_i$. 这是译文~\eqref{eq:cocycleForLambda} 这是译文 $\Lambda$, 这是译文~\eqref{eq:BConsistencyCheck}. 这是译文~\eqref{eq:cocycleForG} 这是译文~\eqref{eq:app_identity_ht(B)/h}, 这是译文         \begin{equation}
            \begin{aligned}
                g^{-1}_{ik}\acton(h_{ijk}\sft(B_i)\acton h^{-1}_{ijk})\ &\overset{!}{=}\ \nabla_k\left(g^{-1}_{jk}\acton \Lambda_{ij}-g^{-1}_{ik}\acton(h_{ijk}\nabla_i h^{-1}_{ijk})\right)-g^{-1}_{jk}\acton(\nabla_j\Lambda_{ij})
                \\
                &\kern1cm+[\Lambda_{jk},g^{-1}_{jk}\acton \Lambda_{ij}-g^{-1}_{ik}\acton(h_{ijk}\nabla_i h^{-1}_{ijk})]
                \\
                &\kern1cm-[g^{-1}_{jk}\Lambda_{ij},g^{-1}_{ik}\acton(h_{ijk}\nabla_i h^{-1}_{ijk})]
                \\
                &\kern1cm+\tfrac12 g^{-1}_{ik}\acton[h_{ijk}\nabla_i h^{-1}_{ijk},h_{ijk}\nabla_i h^{-1}_{ijk}]~.
            \end{aligned}
        \end{equation}  这是译文         \begin{equation}
            \nabla_A(g\acton \Lambda)\ =\ g\acton(\nabla_{g^{-1}Ag+g^{-1}\rmd g}~\Lambda)~,
        \end{equation}  这是译文~\eqref{eq:app_identity_d(gLambda)}, 这是译文 $\acton$ 这是译文 $\sft$ 这是译文         \begin{equation}
            \begin{aligned}
                \nabla_k(g^{-1}_{jk}\acton\Lambda_{ij})\ &=\ g^{-1}_{jk}\acton(\nabla_j\Lambda_{ij})-[\Lambda_{jk},g^{-1}_{jk}\acton\Lambda_{ij}]~,
                \\
                \nabla_k\left(g^{-1}_{ik}\acton(h_{ijk}\nabla_i h^{-1}_{ijk})\right)\ &=\ g^{-1}_{ik}\acton\left(\nabla_i(h_{ijk}\nabla_i h^{-1}_{ijk})\right)-[\Lambda_{ik},g^{-1}_{ik}\acton(h_{ijk}\nabla_i h^{-1}_{ijk})]~, 
            \end{aligned}
        \end{equation}  这是译文 $A$. 这是译文~\eqref{eq:app_identity_AhA/h},~\eqref{eq:app_identity_d(hA/h)}, 这是译文 $\rmd(h\rmd h^{-1})=-\tfrac12[h\rmd h^{-1},h\rmd h^{-1}]$, 这是译文         \begin{equation}
            \nabla_A(h\nabla_A h^{-1})\ =\ h\left(\rmd A+\tfrac12[A,A]\right)\acton h^{-1}-\tfrac12\left[h\nabla_A h^{-1},h\nabla_A h^{-1}\right].
        \end{equation}  这是译文         \begin{equation}
            \begin{aligned}
                &g^{-1}_{ik}\acton\left(h_{ijk}\left(\rmd A_i+\tfrac12[A_i,A_i]+\sft(B_i)\right)\acton h^{-1}_{ijk}\right)
                \\
                &\kern1cm=\ \left[\Lambda_{ik}-\Lambda_{jk}-g_{jk}^{-1}\acton \Lambda_{ij}+g_{ik}^{-1}\acton(h_{ijk}\nabla_i h_{ijk}^{-1}),g_{ik}^{-1}\acton(h_{ijk}\nabla_i h_{ijk}^{-1})\right].
            \end{aligned}
        \end{equation}  这是译文~\eqref{eq:cocycleForLambda} 这是译文         \begin{equation}\label{eq:appFakeFlatUnadjusted}
            g^{-1}_{ik}\acton\left(h_{ijk}(F_i\acton h^{-1}_{ijk})\right)\ =\ 0~.
        \end{equation}  这是译文~\eqref{eq:cocycleForG} 这是译文         \begin{equation}
            (g_{ij}g_{jk})^{-1}\acton\left(h^{-1}_{ijk}(F_i\acton h_{ijk})\right)\ =\ 0~.
        \end{equation}  这是译文~\eqref{eq:BConsistencyCheck} 这是译文 $\kappa(g^{-1},F)$, 这是译文         \begin{equation}\label{eq:appFakeFlatAdjusted}
            (g_{ij}g_{jk})^{-1}\acton\left(h^{-1}_{ijk}(F_i\acton h_{ijk})\right)\ =\ \kappa(g^{-1}_{ik},F_i)-g_{jk}^{-1}\acton\kappa(g^{-1}_{ij},F_i)-\kappa(g^{-1}_{jk},F_j)~.
        \end{equation}
        
 这是译文~\eqref{eq:appFakeFlatUnadjusted} 这是译文~\eqref{eq:appFakeFlatAdjusted} 这是译文 (这是译文) 这是译文.
        
        \section{这是译文}
        
这是译文 (这是译文), 这是译文.
        
        \subsection{这是译文}\label{app:YMCoset}
        
 这是译文 \cref{sec:Spin4Bundle}, 这是译文~\cite[Section II.11]{Kobayashi:1963}. 这是译文~\cite{Trautman:1977im,Harnad:1980ct} 这是译文 (这是译文) 这是译文~\cite{Trautman:1977im,Bais:1985ns}.\footnote{这是译文~\cite{Ivanova:2009yi}.} 
        
        \paragraph{这是译文.}
 这是译文 $\sfG$ 这是译文, $\sfH$ 这是译文 $\sfG$, 这是译文 $\frg$ 这是译文 $\frh$ 这是译文 $\sfG/\sfH$ 这是译文 \uline{这是译文}, 这是译文         \begin{equation}\label{eq:reductiveCosetDecomposition}
            \frg\ \cong\ \frh\oplus\frm
        \end{equation}  这是译文 $\frm$ 这是译文 $\Ad(\sfH)$-这是译文 $[\frh,\frm]\subseteq\frm$, 这是译文 $\sfH$ 这是译文         \begin{equation}\label{eq:reductiveCosetConditions}
            [\frh,\frh]\ \subseteq\ \frh~,
            \quad
            [\frh,\frm]\ \subseteq\ \frm~,
            \eand
            [\frm,\frm]\ \subseteq\ \frh\oplus\frm~.
        \end{equation}  这是译文 $\frg$ 这是译文 $\inner{-}{-}$ 这是译文~\eqref{eq:reductiveCosetDecomposition} 这是译文 $\inner{-}{-}$.
        
 这是译文 $\theta$ 这是译文 (这是译文) 这是译文 $\sfG$. 这是译文~\eqref{eq:reductiveCosetDecomposition}, 这是译文         \begin{equation}
            \theta\ =\ A+m
            \ewith
            A\ \in\ \Omega^1(\sfG,\frh)
            \eand
            m\ \in\ \Omega^1(\sfG,\frm)~.
        \end{equation}  这是译文 $\sfH$-这是译文 $g\mapsto gh$ 这是译文 $\sfG$ 这是译文 $g\in\sfG$ 这是译文 $h\in\sfH$, 这是译文         \begin{equation}
            A\ \mapsto\ h^{-1}Ah+h^{-1}\rmd h
            \eand
            m\ \mapsto\ h^{-1}mh~.
        \end{equation}  这是译文 $\rmd\theta+\frac12[\theta,\theta]=0$ 这是译文         \begin{equation}\label{eq:decompositionMCFormReductive}
            \rmd A+\tfrac12[A,A]\ =\ -\tfrac12[m,m]_\frh
            \eand
            \nabla_Am\ =\ -\tfrac12[m,m]_\frm~,
        \end{equation}  这是译文 $\rmd$ 这是译文 $\sfG$ 这是译文.
        
        \paragraph{这是译文.}
 这是译文 $\sfG$-这是译文 $\pi:P\rightarrow X$ 这是译文 $X$ 这是译文 $\pi^*P\coloneqq P\times_XP=P^{[2]}$ 这是译文         \begin{equation}
            \begin{tikzcd}
                P\times_X P\arrow[r]\arrow[d] & P\arrow[d,"\pi"]
                \\
                P\arrow[r,"\pi"] & X
            \end{tikzcd}
        \end{equation}  这是译文 $P\rightarrow P^{[2]}$ 这是译文 $p\mapsto(p,p)$ 这是译文 $p\in P$. 这是译文, $P^{[2]}\cong P\times\sfG$ 这是译文, $P^{[n]}\cong P\times\sfG^{n-1}$. 这是译文 \v 这是译文 $\pi$ 这是译文 $(p_1,p_2)\mapsto g_{p_1p_2}$ 这是译文 $(p_1,p_2)\in P^{[2]}$, 这是译文 $g_{p_1p_2}\in \sfG$ 这是译文         \begin{equation}
            p_2\ \eqqcolon\ p_1g_{p_1p_2}~,
        \end{equation}  这是译文 $g_{p_1p_2}g_{p_2p_3}=g_{p_1p_3}$ 这是译文 $(p_1,p_2,p_3)\in P^{[3]}$. 这是译文 $\sfH$-这是译文 $\sfG\rightarrow\sfG/\sfH$, 这是译文         \begin{equation}
            \begin{aligned}
                g\,:\,\sfG^{[2]}\ &\rightarrow\ \sfH~,
                \\
                (g_1,g_2)\ &\mapsto\ h_{g_1g_2}\ \coloneqq\ g_1^{-1}g_2~.
            \end{aligned}
        \end{equation}
        
 这是译文 $\{U_i\}_{i\in I}$ 这是译文 $X$ 这是译文 $X=\bigcup_{i\in I}U_i$ 这是译文 $U_i$, $U_i\cap U_j$, $U_i\cap U_j\cap U_k$, 这是译文 $i,j,k,\ldots\in I$ 这是译文 $Y\coloneqq\bigsqcup_{i\in I}U_i$ 这是译文 $X$. 这是译文 $s_i:U_i\rightarrow\sfG$ 这是译文~\eqref{eq:decompositionMCFormReductive}, 这是译文         \begin{equation}\label{eq:canonicalConnection}
            A_{s_i}\ \coloneqq\ s_i^*A\ \in\ \Omega^1(U_i,\frh)
            \eand
            m_{s_i}\ \coloneqq\ s_i^*m\ \in\ \Omega^1(U_i,\frm)~,
        \end{equation}  这是译文 $Y_i$, 这是译文         \begin{equation}\label{eq:decompositionMCFormReductive2}
            \rmd A_{s_i}+\tfrac12[A_{s_i},A_{s_i}]\ =\ -\tfrac12[m_{s_i},m_{s_i}]_\frh
            \eand
            \nabla_{A_{s_i}}m_{s_i}\ =\ -\tfrac12[m_{s_i},m_{s_i}]_\frm~,
        \end{equation}  这是译文 $\rmd$ 这是译文 $Y_i$. 这是译文 $\star$ 这是译文 $\sfG/\sfH$ 这是译文 ${s_i}$ 这是译文 $\sfG$ 这是译文 $\inner{-}{-}$. 这是译文  ~\cite{Trautman:1977im,Harnad:1980ct}         \begin{equation}\label{eq:canonicalCurvature}
            \nabla_{A_{s_i}}{\star}F_{s_i}\ =\ 0
            \ewith
            F_{s_i}\ \coloneqq\ \rmd A_{s_i}+\tfrac12[A_{s_i},A_{s_i}]~.
        \end{equation}  这是译文~\eqref{eq:reductiveCosetConditions} 这是译文~\eqref{eq:decompositionMCFormReductive2}. 这是译文, $A_{s_i}$ 这是译文 $U_i$ 这是译文 $\sfH$. 
        
这是译文 $(g_1,g_2)\in\sfG^{[2]}$ 这是译文 $s_{1,2}:U_{1,2}\rightarrow\sfG$ 这是译文 $(s_{1,2}\circ\pi)(g_{1,2})=g_{1,2}$ 这是译文 $s_2=s_1h_{g_1g_2}$. 这是译文 $A_{s_{1,2}}\coloneqq s_{1,2}^*A\in\Omega^1(U_{1,2},\frh)$, 这是译文         \begin{equation}
            A_{s_2}\ =\ h^{-1}_{g_1g_2}A_{s_1}h_{g_1g_2}+h^{-1}_{g_1g_2}\rmd h_{g_1g_2}~.
        \end{equation}
        
        \subsection{这是译文}\label{app:group_structures}
        
        \paragraph{$\sfG$-这是译文 $\sfH$-这是译文.}  这是译文 $\sfH$-这是译文 $P$ 这是译文 $X$. 这是译文 \uline{这是译文 $\sfH$ 这是译文 $\sfG$} 这是译文 $\phi:\sfG\rightarrow\sfH$, 这是译文 $\sfG$-这是译文 $Q$ 这是译文 $X$, 这是译文 $\phi_Q:Q\times_\sfG\sfH\xrightarrow{\cong}P$.
        
 这是译文 $E\rightarrow X$ 这是译文 $n$, 这是译文 $\sfGL(n)$-这是译文 $X$. 这是译文 \uline{$\sfG$-这是译文} 这是译文 $E$ 这是译文 $\sfGL(n)$ 这是译文 $\sfG$ 这是译文 $\sfG$-这是译文 $\sfG$-这是译文.
        
 这是译文 $\sfG$-这是译文 $E=TX$. 这是译文 \uline{这是译文} 这是译文 $X$ 这是译文 $\sfG$-这是译文 $\sfGL(n)$ 这是译文 $\sfGL^+(n)$ 这是译文 $n=\dim(X)$. 这是译文 \uline{这是译文} 这是译文 $\sfGL(n)$ 这是译文 $\sfO(n)$, 这是译文 \uline{这是译文} 这是译文 $\sfO(n)$ 这是译文 $\sfSO(n)$. 这是译文 \uline{这是译文} 这是译文 $\sfSO(n)$ 这是译文 $\sfSpin(n)$. 这是译文. 
        
        \paragraph{这是译文.}  这是译文 $LX$ 这是译文 $X$. 这是译文 $LX$ 这是译文 $X$ 这是译文~\cite{Witten:1987cg}; 这是译文~\cite{Killingback:1986rd,mclaughlin1992orientation}. 这是译文~\cite{mclaughlin1992orientation}, 这是译文 $LX$ 这是译文 $X$ 这是译文 \uline{这是译文}, 这是译文 $\sfSpin(n)$ 这是译文 $\sfString(n)$, 这是译文 \uline{这是译文}. 

 这是译文, $\sfString(n)$ 这是译文         \begin{equation}\label{eq:whitehead_tower}
            \cdots\ \longrightarrow\ \underbrace{\sfString(n)}_{\eqqcolon\,W_3}\ \longrightarrow\ \underbrace{\sfSpin(n)}_{\eqqcolon\,W_2}\ \longrightarrow\ \underbrace{\sfSpin(n)}_{\eqqcolon\,W_1}\ \longrightarrow\ \underbrace{\sfSO(n)}_{\eqqcolon\,W_0}\ \longrightarrow\ \underbrace{\sfO(n)}_{\eqqcolon\,W}~,
        \end{equation}  这是译文: $\pi_i(W_j)=0$ 这是译文 $i\leq j$ 这是译文 $\pi_i(W_j)$ 这是译文 $i>j$ 这是译文 $\sfString(n)$ 这是译文 $\sfSpin(n)$~\cite{Stolz:2004aa}. 这是译文~\eqref{eq:whitehead_tower}. 这是译文 $\sfString(n)$, 这是译文.~\cite{Nikolaus:2011zg} 这是译文 $\sfString(n)$-这是译文 $\sfString(n)$; 这是译文.~\cite{Demessie:2016ieh}.

 这是译文 $\sfSpin(n)$-这是译文 $P$ 这是译文 $\sfString(n)$, 这是译文 (这是译文 $n\geq 5$) 这是译文 $\tfrac12p_1(P)$ 这是译文 $p_1(P)$ 这是译文 $P$~\cite{mclaughlin1992orientation}. 这是译文~\cite[Section 5]{Stolz:2004aa} 这是译文 $P$ 这是译文 $H^3(X,\IZ)$. 这是译文,~\cite{Waldorf:2009uf} 这是译文 $P$ 这是译文 $P$ 这是译文~\cite{Carey:2004xt}.

        \subsection{这是译文}\label{app:centralExtensions}
        
        \paragraph{这是译文.}
 这是译文 \uline{这是译文} $\sfE$ 这是译文 $\sfG$ 这是译文 $\sfA$ 这是译文         \begin{equation}\label{eq:centralExtension}
            \unit\ \longrightarrow\ \sfA\ \overset{\iota}{\longrightarrow}\ \sfE\ \overset{\pi}{\longrightarrow}\ \sfG\ \longrightarrow\ \unit
        \end{equation}  这是译文 $\im(\iota)$ 这是译文 $\sfE$. 这是译文 $s:\sfG\rightarrow\sfE$ 这是译文 $\pi:\sfE\rightarrow\sfG$, 这是译文, $\pi\circ s=\id$. 这是译文 $g_{1,2}\in\sfG$,         \begin{equation}
            \pi(s(g_1)s(g_2))\ =\ \pi(s(g_1))\pi(s(g_2))\ =\ \pi(s(g_1g_2))~,
        \end{equation}  这是译文 $s$ 这是译文 $c:\sfG\times\sfG\rightarrow\sfA$ 这是译文         \begin{equation}\label{eq:cocycle}
            s(g_1g_2)\iota(c(g_1,g_2))\ \coloneqq\ s(g_1)s(g_2)~.
        \end{equation}  这是译文 $c$ 这是译文         \begin{equation}\label{eq:cocycleConditionapp}
            c(g_1,g_2)c(g_1g_2,g_3)\ =\ c(g_1,g_2g_3)c(g_2,g_3)
        \end{equation}  这是译文 $g_{1,2,3}\in\sfG$. 这是译文 $(g_1,g_2,g_3)=(g,\unit,g)$ 这是译文 $(g_1,g_2,g_3)=(\unit,\unit,g)$ 这是译文 $g\in\sfG$, 这是译文         \begin{equation}\label{eq:normalisation1}
            c(\unit,\unit)\ =\ c(\unit,g)\ =\ c(g,\unit)~.
        \end{equation}  这是译文 $\phi:\sfG\times\sfA\to\sfE$ 这是译文 $\phi(g,a)\coloneqq s(g)\iota(a)$ 这是译文 $(g,a)\in\sfG\times\sfA$. 这是译文 $\phi(g_1,a_1)\phi(g_1,a_1)=s(g_1g_2)\iota(a_1a_2c(g_1,g_2))$ 这是译文 $\sfG\times\sfA$ 这是译文         \begin{equation}\label{eq:product}
            (g_1,a_1)(g_2,a_2)\ \coloneqq\ (g_1g_2,a_1a_2c(g_1,g_2))
        \end{equation}  这是译文 $(g_{1,2},a_{1,2})\in\sfG\times\sfA$. 这是译文~\eqref{eq:cocycleConditionapp} 这是译文~\eqref{eq:product} 这是译文 $\phi$ 这是译文 $(\unit,(c(\unit,\unit))^{-1})$ 这是译文~\eqref{eq:product}, 这是译文 $(g,a)\in\sfG\times\sfA$ 这是译文 $(g,a)^{-1}=(g^{-1},(ac(\unit,\unit)c(g,g^{-1}))^{-1})$.
        
 这是译文 $\tilde s$ 这是译文 $\sfE$, 这是译文 $\pi\circ\tilde s=\id=\pi\circ s$ 这是译文 $\tilde s(g)=s(g)\iota(d(g))$ 这是译文 $g\in\sfG$ 这是译文 $d:\sfG\rightarrow\sfA$. 这是译文         \begin{equation}\label{eq:coboundaryCondition}
            \tilde c(g_1,g_2)\ =\ c(g_1,g_2)(d(g_1g_2))^{-1}d(g_1)d(g_2) 
        \end{equation}  这是译文 $g_{1,2}\in\sfG$. 这是译文~\eqref{eq:coboundaryCondition},~\eqref{eq:normalisation1} 这是译文 $c$ 这是译文         \begin{equation}\label{eq:normalisation2}
            c(\unit,\unit)\ =\ c(\unit,g)\ =\ c(g,\unit)\ =\ \unit
        \end{equation}  这是译文 $g\in\sfG$.
        
 这是译文 $\sfB\sfG$. 这是译文 $\sfG^p\coloneqq\sfG\times\cdots\times\sfG$ ($p$-这是译文) 这是译文 $\sfG^0\coloneqq\unit$ 这是译文 (这是译文) 这是译文 $m_i:\sfG^{p}\rightarrow\sfG^{p-1}$ 这是译文 $i=1,\ldots,p+1$ 这是译文         \begin{equation}
            m_i(g_1,\ldots,g_{p})\ \coloneqq\ 
            \begin{cases}
                (g_2,\ldots,g_{p}) &\efor i=1
                \\
                (g_1,\ldots,g_{i-2},g_{i-1}g_i,g_{i+1},\ldots,g_{p}) &\efor i=2,\ldots,p
                \\
                (g_1,\ldots,g_{p-1}) &\efor i=p+1
            \end{cases}
        \end{equation}  这是译文 $g_{1,\ldots,p}\in\sfG$. 这是译文 $\delta:\sfMap(\sfG^p,\sfA)\rightarrow\sfMap(\sfG^{p+1},\sfA)$ 这是译文         \begin{equation}
            \delta(c)(g_1,\ldots,g_{p+1}) \coloneqq \prod_{i=0}^{\left\lceil p/2\right\rceil}(c\circ m_{2i+1})(g_1,\ldots,g_{p+1})((c\circ m_{2i+2})(g_1,\ldots,g_{p+1}))^{-1}
        \end{equation}  这是译文 $g_{1,\ldots,p+1}\in\sfG$, 这是译文         \begin{equation}\label{eq:extensionComplex}
            \underbrace{\sfMap(\sfG^0,\sfA)}_{\cong\,\sfA}\ \overset{\delta}{\longrightarrow}\ \sfMap(\sfG^1,\sfA)\ \overset{\delta}{\longrightarrow}\ \sfMap(\sfG^2,\sfA)\ \overset{\delta}{\longrightarrow}\ \cdots~.
        \end{equation}  这是译文 $c\in\sfMap(\sfG^2,\sfA)$, 这是译文~\eqref{eq:cocycleConditionapp} 这是译文 $\delta(c)=\unit$ 这是译文 $d\in\sfMap(\sfG^1,\sfA)$, 这是译文~\eqref{eq:coboundaryCondition} 这是译文 $\tilde c(g_1,g_2)=c(g_1,g_2)\delta(d)(g_1,g_2)$ 这是译文 $g_{1,2}\in\sfG$. 这是译文 $H^2(\sfG,\sfA)$ 这是译文~\eqref{eq:extensionComplex}. 这是译文~\eqref{eq:cocycleConditionapp} 这是译文 \uline{这是译文} 这是译文~\eqref{eq:coboundaryCondition} 这是译文 \uline{这是译文}, 这是译文. 
        
 这是译文 $H^1(\sfG,\sfA)$ 这是译文~\eqref{eq:centralExtension} 这是译文 $\sfA$. 这是译文 $H^1(\sfG,\sfA)$ 这是译文~\eqref{eq:centralExtension} 这是译文.~\cite{Azcarraga:2011hqa} 这是译文~\cite[Section 4]{Mickelsson:1989hp} 这是译文.
        
        \begin{remark}\label{rmk:normalisation}  这是译文 $H^2(\sfG,\sfA)$ 这是译文~\eqref{eq:normalisation2} 这是译文~\eqref{eq:product} 这是译文 $(\unit,\unit)$ 这是译文 $(g,a)\in\sfG\times\sfA$ 这是译文 $(g,a)^{-1}=(g^{-1},(ac(g,g^{-1}))^{-1})$. 这是译文~\eqref{eq:cocycleConditionapp} 这是译文 
            \begin{equation}\label{eq:conjugation}
                (g_1,a_1)^{-1}(g_2,a_2)(g_1,a_1)\ =\ (g^{-1}_1g_2g_1,a_2c(g_1^{-1}g_2,g_1)(c(g_1,g_1^{-1}g_2))^{-1})
            \end{equation}  这是译文 $(g_{1,2},a_{1,2})\in\sfG\times\sfA$.
        \end{remark}
        
        \paragraph{这是译文{\'这是译文}这是译文.}
 这是译文 $\sfG$ 这是译文{\'这是译文}这是译文 $\sfA$ 这是译文{\'这是译文}这是译文~\eqref{eq:centralExtension} 这是译文 $\sfE$ 这是译文{\'这是译文}这是译文~\eqref{eq:product} 这是译文 $\pi:\sfE\rightarrow\sfG$ 这是译文 $\sfA$-这是译文, $\sfG$ 这是译文 $c\in H^2(\sfG,\sfA)$ 这是译文 $(\unit,\unit)\in\sfG\times\sfG$. 这是译文~\cite[Prop.~3.11]{Tuynman:1987ij} 这是译文~\cite{neeb2002central} 这是译文.
        
        \paragraph{这是译文 $\sfU(1)$-这是译文.} 这是译文 $\sfA=\sfU(1)$. 这是译文{\'这是译文}这是译文~\eqref{eq:centralExtension} 这是译文 \v 这是译文 $t\in \check H^2(\sfG,\IZ)$ 这是译文.\footnote{这是译文.} 这是译文 $s:\sfG\rightarrow\sfE$ 这是译文 $U$ 这是译文 $\unit\in\sfG$. 这是译文~\eqref{eq:cocycle}, 这是译文 $U\times U$ 这是译文 $(\unit,\unit)\in\sfG\times\sfG$. 这是译文 $\{U_g\}_{g\in\sfG}$ 这是译文 $\sfG$ 这是译文 $U_g\coloneqq L_g(U)$, 这是译文 $L$ 这是译文 $U_g$, 这是译文 $\sigma_g:U_g\rightarrow\pi^{-1}(U_g)$ 这是译文 $\pi:\sfE\rightarrow\sfG$ 这是译文 
        \begin{equation}\label{eq:locSec}
            \sigma_g\ \coloneqq\ L_{s(g)}\circ s\circ L_{g^{-1}}~.
        \end{equation}  这是译文, $\sigma_g(h)=s(g)(s(g^{-1}h))=s(h)\iota(c(g,g^{-1}h))$ 这是译文 $h\in U_g$. 这是译文 $s$ 这是译文 $U$ 这是译文 $\sigma_g$ 这是译文 $U_g$. 这是译文 $t_{g_1g_2}:U_{g_1}\cap U_{g_2}\rightarrow\sfU(1)$ 这是译文 
        \begin{equation}
            \sigma_{g_1}(h)(\sigma_{g_2}(h))^{-1}\ =\ \iota(c(g_1,g_1^{-1}h)(c(g_2,g_2^{-1}h))^{-1})\ \eqqcolon\ \iota(t_{g_1g_2}(h))
        \end{equation}  这是译文 $h\in U_{g_1}\cap U_{g_2}$.\footnote{这是译文 $U_g\ni h\mapsto c(g,g^{-1}h)\in\sfU(1)$ 这是译文 $c$ 这是译文 $(\unit,\unit)\in\sfG\times\sfG$.} 这是译文 $g_{1,2,3}\in\sfG$ 这是译文 $U_{g_1}\cap U_{g_2}\cap U_{g_3}\neq\emptyset$ 这是译文         \begin{equation}
            t_{g_1g_2}(h)t_{g_2g_3}(h)\ =\ t_{g_1g_3}(h)
        \end{equation}  这是译文 $h\in U_{g_1}\cap U_{g_2}\cap U_{g_3}$. 这是译文 $t_{g_1g_2}:U_{g_1}\cap U_{g_2}\rightarrow\sfU(1)$ 这是译文 {\v 这是译文}这是译文 $t_g:U_g\rightarrow\sfU(1)$ 这是译文 $\iota(t_g(h))\coloneqq d(g)d(g^{-1}h)$ 这是译文 $h\in U_g$, 这是译文~\eqref{eq:coboundaryCondition} 这是译文         \begin{equation}
            \tilde t_{g_1g_2}(h)\ =\ t_{g_1}(h)t_{g_1g_2}(h)(t_{g_2}(h))^{-1}~,
        \end{equation}  这是译文, {\v 这是译文}这是译文 {\v 这是译文}这是译文 $\check H^1(\{U_g\}_{g\in\sfG},\sfU(1))$ 这是译文 $\check H^1(\sfG,\sfU(1))\cong\check H^2(\sfG,\IZ)$. 这是译文 $H^2(\sfG,\sfU(1))\rightarrow\check H^2(\sfG,\IZ)$ 这是译文 {\v 这是译文}这是译文 $\sfE\rightarrow\sfG$ 这是译文.  

        \section{这是译文 \texorpdfstring{\cref{sec:stringBundlesWithConnection}}{Section 4}}\label{app:proofs}
        
        \paragraph{这是译文 $(\Ad_g)^* \omega$.}
 这是译文 $\omega$ 这是译文~\eqref{eq:2FormCurvatureKacMoody} 这是译文 $\xi_g$ 这是译文 $g\in\sfG$ 这是译文~\eqref{eq:Baez1Form}. 这是译文 
        \begin{equation}
            (\Ad_g)^* \omega\ =\ \omega+\rmd\xi_g~,
        \end{equation}  这是译文.~\cite[Equation (4.6.6)]{Pressley:1988qk}. 这是译文 $\unit\in L_0\sfG$ 这是译文 $U,V\in L_0\frg$ 这是译文         \begin{equation}
            \begin{aligned}
                \omega_\unit(\Ad_gU,\Ad_gV)\ &=\ \frac{\rmi}{2\pi}\int_0^1\rmd r\innerLarge{\Ad_{g(r)}U(r)}{\parder[\Ad_{g(r)}V(r)]{r}}
                \\ 
                &=\ \frac{\rmi}{2\pi}\int_0^1\rmd r\innerLarge{\Ad_{g(r)}U(r)}{\Ad_{g(r)}\left(\parder[V(r)]{r}+\left[g^{-1}(r)\parder[g(r)]{r},V(r)\right]\right)}.
            \end{aligned}
        \end{equation}  这是译文 $\Ad$-这是译文 $(\rmd\xi_g|_\unit)(U,V)=-\xi_g|_\unit([U,V])$, 这是译文         \begin{equation}
            \begin{aligned}
                \omega_\unit(\Ad_gU,\Ad_gV)\ &=\ \frac{\rmi}{2\pi}\int_0^1\rmd r\innerLarge{U(r)}{\parder[V(r)]{r}}-\frac{\rmi}{2\pi}\int_0^1\rmd r\innerLarge{[U(r),V(r)]}{g^{-1}(r)\parder[g(r)]{r}}
                \\
                &=\ \omega_\unit(U,V)+(\rmd\xi_g|_e)(U,V)~.
            \end{aligned}
        \end{equation}  这是译文.
        
        \paragraph{这是译文.}
 这是译文 $(\sft(\hat h_1)\acton\hat h_2)\hat h_1\hat h_2^{-1}\hat h_1^{-1}\in\sfN$ 这是译文 $\hat h_{1,2}\in P_0L_0\sfG\times\sfU(1)$. 这是译文 $\hat n_{1,2}\in\sfN$, 这是译文 
        \begin{equation}
            \begin{aligned}
                &(\sft(\hat h_1\hat n_1)\acton(\hat h_2\hat n_2))\hat h_1\hat n_1(\hat h_2\hat n_2)^{-1}(\hat h_1\hat n_1)^{-1}
                \\
                &\kern1cm=\ (\sft(\hat h_1)\acton(\hat h_2\hat n_2))\hat h_1\hat n_1\hat h_2^{-1}\hat n_2^{-1}\hat n_1^{-1}\hat h_1^{-1} 
                \\
                &\kern1cm=\ (\sft(\hat h_1)\acton\hat h_2)(\sft(\hat h_1)\acton\hat n_2)\hat h_1\hat n_1\hat h_2^{-1}\hat n_2^{-1}\hat n_1^{-1}\hat h_1^{-1}
                \\
                &\kern1cm=\ (\sft(\hat h_1)\acton\hat h_2)\hat h_1\hat h_2^{-1}\hat h_1^{-1}
                \\
                &\kern2cm\times\underbrace{\hat h_1\hat h_2\hat h_1^{-1}(\sft(\hat h_1)\acton\hat n_2)\hat h_1\hat h_2^{-1}\hat h_1^{-1}}_{\in\,\sfN}
                \\
                &\kern2cm\times\underbrace{\hat h_1\hat h_2\hat n_1\hat h_2^{-1}\hat n_2^{-1}\hat n_1^{-1}\hat h_1^{-1}}_{\in\,\sfN}~,
            \end{aligned}
        \end{equation}  这是译文~\eqref{eq:tMomorphismString2Group}, 这是译文~$\sfN$, 这是译文~$\sfN$. 这是译文, $(\sft(\hat h_1)\acton\hat h_2)\hat h_1\hat h_2^{-1}\hat h_1^{-1}\in\sfN$ 这是译文 $\sft([\hat h_1])\acton[\hat h_2]=[\hat h_1][\hat h_2][\hat h_1]^{-1}$ 这是译文 $[\hat h_1],[\hat h_2]\in\big(P_0L_0\sfG\times\sfU(1)\big)/\sfN$ 这是译文~\eqref{eq:tMomorphismString2Group} 这是译文~\eqref{eq:actionOfP0G}. 
        
 这是译文 $(\sft(\hat h_1)\acton\hat h_2)\hat h_1\hat h_2^{-1}\hat h_1^{-1}\in\sfN$, 这是译文 $\hat h_{1,2}=(f_{1,2},z_{1,2})$. 这是译文~\eqref{eq:conjugation} 这是译文~\eqref{eq:tMomorphismString2Group} 这是译文~\eqref{eq:actionOfP0G}, 这是译文         \begin{subequations}
            \begin{equation}
                f\ \coloneqq\ \boundary(f_1)f_2\boundary(f_1^{-1})f_1f_2^{-1}f_1^{-1}
            \end{equation}
            for the $P_0L_0\sfG$-component of $(\sft(\hat h_1)\acton\hat h_2)\hat h_1\hat h_2^{-1}\hat h_1^{-1}$. Evidently, $\boundary(f)=\unit\in P_0L_0\sfG$ and so $f\in L_0L_0\sfG$. Likewise, again using~\eqref{eq:conjugation} together with~\eqref{eq:tMomorphismString2Group} and~\eqref{eq:actionOfP0G}, the $\sfU(1)$-component of $(\sft(\hat h_1)\acton\hat h_2)\hat h_1\hat h_2^{-1}\hat h_1^{-1}$ is
            \begin{equation}\label{eq:peiffer_proof}
                \begin{aligned}
                    z\ &\coloneqq\ c(\boundary(f_1)f_2\boundary(f_1^{-1}),f_1f_2^{-1}f_1^{-1})\,c^{-1}(f_1f_2,f_1^{-1})
                    \\ 
                    &\kern1cm\times c(f_1^{-1},f_1f_2)\,c^{-1}(f_1f_2f_1^{-1},f_1f_2^{-1}f_1^{-1})\,\exp\left(\int_0^1\rmd t\,\xi_{\boundary(f_1)}\left(f_2^{-1}\parder[f_2]{t}\right)\right)
                \end{aligned}
            \end{equation}
            with $\xi_g$ defined in~\eqref{eq:Baez1Form}.
        \end{subequations}  这是译文         \begin{equation}
            \sfhol^{-1}(f)\ =\ z~.
        \end{equation}  这是译文~\eqref{eq:triangle}. 这是译文 $c(\boundary(f_1)f_2\boundary(f_1^{-1}),f_1f_2^{-1}f_1^{-1})$, 这是译文 $\sfhol^{-1}(f)$. 这是译文~\eqref{eq:peiffer_proof} 这是译文~\eqref{eq:holonomy_shift}. 这是译文 $\omega$ 这是译文~\eqref{eq:2FormCurvatureKacMoody} 这是译文 $\Box$ 这是译文 
        \begin{equation}
            \left(\boundary(f_1)f_1^{-1},\boundary(f_1)\overline{f_2},\boundary(f_1f_2)\overline{f_1^{-1}},\boundary(f_1)f_2\boundary(f_1^{-1})\right)
        \end{equation}  这是译文         \begin{equation}
            (s,t)\ \mapsto\ (\boundary(f_1))(r)f_2(s,r)f_1^{-1}(1-t,r)
            \eforall
            s,t\ \in\ [0,1]~,
        \end{equation}  这是译文 $r$ 这是译文         \begin{equation}
            \int_\Box\omega\ =\ \frac{\rmi}{2\pi}\int_0^1\rmd r\int_0^1\rmd s\innerLarge{f_2^{-1}(s,r)\parder[f_2(s,r)]{s}}{\boundary(f_1^{-1}(s,r))\parder[\,\boundary(f_1(s,r))]{r}}.
        \end{equation}  这是译文 $\omega$ 这是译文 $\boundary(f_1)$ 这是译文~\eqref{eq:helpful_identity} 这是译文 $t$ 这是译文~\eqref{eq:peiffer_proof}.
        
        \paragraph{这是译文.}
 这是译文 $\gamma\in L_0L_0\frg$. 这是译文 $\exp(t\gamma)\in L_0L_0\sfG$ 这是译文 $t$ 这是译文         \begin{equation}\label{eq:holonomy}
            \sfhol(\exp(t\gamma))\ =\ \exp\left(-\int_{D_{\exp(t\gamma)}}\omega\right)\ =\ \exp\left(-\int_{\tilde D_{t\gamma}}\tilde\omega\right),
        \end{equation}  这是译文 $\tilde D_{t\gamma}$ 这是译文 $L_0\frg$ 这是译文 $t\gamma$ 这是译文 $\exp(\tilde D_{t\gamma})=D_{\exp t\gamma}$, 这是译文 
        \begin{equation}
            \tilde\omega\ \coloneqq\ \exp^*\omega\ =\ \frac{\rmi}{4\pi}\int_0^1\rmd r\innerLarge{\theta_{V(r)}}{\parder[\theta_{V(r)}]{r}}\ \in\ \Omega^2(L_0\frg)~,
        \end{equation}  这是译文         \begin{equation}
            \theta_V\ \coloneqq\ \frac{1-\rme^{-\ad_V}}{\ad_V}\rmd V
        \end{equation}  这是译文 $V$ 这是译文 $V\mapsto V$ 这是译文 $L_0\frg$. 这是译文~\eqref{eq:holonomy} 这是译文 $t$ 这是译文         \begin{equation}
            \begin{aligned}
                \sfhol(\exp(t\gamma))\ &=\ 1-\frac{\rmi}{4\pi}\int_{\tilde D_{t\gamma}}\int_0^1\rmd r\innerLarge{V(r)}{\parder[V(r)]{r}}+\caO(t^3)
                \\
                &=\ 1+\frac{\rmi}{4\pi}\int_{t\gamma}\int_0^1\rmd r\innerLarge{\rmd V(r)}{\parder[V(r)]{r}}+\caO(t^3)
                \\
                &=\ 1+\frac{\rmi t^2}{4\pi}\int_0^1\rmd s\int_0^1\rmd r\innerLarge{\parder[\gamma(s,r)]{s}}{\parder[\gamma(s,r)]{r}}+\caO(t^3)~.
            \end{aligned} 
        \end{equation}  这是译文 $L_0\frg$ 这是译文 $s\mapsto t\gamma(s) $. 这是译文         \begin{equation}\label{eq:holder}
            \dder{t}\bigg|_0\sfhol(\exp(t\gamma))\ =\ 0~,
        \end{equation}  这是译文~\eqref{eq:derhol}.\footnote{这是译文~\cite{Mackaay:2000ac} 这是译文~\cite{Barrett:1991aj}.}
        
 这是译文 $ \ell $ 这是译文 $ L_0 L_0 \sfG $ 这是译文 $ t $ 这是译文 $ \ell_0$. 这是译文 
        \begin{equation}\label{eq:holderhol}
            \begin{aligned}
                &\sfhol(\ell_0)\dder{t}\bigg|_0\sfhol^{-1}(\ell(t)) 
                \\
                &\kern.5cm=\ -\frac{\rmi}{2\pi}\int_0^1\rmd s\int_0^1\rmd r\innerLarge{\parder{s}\left(\dder{t}\bigg|_0\ell(t,s,r)\ell_0^{-1}(s,r)\right)}{\parder[\ell_0(s,r)]{r}\ell_0^{-1}(s,r)}
                \\
                &\kern.5cm=\ -\frac{\rmi}{2\pi}\int_0^1\rmd s\int_0^1\rmd r\innerLarge{\ell_0^{-1}(s,r)\parder[\ell_0(s,r)]{s}}{\parder{r}\left(\ell_0^{-1}(s,r)\dder{t}\bigg|_0\ell(t,s,r)\right)}.
            \end{aligned}
        \end{equation}  这是译文         \begin{equation}
            \dder{t}\bigg|_0\big(\sfhol(\ell_0)\,\sfhol^{-1}(\ell(t))\,c^{-1}(\ell_0,\ell_0^{-1})\,c(\ell_0^{-1},\ell(t))\big)\ =\ \dder{t}\bigg|_0\sfhol^{-1}(\ell_0^{-1}\ell(t))\ =\ 0~,
        \end{equation}  这是译文 $ (\ell_0, \sfhol^{-1}(\ell_0))^{-1}\cdot (\ell(t),\sfhol^{-1}(\ell(t))) \in \sfN $, 这是译文~\eqref{eq:holder}.
        
        \paragraph{这是译文 $\widehat{L_0\sfG}$.}
 这是译文 $ P_0L_0\sfG\times\sfU(1)$ 这是译文 $\omega$ 这是译文~\eqref{eq:2FormCurvatureKacMoody}. 这是译文 $ \omega \in \Omega^2(L_0\sfG,\fru(1))$ 这是译文 $P_0L_0\sfG$ 这是译文 $\Phi_t:P_0L_0\sfG\rightarrow P_0L_0\sfG$ 这是译文 $t\in[0,\infty)$ 这是译文 $(\Phi_t f)(s,r)\coloneqq f(\rme^{-t}s,r)$ 这是译文 $f\in P_0L_0\sfG$. 这是译文         \begin{equation}
            \hat\mu\ \coloneqq\ \mu'-\int_0^\infty\rmd t\,\Phi^*_t(\iota_\xi\pi_1^*\boundary^*\omega)\ \in\ \Omega^1(P_0L_0\sfG\times\sfU(1),\fru(1))~,
        \end{equation}  这是译文 $\mu'$ 这是译文 $P_0L_0\sfG\times\sfU(1)$ 这是译文 (这是译文) 这是译文 $\sfU(1)$, $\xi$ 这是译文 $\Phi_t$, $\iota_\xi$ 这是译文 $\xi$, $ \pi_1 $ 这是译文 $ P_0L_0\sfG \times \sfU(1)$, 这是译文 $\boundary$ 这是译文 $P_0L_0\sfG\rightarrow L_0\sfG$. 这是译文 $(f,z)\in P_0L_0\sfG\times\sfU(1)$ 这是译文         \begin{equation}\label{eq:hatmu}
            \hat{\mu}_{(f,z)}\ =\ z^{-1}\rmd z+\frac{\rmi}{2\pi}\int_0^1\rmd s\int_0^1\rmd r\,\innerLarge{f^{-1}(s,r)\parder[f(s,r)]{s}}{\parder{r}(f^{-1}(s,r)\rmd f(s,r))}.
        \end{equation}  这是译文         \begin{subequations}
            \begin{equation}
                \rmd\left(f^{-1}(s,r)\parder[f(s,r)]{s}\right)\ =\ \parder{s}(f^{-1}(s,r)\rmd f(s,r))+\left[f^{-1}(s,r)\parder[f(s,r)]{s},f^{-1}(s,r)\rmd f(s,r)\right],    
            \end{equation}
            \begin{equation}
                \begin{aligned}
                    &\rmd\innerLarge{f^{-1}(s,r)\parder[f(s,r)]{s}}{\parder{r}(f^{-1}(s,r)\rmd f(s,r))}
                    \\
                    &\kern1cm=\ \innerLarge{\parder{s}(f^{-1}(s,r)\rmd f(s,r))}{\parder{r}(f^{-1}(s,r)\rmd f(s,r))},
                \end{aligned}
            \end{equation}
            and
            \begin{equation}
                \begin{aligned}
                    &\int_0^1\rmd s\int_0^1\rmd r\innerLarge{\parder{s}(f^{-1}(s,r)\rmd f(s,r))}{\parder{r}(f^{-1}(s,r)\rmd f(s,r))}
                    \\
                    &\kern2cm=\ \frac{1}{2}\int_0^1\rmd r\innerLarge{f^{-1}(1,r)\rmd f(1,r)}{\parder{r}(f^{-1}(1,r)\rmd f(1,r))},
                \end{aligned}
            \end{equation}
        \end{subequations}  这是译文 
        \begin{equation}\label{eq:dofmu}
            \rmd\hat\mu\ =\ \pi_1^*\boundary^*\omega~.
        \end{equation}  这是译文 $\hat\mu$ 这是译文 $(P_0L_0\sfG\times\sfU(1))/\sfN\cong\widehat{L_0\sfG}$. 这是译文 $\sfN$. 这是译文 $\hat h\in P_0L_0\sfG\times\sfU(1)$ 这是译文 $\hat n\in\sfN$, 这是译文         \begin{equation}\label{eq:muinvar}
            R^*_{\hat n}\hat\mu_{\hat h\hat n}\ =\ \hat\mu_{\hat h}\ =\ L^*_{\hat n}\hat\mu_{\hat n \hat h}~,
        \end{equation}  这是译文 $R_{\hat n}$ 这是译文 $ L_{\hat n} $ 这是译文 $\hat n$, 这是译文~\eqref{eq:holderhol}. 这是译文\footnote{这是译文 $\Phi_t^*\alpha=\rme^{t\caL_V}\alpha$ 这是译文 $\Phi_t$ 这是译文 $V$ 这是译文 $\alpha$.} 这是译文 
        \begin{equation}
            \caL_{V^{R,L}}\hat \mu\ =\ 0 ~,
        \end{equation}  这是译文 $V^R$ 这是译文 $V^L$ 这是译文 $\sfN$, 这是译文 $\caL_V=\rmd\iota_V+\iota_V\rmd$ 这是译文 $\iota_V\rmd\hat\mu =\iota_V\pi_1^*\boundary^*\omega=0$, 这是译文 $V^{R,L}$ 这是译文 $\boundary$, 这是译文         \begin{equation}\label{eq:hatmuX}
            \iota_V\hat\mu\ =\ 0~.
        \end{equation}  这是译文 $L(t)=(\ell(t),\sfhol^{-1}(\ell(t)))\in\sfN$ 这是译文 $L(0)=(\unit_{P_0L_0\sfG},\unit_{\sfU(1)})$. 这是译文 $\dot\ell(0)\coloneqq\dder{t}\big|_0\ell(t)\in L_0L_0\frg$. 这是译文 $F=(f,z)\in P_0L_0\sfG$, 这是译文 $\sfN$ 这是译文         \begin{equation}
            V^R_F\ =\ \dder{t}\bigg|_0F\cdot L(t)\ =\ \left(f\dot\ell(0),z\dder{t}\bigg|_0c(f,\ell(t))\right),
        \end{equation}  这是译文~\eqref{eq:holder} 这是译文~\eqref{eq:hatmu}, 这是译文         \begin{equation}\label{eq:hatmuX_explicit}
            \hat\mu_F (X_F^R)\ =\ \dder{t}\bigg|_0 c(f,\ell(t))+\frac{\rmi}{2\pi}\int_0^1\rmd s\int_0^1\rmd r\,\innerLarge{f^{-1}(s,r)\parder[f(s,r)]{s}}{\parder[\dot \ell(0)]{r}}\ =\ 0~,
        \end{equation}  这是译文~\eqref{eq:cocyclePLGUBaez} 这是译文~\eqref{eq:helpful_identity}. 这是译文~\eqref{eq:hatmuX} 这是译文 $ V^L $ 这是译文 $ \sfN $. 
        
 这是译文~\eqref{eq:hatmuX_explicit} (这是译文~\eqref{eq:hatmuX}), 这是译文 $\hat\mu$ 这是译文 $\widehat{L_0\sfG}$, 这是译文~\eqref{eq:dofmu} 这是译文 $\omega$. 这是译文~\eqref{eq:hatmu} 这是译文.
        
        \paragraph{这是译文 $\widehat{L_0\sfG}$.}
 这是译文 (这是译文) 这是译文 $\widehat{L_0\sfG}$. 这是译文 $P_0L_0\sfG\times\sfU(1)$. 这是译文 $\hat h(t)=(f(t),z(t))$ 这是译文 $P_0L_0\sfG\times\sfU(1)$ 这是译文 $t\in[0,1]$ 这是译文 $\dder{t}\big|_0\hat h(t)\in T_{\hat h(0)}(P_0L_0\sfG\times\sfU(1))$ 这是译文 $\hat h(0)$. 这是译文 $\theta^L$ 这是译文 $\dder{t}\big|_0\hat h(t)$,         \begin{equation}
            \hat\theta^L_{\hat h(0)}\left(\dder{t}\bigg|_0\hat h(t)\right)\ \coloneqq\ \hat h^{-1}(0)\dder{t}\bigg|_0\hat h(t)~,
        \end{equation}  这是译文         \begin{equation}\label{eq:mcl}
            \hat\theta^L_{(f,z)}\ =\ \big(\theta^L_f,\hat\mu_{(f,z)}\big)
        \end{equation}  这是译文 $(f,z)\in P_0L_0\sfG\times\sfU(1)$ 这是译文 $\hat\mu$ 这是译文~\eqref{eq:hatmu}. 这是译文 
        \begin{equation}\label{eq:mcr}
            \hat\theta^R_{(f,z)}\ =\ \left(\theta^R_f,\hat\mu_{(f,z)}+\frac{\rmi}{2\pi}\int_0^1\rmd r\innerLarge{f^{-1}(1,r)\parder[\,f(1,r)]{r}}{\theta^L_{f(1,r)}}\right),
        \end{equation}  这是译文 $\theta^R$ 这是译文 $P_0L_0\sfG$. 这是译文~\eqref{eq:dofmu}, 这是译文         \begin{equation}
            \rmd\hat\theta^{L,R}\ =\ \mp\tfrac12[\hat\theta^{L,R},\hat\theta^{L,R}]~,
        \end{equation}  这是译文 $ L/R $. 这是译文, $\hat\mu$ 这是译文 $\widehat{L_0\sfG}$. 这是译文 $\hat\theta^L$ 这是译文 $\hat\theta^R $ 这是译文 $\widehat{L_0\sfG}$. 
        
        \paragraph{这是译文 $\hat h\nabla\hat h^{-1}$.}  这是译文~\eqref{eq:adjustedCocycleConditions},~\eqref{eq:adjustedCoboundaryConditions}, 这是译文~\eqref{eq:adjustedHigherGaugeTransformations} 这是译文 $ \sfString(\sfG) $-这是译文         \begin{equation}
            \hat h\nabla\hat h^{-1}\ =\ \hat h\rmd\hat h^{-1}+ \hat h(A\acton\hat h^{-1})~,
        \end{equation}  这是译文 $\hat h=[(f,z)]\in (P_0L_0\sfG\times\sfU(1))/\sfN$ 这是译文 $A$ 这是译文 $P_0\frg$. 这是译文 $\hat h\rmd\hat h^{-1}=-\hat \theta^R_{\hat h}$ 这是译文 $P_0\sfG$ 这是译文         \begin{equation}
            \begin{aligned}
                &\hat h(g\acton\hat h^{-1})
                \\
                &\kern5pt\ =\  [(f,z)](g\acton[(f,z)]^{-1})
                \\
                &\kern5pt\ =\ \bigg[\bigg(fgf^{-1}g^{-1},c^{-1}(f,f^{-1})c(f,gf^{-1}g^{-1})
                \\
                &\kern3cm\times\exp\left(-\frac{\rmi}{2\pi}\int_0^1\rmd s\int_0^1\rmd r\innerLarge{g^{-1}(r)\parder[g(r)]{r}}{\parder[f(s,r)]{s}f^{-1}(s,r)}\right)\bigg)\bigg],
            \end{aligned}
        \end{equation}  这是译文 $g\in P_0\sfG$, 这是译文~\eqref{eq:actionOfP0G} 这是译文 $P_0L_0\sfG\times \sfU(1)$. 这是译文 $ g=1+A $ 这是译文 $\parder{s}A=0$, 这是译文 $\frn$, 这是译文         \begin{equation}\label{eq:hAhinv}
            \hat h (A\acton\hat h^{-1})\ =\ \left(\boundary(f)[A,\boundary(f^{-1})],\frac{\rmi}{2\pi}\int_0^1\rmd r\innerLarge{\boundary(f^{-1}(s,r))\parder[\,\boundary(f(s,r))]{r}}{A(r)}\right).
        \end{equation}

    \end{body}

\end{document}
